% Les lettres de motivation et d'encouragement — Allemand 1ère A — Cours signé OnBuch+
% Programme officiel MINESEC. Compiler avec : tectonic ALA19-les-lettres-de-motivation-et-d-encouragement.tex
\newcommand{\DOCMATIERE}{Allemand}
\newcommand{\DOCNIVEAU}{1ère A}
\newcommand{\DOCMODULE}{Chapitre 4 : Média et communication}
\newcommand{\DOCLECON}{ALA19}
\newcommand{\DOCTITRE}{Les lettres de motivation et d'encouragement}
% OnBuch+ — préambule « cours signés » ENRICHI (compilé avec tectonic / XeLaTeX)
% Système visuel « Pop » d'OnBuch+ : fond crème (jamais blanc), bordures encre
% épaisses, ombres « dures » décalées, titres Archivo Black en capitales.
% Version MATHÉMATIQUES : graphes (pgfplots/tikz), boîtes pédagogiques
% (définition, propriété, méthode, attention, le savais-tu, exercice résolu,
% exercice, TP, bilan), page de garde et signature OnBuch+ ancrée en bas.
%
% Macros attendues dans main.tex AVANT \input{preamble} :
%   \DOCMATIERE  \DOCNIVEAU  \DOCMODULE  \DOCLECON (numéro)  \DOCTITRE
\documentclass[11pt]{article}

\usepackage{fontspec}
\usepackage[ngerman,french]{babel} % français = langue principale ; l'allemand via \all{..} / textebox
\usepackage[a4paper,margin=2cm,top=3.4cm,bottom=2.8cm,headheight=1.4cm,footskip=1.3cm]{geometry}
\usepackage{amsmath,amssymb,amsfonts}
\usepackage{mathtools} % \xmapsto, \coloneqq... souvent utilisés par les modèles
\usepackage{cancel} % simplifications barrées (bilans de forces)
% \abs{z} (module, valeur absolue) et \norm{u} (norme), utilisés par les modèles
\providecommand{\abs}{}\let\abs\relax\DeclarePairedDelimiter{\abs}{\lvert}{\rvert}
\providecommand{\norm}{}\let\norm\relax\DeclarePairedDelimiter{\norm}{\lVert}{\rVert}
\usepackage[table]{xcolor} % [table] : fournit aussi \rowcolors (alternance de lignes)
\usepackage{pagecolor}
\usepackage{tikz}
\usetikzlibrary{arrows.meta,positioning,calc,shapes.geometric,shapes.misc,shapes.arrows,decorations.pathmorphing,decorations.pathreplacing,decorations.markings,patterns,shadows,mindmap,trees,fit,backgrounds,matrix,intersections,angles,quotes}
\usepackage{pgfplots}
\usepackage[version=4]{mhchem} % équations nucléaires \ce{^{235}_{92}U}
\usepackage[european,siunitx]{circuitikz} % schémas électriques
\pgfplotsset{compat=1.18}
\usepgfplotslibrary{fillbetween,groupplots} % aires entre courbes (calcul intégral)
\usepackage{enumitem}
\usepackage{fancyhdr}
% [most] charge toutes les bibliothèques tcolorbox (listings, minted, raster,
% documentation...) et gonfle inutilement la mémoire TeX sur un document long
% (~300 boîtes) : on ne charge que ce qui est réellement utilisé.
\usepackage[skins,breakable]{tcolorbox}
\usepackage{graphicx}
\usepackage{colortbl}
\usepackage{booktabs}
\usepackage{tabularx}
\usepackage{diagbox} % case d'en-tête diagonale des tableaux à double entrée
% Colonnes extensibles centrée / gauche / droite (C, L, R), souvent utilisées par les modèles
\newcolumntype{C}{>{\centering\arraybackslash}X}
\newcolumntype{L}{>{\raggedright\arraybackslash}X}
\newcolumntype{R}{>{\raggedleft\arraybackslash}X}
\usepackage{array}
\usepackage{multicol}
\usepackage{multirow}
\usepackage{siunitx}
% parse-numbers=false : accepte aussi \SI{2,5 \times 10^{-3}}{...}
\sisetup{locale=FR,output-decimal-marker={,},group-minimum-digits=4,parse-numbers=false}
\DeclareSIUnit\ohm{\ensuremath{\symup{\Omega}}} % U+2126 absent de la police
\DeclareSIUnit\division{div} % réglages d'oscilloscope (ms/div, V/div)
\DeclareSIUnit\tr{tr} % tours (vitesse de rotation, ex. \SI{1500}{\tr\per\minute})
\DeclareSIUnit\molar{M} % concentration molaire (ex. \SI{3}{\molar}, \SI{1}{\micro\molar})
\DeclareSIUnit\an{an} % année (le modèle confond parfois avec \year, un compteur TeX) — ex. \SI{5}{\kilo\meter\per\an}
\DeclareSIUnit\FCFA{FCFA} % franc CFA (ex. \SI{500}{\FCFA\per\kilo\gram})
\DeclareSIUnit\degreeBrix{\textdegree Brix} % teneur en sucre d'un sirop/jus (réfractométrie)
\DeclareSIUnit\month{mois} % le modèle utilise parfois \month, un compteur TeX, comme unité de durée
\usepackage{qrcode}
% \degree : commande fréquemment utilisée par les modèles pour un angle ou une
% température en dehors de \SI{}{} (non fournie par les paquets chargés).
\providecommand{\degree}{^{\circ}}
% \qed (fin de démonstration), utilisé par les modèles sans amsthm
\providecommand{\qed}{\hfill$\square$}
% \barre : double barre des valeurs interdites dans un tableau de variations (tkz-tab)
\providecommand{\barre}{\ensuremath{\|}}
% environnement proof (amsthm non chargé) utilisé par les modèles
\ifdefined\proof\else\newenvironment{proof}[1][Démonstration]{\par\noindent\textit{#1.}\ }{\qed\par}\fi
\usepackage{lastpage}
\usepackage{hyperref}
\hypersetup{hidelinks}

% ============================================================
% Palette « Pop » OnBuch+ — mêmes hex que la classe P (plus_ui.dart)
% ============================================================
\definecolor{poporange}{HTML}{F59321}
\definecolor{popdark}{HTML}{E07A0C}
\definecolor{popcream}{HTML}{FFF6E8}
\definecolor{popink}{HTML}{1C1714}
\definecolor{poppurple}{HTML}{7A5AE0}
\definecolor{popgreen}{HTML}{2E9E6A}
\definecolor{poppink}{HTML}{E0457B}
\definecolor{popblue}{HTML}{2F7DE1}
\definecolor{popgold}{HTML}{C98A12}
\definecolor{popmuted}{HTML}{8A7F76}
\definecolor{popline}{HTML}{EADFD2}
\definecolor{popgray}{HTML}{8A7F76}
\colorlet{popgrey}{popgray}
% Alias tolérants (le modèle nomme parfois les couleurs différemment)
\colorlet{popwhite}{white}
\colorlet{popblack}{popink}
\colorlet{popred}{poppink}
\colorlet{popyellow}{popgold}
% Teintes claires (fonds de tableaux/encadrés) — le modèle nomme parfois
% les couleurs avec un suffixe L (ex. popblueL) sans les avoir définies.
\colorlet{poporangeL}{poporange!20}
\colorlet{popdarkL}{popdark!20}
\colorlet{poppurpleL}{poppurple!20}
\colorlet{popgreenL}{popgreen!20}
\colorlet{poppinkL}{poppink!20}
\colorlet{popblueL}{popblue!20}
\colorlet{popgoldL}{popgold!20}
\colorlet{popredL}{popred!20}
\colorlet{popgrayL}{popgray!20}
% Teintes claires pour fonds de tableaux / zones
\colorlet{popblueL}{popblue!10!white}
\colorlet{poporangeL}{poporange!14!white}
\colorlet{popgreenL}{popgreen!12!white}
\colorlet{poppurpleL}{poppurple!12!white}
\colorlet{poppinkL}{poppink!12!white}
\colorlet{popgoldL}{popgold!14!white}

\pagecolor{popcream}

% ============================================================
% Typographie
% ============================================================
\defaultfontfeatures{Ligatures=TeX}
\setmainfont{PlusJakartaSans-Medium.ttf}[Path=../../../fonts/,BoldFont=PlusJakartaSans-Bold.ttf]
% Symboles, API et emojis absents de la police principale : repli sur DejaVu Sans (caractères actifs)
\newfontface\symfont{DejaVuSans.ttf}[Path=../../../fonts/]
\makeatletter
\newcommand\symactive[1]{\catcode#1=13 \begingroup\lccode`\~=#1 \lowercase{\endgroup\def~}{{\symfont\char#1}}}
\newcommand\symblank[1]{\catcode#1=13 \begingroup\lccode`\~=#1 \lowercase{\endgroup\def~}{}}
\newcommand\symhyph[1]{\catcode#1=13 \begingroup\lccode`\~=#1 \lowercase{\endgroup\def~}{\mbox{-}}}
\newcount\symc
\newcommand\symrange[2]{\symc=#1 \@whilenum\symc<#2 \do{\expandafter\symactive\expandafter{\the\symc}\advance\symc by 1 }}
\newcommand\symblankrange[2]{\symc=#1 \@whilenum\symc<#2 \do{\expandafter\symblank\expandafter{\the\symc}\advance\symc by 1 }}
\makeatother
\symrange{"0250}{"02FF}\symactive{"2713}\symactive{"2714}\symactive{"2717}\symactive{"2718}\symactive{"2610}\symactive{"2611}\symactive{"2612}
\symhyph{"2011}\symhyph{"2E1A}\symblankrange{"1F300}{"1FAFF}\symblank{"FE0F}
% Maths en sans-serif (Fira Math) avec chiffres et lettres droites de Plus
% Jakarta, pour que les formules se fondent dans le texte.
\usepackage[math-style=ISO,bold-style=ISO]{unicode-math}
\setmathfont{FiraMath-Regular.otf}[Path=../../../fonts/]
\setmathfont{PlusJakartaSans-Medium.ttf}[Path=../../../fonts/,range={up/num,up/{Latin,latin},\mathup}]
\usepackage{icomma} % virgule décimale sans espace en mode math
% Caractères mathématiques Unicode absents des polices de texte (Plus Jakarta,
% Archivo Black) : rendus par leur équivalent mathématique, en texte comme en
% formule — sinon ils s'affichent en carré vide (ex. « Arithmétique dans ℤ »).
\usepackage{newunicodechar}
\newunicodechar{ℝ}{\ensuremath{\mathbb{R}}}
\newunicodechar{ℤ}{\ensuremath{\mathbb{Z}}}
\newunicodechar{ℂ}{\ensuremath{\mathbb{C}}}
\newunicodechar{ℕ}{\ensuremath{\mathbb{N}}}
\newunicodechar{ℚ}{\ensuremath{\mathbb{Q}}}
\newunicodechar{𝔻}{\ensuremath{\mathbb{D}}}
\newunicodechar{α}{\ensuremath{\alpha}}
\newunicodechar{β}{\ensuremath{\beta}}
\newunicodechar{γ}{\ensuremath{\gamma}}
\newunicodechar{δ}{\ensuremath{\delta}}
\newunicodechar{ε}{\ensuremath{\varepsilon}}
\newunicodechar{θ}{\ensuremath{\theta}}
\newunicodechar{λ}{\ensuremath{\lambda}}
\newunicodechar{μ}{\ensuremath{\mu}}
\newunicodechar{σ}{\ensuremath{\sigma}}
\newunicodechar{τ}{\ensuremath{\tau}}
\newunicodechar{φ}{\ensuremath{\varphi}}
\newunicodechar{ψ}{\ensuremath{\psi}}
\newunicodechar{ω}{\ensuremath{\omega}}
\newunicodechar{Σ}{\ensuremath{\Sigma}}
% majuscules produites par \MakeUppercase dans les titres (α-aminés -> Α)
\newunicodechar{Α}{\ensuremath{\alpha}}
\newunicodechar{Β}{\ensuremath{\beta}}
\newunicodechar{Γ}{\ensuremath{\Gamma}}
\newunicodechar{Δ}{\ensuremath{\Delta}}
\newunicodechar{Ε}{\ensuremath{\varepsilon}}
\newunicodechar{Θ}{\ensuremath{\Theta}}
\newunicodechar{Λ}{\ensuremath{\Lambda}}
\newunicodechar{Μ}{\ensuremath{\mu}}
\newunicodechar{Τ}{\ensuremath{\tau}}
\newunicodechar{Φ}{\ensuremath{\Phi}}
\newunicodechar{Ψ}{\ensuremath{\Psi}}
\newunicodechar{Ω}{\ensuremath{\Omega}}
\newunicodechar{↦}{\ensuremath{\mapsto}}
\newunicodechar{∈}{\ensuremath{\in}}
\newunicodechar{∉}{\ensuremath{\notin}}
\newunicodechar{∧}{\ensuremath{\wedge}}
\newunicodechar{⇔}{\ensuremath{\Leftrightarrow}}
\newunicodechar{⟺}{\ensuremath{\Longleftrightarrow}}
\newunicodechar{⇒}{\ensuremath{\Rightarrow}}
\newunicodechar{⟹}{\ensuremath{\Longrightarrow}}
\newunicodechar{∘}{\ensuremath{\circ}}
\newunicodechar{∪}{\ensuremath{\cup}}
\newunicodechar{∩}{\ensuremath{\cap}}
\newunicodechar{≡}{\ensuremath{\equiv}}
\newunicodechar{⊂}{\ensuremath{\subset}}
\newunicodechar{⊥}{\ensuremath{\perp}}
\newunicodechar{∃}{\ensuremath{\exists}}
\newunicodechar{∀}{\ensuremath{\forall}}
\newunicodechar{∛}{\ensuremath{\sqrt[3]{\,}}}
\newunicodechar{∜}{\ensuremath{\sqrt[4]{\,}}}
\newunicodechar{⃗}{\ensuremath{{}^{\to}}}
\newunicodechar{ⁿ}{\ifmmode^{n}\else\textsuperscript{\ensuremath{n}}\fi}
\newunicodechar{ˣ}{\ifmmode^{x}\else\textsuperscript{\ensuremath{x}}\fi}
\newunicodechar{ᵇ}{\ifmmode^{b}\else\textsuperscript{\ensuremath{b}}\fi}
\newunicodechar{ʳ}{\ifmmode^{r}\else\textsuperscript{\ensuremath{r}}\fi}
\newunicodechar{ᵘ}{\ifmmode^{u}\else\textsuperscript{\ensuremath{u}}\fi}
\newunicodechar{ʸ}{\ifmmode^{y}\else\textsuperscript{\ensuremath{y}}\fi}
\newunicodechar{ᵃ}{\ifmmode^{a}\else\textsuperscript{\ensuremath{a}}\fi}
\newunicodechar{ᵉ}{\ifmmode^{e}\else\textsuperscript{\ensuremath{e}}\fi}
\newunicodechar{ᵏ}{\ifmmode^{k}\else\textsuperscript{\ensuremath{k}}\fi}
\newunicodechar{ᵖ}{\ifmmode^{p}\else\textsuperscript{\ensuremath{p}}\fi}
\newunicodechar{ᴺ}{\ifmmode^{N}\else\textsuperscript{\ensuremath{N}}\fi}
\newunicodechar{ᶜ}{\ifmmode^{c}\else\textsuperscript{\ensuremath{c}}\fi}
\newunicodechar{ʰ}{\ifmmode^{h}\else\textsuperscript{\ensuremath{h}}\fi}
\newunicodechar{ᵠ}{\ifmmode^{q}\else\textsuperscript{\ensuremath{q}}\fi}
\newunicodechar{ₙ}{\ifmmode_{n}\else\textsubscript{\ensuremath{n}}\fi}
\newunicodechar{ₐ}{\ifmmode_{a}\else\textsubscript{\ensuremath{a}}\fi}
\newunicodechar{ᵢ}{\ifmmode_{i}\else\textsubscript{\ensuremath{i}}\fi}
\newunicodechar{ᵣ}{\ifmmode_{r}\else\textsubscript{\ensuremath{r}}\fi}
\newunicodechar{ₖ}{\ifmmode_{k}\else\textsubscript{\ensuremath{k}}\fi}
\newunicodechar{ᵦ}{\ifmmode_{\beta}\else\textsubscript{\ensuremath{\beta}}\fi}
\newunicodechar{ₓ}{\ifmmode_{x}\else\textsubscript{\ensuremath{x}}\fi}
\newfontfamily\popdisplay{ArchivoBlack-Regular.ttf}[Path=../../../fonts/]
\newfontfamily\poplogo{SpaceGrotesk-Bold.ttf}[Path=../../../fonts/]
\newfontfamily\popsemi{PlusJakartaSans-SemiBold.ttf}[Path=../../../fonts/]
\newfontfamily\popxbold{PlusJakartaSans-ExtraBold.ttf}[Path=../../../fonts/]
\newfontfamily\popxboldls{PlusJakartaSans-ExtraBold.ttf}[Path=../../../fonts/,LetterSpace=3.0]

\setlist[enumerate]{leftmargin=1.7em,itemsep=5pt,topsep=4pt}
\setlist[itemize]{leftmargin=1.7em,itemsep=4pt,topsep=4pt,label={\color{poporange}\textbullet}}
\setlength{\parindent}{0pt}
\setlength{\parskip}{5pt}
\renewcommand{\arraystretch}{1.25}

% pgfplots : style Pop par défaut
\pgfplotsset{
  every axis/.append style={axis line style={popink,line width=1pt},tick style={popink},
    grid style={popline,line width=0.5pt},label style={font=\small},tick label style={font=\footnotesize}},
  popcurve/.style={poporange,line width=1.8pt,no markers,smooth},
}
% Même style disponible dans un \draw TikZ simple (hors pgfplots)
\tikzset{popcurve/.style={poporange,line width=1.8pt}}
\tikzset{popgrid/.style={popline,very thin}}
\colorlet{popcurve}{poporange} % le nom du style est parfois utilisé comme couleur

% ============================================================
% Pill / squiggle
% ============================================================
\newcommand{\poppill}[3]{%
  \tikz[baseline=-0.55ex]\node[draw=popink,line width=1.2pt,rounded corners=2.6pt,%
    fill=#2,inner xsep=6.5pt,inner ysep=3pt]{%
    \popxboldls\fontsize{7}{7}\selectfont\color{#3}\MakeUppercase{#1}};%
}
\newcommand{\popsquiggle}[1]{%
  \tikz[baseline=-0.4ex]{
    \draw[#1,line width=2.6pt,line cap=round]
      plot [smooth] coordinates {(0,0.09) (0.34,0.01) (0.68,0.21) (1.02,0.01) (1.36,0.10)};
  }%
}
% Mot-clé surligné (marqueur orange sous le texte)
\newcommand{\cle}[1]{{\popxbold\color{popdark}#1}}

% ============================================================
% En-tête / pied
% ============================================================
\pagestyle{fancy}
\fancyhf{}
\renewcommand{\headrulewidth}{0pt}
\renewcommand{\footrulewidth}{0pt}
\lhead{\raisebox{-0.15ex}{\poplogo\fontsize{13}{13}\selectfont\color{popink}OnBuch\color{poporange}+}}
\rhead{\poppill{\DOCMATIERE\ \textbullet\ \DOCNIVEAU\ \textbullet\ Leçon \DOCLECON}{white}{popink}}
\lfoot{\popxbold\fontsize{7.6}{7.6}\selectfont\color{popmuted}\MakeUppercase{Cours signé OnBuch+}\ \textbullet\ {\popsemi\DOCTITRE}}
\rfoot{\tikz[baseline=-0.6ex]\node[rounded corners=8.5pt,fill=popink,minimum height=17pt,minimum width=17pt,inner xsep=5pt,inner sep=0]{\popxbold\fontsize{8}{8}\selectfont\color{popcream}\thepage\,/\,\pageref*{LastPage}};}

% ============================================================
% Titres
% ============================================================
\newcommand{\chaptitre}[2]{%
  \begin{tcolorbox}[enhanced,colback=poporange,colframe=popink,boxrule=2.6pt,arc=11pt,
      drop shadow=popink,left=15pt,right=15pt,top=12pt,bottom=11pt,before skip=3pt,after skip=18pt]
    \poppill{#2}{popink}{popcream}\par\vspace{9pt}
    {\raggedright\hyphenpenalty=10000\exhyphenpenalty=10000\popdisplay\fontsize{22}{24}\selectfont\color{popcream}\MakeUppercase{#1}\par}
  \end{tcolorbox}
}
% Section numérotée : pastille orange avec le numéro + titre Archivo
\newcounter{popsec}
\newcommand{\coursec}[1]{%
  \stepcounter{popsec}\setcounter{popsub}{0}%
  \par\vspace{16pt}\noindent
  \begin{minipage}{\linewidth}
  \tikz[baseline=-0.7ex]\node[draw=popink,line width=1.6pt,fill=poporange,rounded corners=4pt,minimum size=22pt,inner sep=0]{\popdisplay\fontsize{12}{12}\selectfont\color{popcream}\thepopsec};%
  \hspace{8pt}{\popdisplay\fontsize{15}{17}\selectfont\color{popink}\MakeUppercase{#1}}\par
  \vspace{4pt}\noindent{\color{poporange}\rule{36pt}{3.6pt}}
  \end{minipage}\par\nopagebreak\vspace{8pt}
}
\newcounter{popsub}
\newcommand{\courssub}[1]{%
  \stepcounter{popsub}%
  \par\vspace{9pt}\noindent{\popxbold\fontsize{11.5}{13}\selectfont\color{popink}{\color{poporange}\thepopsec.\thepopsub}\hspace{6pt}#1}\par\nopagebreak\vspace{4pt}
}

% ============================================================
% Boîtes pédagogiques — même recette : fond blanc, bordure épaisse colorée,
% ombre dure encre, badge plein. Titre optionnel [#1].
% ============================================================
\tcbset{popbase/.style={breakable,enhanced,colback=white,boxrule=2.3pt,arc=9pt,
  shadow={3pt}{-3pt}{0pt}{popink},left=14pt,right=14pt,top=11pt,bottom=11pt,before skip=11pt,after skip=15pt,
  fontupper=\popsemi\fontsize{10.6}{14.5}\selectfont\color{popink},
  fontlower=\popsemi\fontsize{10.6}{14.5}\selectfont\color{popink}}}
% Coche dessinée (le glyphe ✓ n'existe pas dans les polices du document)
\newcommand{\popcheck}[1]{\tikz[baseline=-0.5ex]\draw[#1,line width=1.6pt,line cap=round,line join=round](-0.13,0.0)--(-0.04,-0.09)--(0.13,0.1);}
\providecommand{\coche}{\popcheck{popgreen}}
\providecommand{\resultat}[1]{\par\smallskip\noindent\fcolorbox{popgreen}{popgreenL}{\parbox{\dimexpr\linewidth-2\fboxsep-2\fboxrule\relax}{\textbf{Résultat.} #1}}\par\smallskip}
\newcommand{\popboxhead}[3]{\poppill{#1}{#2}{white}\ifx\relax#3\relax\else\hspace{6pt}{\popxbold\fontsize{10.6}{12}\selectfont\color{popink}#3}\fi\par\vspace{7pt}}

\newtcolorbox{aretenir}[1][]{popbase,colframe=popgreen,before upper={\popboxhead{À retenir}{popgreen}{#1}}}
% Texte en allemand (document de lecture, dialogue, extrait) : cadre
% d'étude. Un titre optionnel (source, auteur) s'affiche dans l'en-tête.
\newtcolorbox{textebox}[1][]{popbase,colframe=popblue,before upper={\popboxhead{Text}{popblue}{#1}\begin{otherlanguage}{ngerman}},after upper={\end{otherlanguage}}}
% Marques ✓ / ✗ (forme correcte / fautive) souvent utilisées par les modèles
\providecommand{\cross}{\textcolor{poppink}{\ensuremath{\times}}}
\providecommand{\xmark}{\cross}\providecommand{\crossmark}{\cross}\providecommand{\cmark}{\ensuremath{\checkmark}}
% Ligne de réponse / justification à compléter (inventée par les modèles)
\providecommand{\justification}[1]{\par\noindent\textit{Justification~:} \dotfill\par}
% \textipa (package tipa non chargé) : la police ne couvre pas l'API -> simple passage
\providecommand{\textipa}[1]{#1}
% Soulignement (ulem non chargé) : \uline, \ul
\providecommand{\uline}[1]{\underline{#1}}\providecommand{\ul}[1]{\underline{#1}}
% \glossaire{\item ...} inventée par les modèles : liste de glossaire
\providecommand{\glossaire}[1]{\begin{itemize}#1\end{itemize}}
% \alert (beamer) inventée par les modèles : texte mis en évidence
\providecommand{\alert}[1]{\textbf{\textcolor{poppink}{#1}}}
% variantes ASCII d'ordinaux écrites en macro
\providecommand{\iere}{\textsuperscript{re}}\providecommand{\ieme}{\textsuperscript{e}}\providecommand{\ere}{\textsuperscript{re}}\providecommand{\eme}{\textsuperscript{e}}\providecommand{\er}{\textsuperscript{er}}\providecommand{\ie}{\textsuperscript{e}}
% Ordinaux français écrits en macro par les modèles : 1\ière, 3\ième
\newcommand{\ière}{\textsuperscript{re}}\newcommand{\ième}{\textsuperscript{e}}
% \exemple (inventée par les modèles) : étiquette « Exemple : »
\providecommand{\exemple}{\textbf{Exemple~:}\ }
% \square utilisable aussi hors mode math (cases à cocher)
\AtBeginDocument{\let\squareMath\square\renewcommand{\square}{\ensuremath{\squareMath}}}
% Mot ou phrase en allemand à l'intérieur d'un paragraphe français
\newcommand{\all}[1]{\ifmmode\text{\textit{#1}}\else\begingroup\selectlanguage{ngerman}\itshape #1\endgroup\fi}
\let\esp\all % alias toléré
\newtcolorbox{exemplebox}[1][]{popbase,colframe=poporange,before upper={\popboxhead{Exemple}{poporange}{#1}}}
\newtcolorbox{definition}[1][]{popbase,colframe=popblue,before upper={\popboxhead{Définition}{popblue}{#1}}}
\newtcolorbox{propriete}[1][]{popbase,colframe=poppurple,before upper={\popboxhead{Propriété}{poppurple}{#1}}}
\newtcolorbox{methode}[1][]{popbase,colframe=popgold,before upper={\popboxhead{Méthode}{popgold}{#1}}}
\newtcolorbox{attention}[1][]{popbase,colframe=poppink,before upper={\popboxhead{Attention}{poppink}{#1}}}
\newtcolorbox{experience}[1][]{popbase,colframe=popblue,colback=popblueL,before upper={\popboxhead{Expérience}{popblue}{#1}}}
% « Le savais-tu ? » : carte violette pleine (contexte, culture, Cameroun)
\newtcolorbox{savaistu}[1][]{popbase,colback=poppurple,colframe=popink,
  fontupper=\popsemi\fontsize{10.4}{14.2}\selectfont\color{white},
  before upper={\poppill{Le savais-tu\,?}{white}{poppurple}\ifx\relax#1\relax\else\hspace{6pt}{\popxbold\color{white}#1}\fi\par\vspace{7pt}}}
% Exercice résolu : énoncé en haut, \tcblower puis la solution sur fond crème
\newcounter{popexo}\newcounter{popexr}
\newtcolorbox{exoresolu}[1][]{popbase,colframe=poporange,colbacklower=popcream,
  segmentation style={popink,line width=1.2pt,dashed},
  before upper={\stepcounter{popexr}\popboxhead{Exercice résolu \thepopexr}{poporange}{#1}},
  before lower={\poppill{Solution}{popink}{popcream}\par\vspace{6pt}}}
% Exercice (non résolu en place) — la correction est dans la partie Corrigés
\newtcolorbox{exercice}[1][]{popbase,colframe=popink,
  before upper={\stepcounter{popexo}\popboxhead{Exercice \thepopexo}{popink}{#1}}}
% Corrigé
\newtcolorbox{corrige}[1][]{popbase,colframe=popgreen,colback=popgreenL,
  before upper={\popboxhead{Corrigé}{popgreen}{#1}}}
% Objectifs / prérequis / bilan
\newtcolorbox{objectifs}{popbase,colframe=popink,colback=poporangeL,
  before upper={\popboxhead{Objectifs de la leçon}{popink}{}}}
\newtcolorbox{prerequis}{popbase,colframe=popink,colback=popblueL,
  before upper={\popboxhead{Prérequis}{popblue}{}}}
\newtcolorbox{bilan}{popbase,colframe=popink,colback=white,boxrule=2.8pt,
  before upper={\popboxhead{Fiche bilan}{poporange}{L'essentiel à connaître pour l'examen}}}
% Figure encadrée avec légende
\newtcolorbox{popfigure}[1][]{enhanced,breakable=false,colback=white,colframe=popline,boxrule=1.4pt,arc=8pt,
  left=10pt,right=10pt,top=10pt,bottom=8pt,before skip=10pt,after skip=12pt,halign=center,
  lower separated=false,halign lower=center,
  fontlower=\popsemi\fontsize{9}{11.5}\selectfont\color{popmuted},
  #1}
\newcommand{\legende}[1]{\par\vspace{6pt}{\popsemi\fontsize{9}{11.5}\selectfont\color{popmuted}#1}}

% ============================================================
% Signature OnBuch+
% ============================================================
\newcommand{\poplogoinline}[1]{{\poplogo\fontsize{#1}{#1}\selectfont\color{popink}OnBuch\color{poporange}+}}
\newcommand{\signatureonbuch}{%
  \begin{tcolorbox}[enhanced,colback=popink,colframe=popink,boxrule=2.2pt,arc=10pt,
      drop shadow=poporange,left=14pt,right=14pt,top=10pt,bottom=10pt,before skip=0pt,after skip=0pt]
    \tikz[baseline=-0.7ex]\node[circle,fill=popgreen,draw=popcream,line width=1.3pt,minimum size=22pt,inner sep=0]{\popcheck{white}};%
    \hspace{9pt}\begin{minipage}[c]{0.62\linewidth}
      {\popxbold\fontsize{12}{13}\selectfont\color{popcream}Signé\ \textbullet\ {\poplogo OnBuch\color{poporange}+}}\par\vspace{2pt}
      {\raggedright\popsemi\fontsize{8}{10.5}\selectfont\color{popline}\MakeUppercase{Équipe pédagogique OnBuch+}\\\MakeUppercase{Conforme au programme officiel MINESEC}\par}
    \end{minipage}\hfill
    {\poplogo\fontsize{22}{22}\selectfont\color{popcream}OnBuch\color{poporange}+}
  \end{tcolorbox}%
}

% ============================================================
% Page de garde
% #1 titre · #2 matière · #3 niveau · #4 module · #5 n° leçon · #6 durée ·
% #7 description / objectifs (texte ou itemize)
% ============================================================
\newcommand{\pagegarde}[7]{%
  \thispagestyle{empty}
  \newgeometry{a4paper,margin=1.7cm,top=1.6cm,bottom=1.4cm}%
  \begin{tikzpicture}[remember picture,overlay]
    \fill[poporange!18] ([xshift=-3.2cm,yshift=-3.6cm]current page.north east) circle (4.6cm);
    \fill[poppurple!14] ([xshift=2.4cm,yshift=7.6cm]current page.south west) circle (3.8cm);
  \end{tikzpicture}%
  {\poplogo\fontsize{30}{30}\selectfont\color{popink}OnBuch\color{poporange}+}\hfill
  \raisebox{0.6ex}{\poppill{Cours signé \textbullet\ Premium}{popink}{popcream}}\par
  \vspace{1pt}\popsquiggle{poppurple}\par
  \vspace{8pt}
  \begin{tcolorbox}[enhanced,colback=poporange,colframe=popink,boxrule=2.8pt,arc=13pt,
      drop shadow=popink,left=18pt,right=18pt,top=12pt,bottom=13pt,after skip=10pt]
    \poppill{#2 \textbullet\ #3}{popink}{popcream}\hspace{5pt}\poppill{#4}{white}{popink}\par\vspace{8pt}
    {\popdisplay\fontsize{14}{14}\selectfont\color{popink}LEÇON #5}\par\vspace{4pt}
    {\raggedright\hyphenpenalty=10000\exhyphenpenalty=10000\popdisplay\fontsize{25}{27}\selectfont\color{popcream}\MakeUppercase{#1}\par}
    \vspace{8pt}
    {\popxbold\fontsize{9.5}{9.5}\selectfont\color{popink}\MakeUppercase{Durée conseillée : #6}}
  \end{tcolorbox}
  \begin{tcolorbox}[enhanced,colback=white,colframe=popink,boxrule=2.2pt,arc=10pt,
      drop shadow=popink,left=16pt,right=16pt,top=10pt,bottom=10pt,after skip=8pt]
    \poppill{Dans cette leçon}{poppurple}{white}\par\vspace{5pt}
    \popsemi\fontsize{9.3}{12.6}\selectfont\color{popink}#7
  \end{tcolorbox}
  \noindent\begin{minipage}[t]{0.487\linewidth}
    \begin{tcolorbox}[enhanced,colback=popgreen,colframe=popink,boxrule=2.2pt,arc=10pt,
        drop shadow=popink,left=11pt,right=11pt,top=10pt,bottom=10pt,height=3.1cm,valign=center]
      \tcbox[colback=white,colframe=popink,boxrule=1.3pt,arc=5pt,boxsep=0pt,nobeforeafter,
        left=3pt,right=3pt,top=3pt,bottom=3pt,valign=center]{\qrcode[height=1.9cm]{https://play.google.com/store/apps/details?id=com.luvvix.onbuch&hl=fr}}%
      \hspace{8pt}\begin{minipage}[c]{0.5\linewidth}
        {\popdisplay\fontsize{11}{12.5}\selectfont\color{white}\MakeUppercase{Télécharge OnBuch}}\par\vspace{4pt}
        {\raggedright\popsemi\fontsize{8.4}{11}\selectfont\color{white}Gratuit \textbullet\ Google Play\\Tuteur IA, annales, concours, résultats\par}
      \end{minipage}
    \end{tcolorbox}
  \end{minipage}\hfill
  \begin{minipage}[t]{0.487\linewidth}
    \begin{tcolorbox}[enhanced,colback=poppurple,colframe=popink,boxrule=2.2pt,arc=10pt,
        drop shadow=popink,left=11pt,right=11pt,top=10pt,bottom=10pt,height=3.1cm,valign=center]
      \tikz[baseline=-0.7ex]\node[circle,fill=white,draw=popink,line width=1.3pt,minimum size=1.6cm,inner sep=0]{\popdisplay\fontsize{24}{24}\selectfont\color{poppurple}+};%
      \hspace{8pt}\begin{minipage}[c]{0.6\linewidth}
        {\popdisplay\fontsize{11}{12.5}\selectfont\color{white}\MakeUppercase{Passe à OnBuch+}}\par\vspace{4pt}
        {\raggedright\popsemi\fontsize{8.4}{11}\selectfont\color{white}Tous les cours signés, Tuteur IA illimité, fiches PDF — onglet OnBuch+ de l'app\par}
      \end{minipage}
    \end{tcolorbox}
  \end{minipage}
  \vspace{8pt}
  \popcontenu
  \vfill
  \signatureonbuch
  \restoregeometry
  \newpage
}

% Rangée « Ce cours contient » (page de garde)
\newcommand{\poptile}[3]{% #1 couleur #2 gros texte #3 légende
  \begin{tcolorbox}[enhanced,colback=#1,colframe=popink,boxrule=2pt,arc=9pt,shadow={2.5pt}{-2.5pt}{0pt}{popink},
      left=6pt,right=6pt,top=9pt,bottom=9pt,halign=center,valign=center,height=1.75cm,width=\linewidth,nobeforeafter]
    {\popdisplay\fontsize{12.5}{13}\selectfont\color{white}\MakeUppercase{#2}}\par\vspace{3pt}
    {\centering\popsemi\fontsize{7.8}{9.5}\selectfont\color{white}#3\par}
  \end{tcolorbox}}
\newcommand{\popcontenu}{%
  \par\noindent{\popxboldls\fontsize{7.5}{7.5}\selectfont\color{popmuted}CE COURS SIGNÉ CONTIENT}\par\vspace{7pt}
  \noindent\begin{minipage}[t]{0.235\linewidth}\poptile{popblue}{Cours}{complet, illustré, pas à pas}\end{minipage}\hfill
  \begin{minipage}[t]{0.235\linewidth}\poptile{poporange}{Méthodes}{et exercices résolus}\end{minipage}\hfill
  \begin{minipage}[t]{0.235\linewidth}\poptile{poppink}{Exercices}{type examen + corrigés}\end{minipage}\hfill
  \begin{minipage}[t]{0.235\linewidth}\poptile{popgreen}{Fiche bilan}{l'essentiel à retenir}\end{minipage}\par}

% Sommaire
\newtcolorbox{sommaire}{enhanced,colback=white,colframe=popink,boxrule=2.2pt,arc=10pt,
  drop shadow=popink,left=18pt,right=18pt,top=13pt,bottom=13pt,before skip=6pt,after skip=20pt,
  before upper={\poppill{Sommaire}{poppurple}{white}\par\vspace{9pt}\popsemi\fontsize{10.6}{14.5}\selectfont\color{popink}}}

% Signature de fin de cours — poussée tout en bas de la dernière page
\newcommand{\findecours}{%
  \par\vfill\nopagebreak
  \begin{center}\popsquiggle{poporange}\end{center}\vspace{4pt}
  \signatureonbuch
}
\AtBeginDocument{\def\llbracket{\mathopen{[\mkern-3.5mu[}}\def\rrbracket{\mathclose{]\mkern-3.5mu]}}} % intervalles d'entiers (glyphe ⟦ absent de la police)
% Glyphes absents de Fira Math : redéfinis à partir de glyphes disponibles
\AtBeginDocument{%
  \def\setminus{\mathbin{\backslash}}\let\smallsetminus\setminus
  \def\vdots{\mathord{\vbox{\baselineskip3.2pt\lineskiplimit0pt\kern4pt\hbox{.}\hbox{.}\hbox{.}}}}%
  \def\ddots{\mathinner{\mkern1mu\raise6.5pt\hbox{.}\mkern2mu\raise3.6pt\hbox{.}\mkern2mu\raise0.7pt\hbox{.}\mkern1mu}}%
  \def\triangle{\mathord{\scalebox{0.9}{$\symup{\Delta}$}}}%
  \def\nabla{\mathord{\raisebox{\depth}{\scalebox{1}[-1]{$\symup{\Delta}$}}}}%
  \def\checkmark{\popcheck{popdark}}%
  \def\star{\mathbin{\ast}}\def\bigstar{\mathord{\ast}}%
  \def\diamond{\mathbin{\scalebox{0.6}{$\Diamond$}}}%
  \def\bigcup{\mathop{\mathchoice{\raisebox{-0.3em}{\scalebox{1.7}{$\cup$}}}{\scalebox{1.25}{$\cup$}}{\cup}{\cup}}\displaylimits}%
  \def\bigcap{\mathop{\mathchoice{\raisebox{-0.3em}{\scalebox{1.7}{$\cap$}}}{\scalebox{1.25}{$\cap$}}{\cap}{\cap}}\displaylimits}%
  \def\lesseqqgtr{\mathrel{\lessgtr}}\def\bigoplus{\mathop{\oplus}\displaylimits}%
}
\providecommand{\ssi}{si et seulement si} % abréviation fréquente
\AtBeginDocument{\providecommand{\wideparen}{\overparen}} % arc de cercle

%% FIGFIX-BEGIN v5 — correctifs globaux des figures (docs/onbuch-generation/tools/figfix)
\usepackage{adjustbox}\usepackage{etoolbox}\tcbuselibrary{raster}
% toute figure TikZ/pgfplots plus large que la ligne est réduite à la largeur disponible
\BeforeBeginEnvironment{tikzpicture}{\begin{adjustbox}{max width=\linewidth}}
\AfterEndEnvironment{tikzpicture}{\end{adjustbox}}
% légendes pgfplots lisibles par-dessus les courbes
\pgfplotsset{every axis/.append style={legend style={fill=white,fill opacity=0.92,text opacity=1,draw=black!15,inner sep=3pt}}}
% étiquettes d'arêtes (syntaxe edge["..."]) avec fond blanc
\tikzset{every edge quotes/.append style={fill=white,inner sep=1.2pt}}
%% FIGFIX-END
\begin{document}

\pagegarde{Les lettres de motivation et d'encouragement}{Allemand}{1ère A}{Chapitre 4 : Média et communication}{ALA19}{2 h}{Cette leçon mixte permet de comprendre et d'exploiter un texte sur les lettres de motivation et d'encouragement. Elle développe le vocabulaire de la candidature, la structure de la lettre formelle et informelle, ainsi que les formes de politesse avec « möchte » et « würde ».
\vspace{4pt}
\begin{multicols}{2}\small
\begin{itemize}
\item À la fin de cette leçon, je sais comprendre globalement puis en détail une lettre de motivation ou d'encouragement en allemand.
\item Je sais identifier l'expéditeur, le destinataire, l'objet et l'intention d'une lettre formelle.
\item Je sais employer correctement le vocabulaire de la candidature (die Bewerbung, der Lebenslauf, die Stelle, sich bewerben...) avec articles et pluriels.
\item Je sais distinguer la structure d'une lettre formelle et d'une lettre informelle (Anrede, Grußformel).
\item Je sais conjuguer et employer « möchte » et « würde » pour exprimer la politesse.
\item Je sais comparer les usages de politesse en français et en allemand et éviter les calques.
\end{itemize}\end{multicols}}

\begin{sommaire}
\begin{multicols}{2}
\begin{enumerate}
\item Texte support : Eine Bewerbung und ein Ermutigungsbrief
\item Lexique thématique : die Bewerbung
\item Grammaire 1 : la lettre formelle et informelle
\item Grammaire 2 : la politesse avec « möchte » et « würde »
\item Méthode du commentaire et commentaire modèle
\item Production guidée et entraînement à la traduction
\item Activité d'intégration
\item Méthodes et exercices résolus
\item Exercices
\item Corrigés des exercices
\item Fiche bilan
\end{enumerate}
\end{multicols}
\end{sommaire}

\begin{objectifs}
\begin{itemize}
\item À la fin de cette leçon, je sais comprendre globalement puis en détail une lettre de motivation ou d'encouragement en allemand.
\item Je sais identifier l'expéditeur, le destinataire, l'objet et l'intention d'une lettre formelle.
\item Je sais employer correctement le vocabulaire de la candidature (die Bewerbung, der Lebenslauf, die Stelle, sich bewerben...) avec articles et pluriels.
\item Je sais distinguer la structure d'une lettre formelle et d'une lettre informelle (Anrede, Grußformel).
\item Je sais conjuguer et employer « möchte » et « würde » pour exprimer la politesse.
\item Je sais comparer les usages de politesse en français et en allemand et éviter les calques.
\item Je sais répondre à des questions de compréhension et résumer une lettre en quelques phrases.
\item Je sais rédiger une courte lettre de motivation ou d'encouragement correcte et polie.
\item Je sais traduire des phrases de thème et de version utilisant les formes de politesse.
\end{itemize}
\end{objectifs}

\begin{prerequis}
\begin{itemize}
\item Le présent et le prétérit des verbes modaux (können, müssen, wollen) vus en Seconde.
\item La place du verbe conjugué en position 2 et du verbe à l'infinitif en fin de phrase.
\item Les pronoms personnels et possessifs (ich, mein, Ihr, Sie).
\item Le vocabulaire de base de l'école et des études (die Schule, das Studium, der Beruf).
\item La distinction du/tu/vous : du, ihr, Sie.
\end{itemize}
\end{prerequis}

\newpage

\coursec{Texte support : Eine Bewerbung und ein Ermutigungsbrief}

\courssub{Situation d'accroche et problématique}

Imagine : tu viens d'obtenir ton baccalauréat à Yaoundé et tu rêves d'une bourse pour étudier à Munich, à Vienne ou à Zurich. Sur le site de l'université, on te demande un \all{Bewerbungsschreiben} et un \all{Lebenslauf} en allemand. Comment commencer ? \all{Sehr geehrte Damen und Herren} ou \all{Liebe Frau Müller} ? Et comment dire poliment « je voudrais » sans paraître impoli ? En Allemagne, une lettre mal structurée ou un ton trop direct peut coûter une place.

\textbf{Problématique :} comment une lettre de motivation et une lettre d'encouragement sont-elles construites en allemand, et quelles formules permettent d'exprimer la politesse et le soutien ?

Dans cette première partie, nous allons lire deux lettres authentiques : la lettre de motivation de Merveille, élève camerounaise qui demande un stage dans une ONG à Bonn, puis la lettre d'encouragement que lui écrit son amie Amina. Nous apprendrons à repérer toutes les parties d'une lettre allemande et à comprendre le texte globalement puis en détail.

\courssub{Lecture globale : la lettre de motivation de Merveille}

\begin{textebox}[Eine Bewerbung — Merveille Ngo Bell]
Merveille Ngo Bell\\
Quartier Nlongkak\\
Yaoundé, Kamerun\\[0.5em]
Hilfswerk „Brücke der Hoffnung“ e.V.\\
Abteilung Praktika\\
Bonn, Deutschland\\[0.5em]
Yaoundé, den 12. März 2025\\[0.5em]
\textbf{Bewerbung um ein Praktikum}\\[0.5em]
Sehr geehrte Damen und Herren,\\[0.5em]
mit großem Interesse habe ich Ihre Anzeige auf Ihrer Internetseite gelesen. Ich möchte mich um das Praktikum in Ihrer Organisation bewerben.\\[0.5em]
Zurzeit besuche ich die Klasse Première A am Gymnasium in Yaoundé. Ich lerne seit vier Jahren Deutsch und ich spreche auch Französisch und Englisch. In der Schule bin ich die Vorsitzende des Deutschklubs. Dort organisiere ich Spiele, Lieder und kleine Theaterstücke auf Deutsch. Diese Erfahrung hat mir gezeigt, wie wichtig Teamarbeit und Kommunikation sind.\\[0.5em]
Ich interessiere mich sehr für Ihre Projekte in Kamerun, besonders für das Schulprojekt in der Region West. Ich bin motiviert, pünktlich und ich arbeite gern mit Kindern und Jugendlichen.\\[0.5em]
Ich würde mich sehr freuen, von Ihnen zu hören. Für ein Gespräch stehe ich gern zur Verfügung. Meinen Lebenslauf finden Sie im Anhang.\\[0.5em]
Mit freundlichen Grüßen\\
Merveille Ngo Bell
\end{textebox}

\textbf{Glossaire (marge de lecture) :}

\begin{center}
\begin{tabularx}{\linewidth}{|l|X|}
\hline
\rowcolor{popblueL} \textbf{Deutsch} & \textbf{Français} \\
\hline
die Bewerbung, die Bewerbungen & la candidature \\
\hline
der Lebenslauf, die Lebensläufe & le curriculum vitae (CV) \\
\hline
die Stelle, die Stellen & le poste, l'emploi \\
\hline
sich bewerben (bewirbt, bewarb, hat beworben) & poser sa candidature, postuler \\
\hline
die Erfahrung, die Erfahrungen & l'expérience \\
\hline
die Ermutigung, die Ermutigungen & l'encouragement \\
\hline
die Anzeige, die Anzeigen & l'annonce \\
\hline
das Praktikum, die Praktika & le stage \\
\hline
der Anhang, die Anhänge & la pièce jointe \\
\hline
zur Verfügung stehen & être à disposition \\
\hline
\end{tabularx}
\end{center}

\begin{attention}[Prononciation]
\all{Bewerbung} se prononce « bé-vèr-boungk » (le \all{w} allemand se dit [v], le \all{g} final se durcit en [k]). \all{möchte} se prononce « meu-r'té » (ö comme dans « œuf », ch doux [ç]). \all{würde} se prononce « vur-dé » (ü comme le « u » français « tu »). À l'oral, une intonation montante douce sur \all{Ich würde mich sehr freuen...} marque la politesse : ne lis jamais cette phrase d'un ton sec ou descendant.
\end{attention}

\courssub{Lecture globale : la lettre d'encouragement d'Amina}

\begin{textebox}[Ein Ermutigungsbrief — Amina an Merveille]
Yaoundé, den 15. März 2025\\[0.5em]
Liebe Merveille,\\[0.5em]
ich habe gehört, dass du dich um das Praktikum in Bonn bewirbst. Das ist eine tolle Idee! Du sprichst gut Deutsch, du bist die Vorsitzende des Deutschklubs und du bist sehr motiviert.\\[0.5em]
Ich weiß, dass du ein bisschen nervös bist. Aber keine Angst! Deine Bewerbung ist sehr gut, und dein Lebenslauf ist klar und vollständig. Ich bin sicher, dass du die Stelle bekommst.\\[0.5em]
Und wenn es diesmal nicht klappt, dann probieren wir es zusammen noch einmal. Du bist stark, und ich glaube an dich.\\[0.5em]
Schreib mir bald und erzähl mir alles!\\[0.5em]
Viele liebe Grüße\\
Amina
\end{textebox}

\textbf{Glossaire :} \all{nervös} = nerveux ; \all{klappen} = marcher, réussir (familier) ; \all{probieren} = essayer ; \all{an jemanden glauben} = croire en quelqu'un ; \all{erzählen} = raconter.

\courssub{Observation : les deux lettres sont-elles identiques ?}

Avant toute règle, observons. Compare les deux textes : que remarques-tu ?

\begin{center}
\begin{tabularx}{\linewidth}{|X|X|X|}
\hline
\rowcolor{popblueL} \textbf{Élément} & \textbf{Lettre de Merveille (formelle)} & \textbf{Lettre d'Amina (informelle)} \\
\hline
Anrede & \all{Sehr geehrte Damen und Herren,} & \all{Liebe Merveille,} \\
\hline
Pronom & \all{Sie / Ihre} (vouvoiement) & \all{du / dein} (tutoiement) \\
\hline
Ton & poli, distant, objectif & chaleureux, direct, affectif \\
\hline
Grußformel & \all{Mit freundlichen Grüßen} & \all{Viele liebe Grüße} \\
\hline
Objectif & obtenir un stage & encourager une amie \\
\hline
\end{tabularx}
\end{center}

\begin{definition}[Les deux registres de la lettre]
La \cle{lettre formelle} (\all{der formelle Brief}) s'adresse à une institution, à une personne inconnue ou à un supérieur : on vouvoie (\all{Sie} avec majuscule), l'\cle{Anrede} est \all{Sehr geehrte Damen und Herren} ou \all{Sehr geehrte Frau Müller}, et la \cle{Grußformel} finale est \all{Mit freundlichen Grüßen} (souvent abrégé \all{MfG} dans les e-mails). La \cle{lettre informelle} (\all{der persönliche Brief}) s'adresse à la famille et aux amis : on tutoie (\all{du}), l'Anrede est \all{Liebe Merveille} / \all{Lieber Paul}, et la formule finale est \all{Viele liebe Grüße} ou \all{Herzliche Grüße}.
\end{definition}

\begin{attention}[Majuscule obligatoire]
En allemand, \all{Sie} (vous) et \all{Ihre} (votre) prennent \textbf{toujours} la majuscule dans une lettre, même au milieu de la phrase. Écrire \all{ich habe ihre Anzeige gelesen} (avec minuscule) est une faute grave de politesse : cela signifie « j'ai lu \emph{son} annonce » (à une tierce personne). Pour le tutoiement, la majuscule à \all{Du} / \all{Dein} est aujourd'hui facultative (Duden), mais elle reste appréciée dans les lettres privées pour marquer le respect.
\end{attention}

\courssub{Repérage : l'anatomie d'une lettre allemande}

Une lettre allemande, formelle ou informelle, suit une architecture fixe (norme DIN 5008). Apprends le nom et l'ordre de chaque partie :

\begin{itemize}
\item \cle{der Absender} : l'expéditeur (celui qui écrit), en haut à gauche (nom, rue, code postal, ville).
\item \cle{der Empfänger} : le destinataire (celui qui reçoit), sous l'expéditeur (nom, rue, code postal, ville).
\item \cle{Ort und Datum} : le lieu et la date, à droite, séparés par une virgule. Modèle : \all{Yaoundé, den 12. März 2025}.
\item \cle{der Betreff} : l'objet de la lettre, en gras, \textbf{sans verbe conjugué} : \all{Bewerbung um ein Praktikum}.
\item \cle{die Anrede} : la formule d'appel, suivie d'une virgule ; la phrase suivante commence par une \textbf{minuscule}.
\item \cle{der Textkörper} : le corps du texte (motivation, compétences, demande).
\item \cle{die Grußformel} : la formule de politesse finale, \textbf{sans virgule} après elle.
\item \cle{die Unterschrift} : la signature (prénom + nom).
\end{itemize}

\begin{popfigure}
\begin{tikzpicture}[
  part/.style={draw=popblue, fill=popblueL, rounded corners=2pt, align=left, text width=6.2cm, inner sep=4pt, font=\small},
  lab/.style={draw=poporange, fill=poporangeL, rounded corners=2pt, align=center, text width=3.4cm, inner sep=3pt, font=\small\bfseries},
  arr/.style={-{Stealth[length=2.5mm]}, thick, popdark}
]
\node[part, align=center] (abs) at (0,0){Merveille Ngo Bell\\ Quartier Nlongkak\\ Yaoundé, Kamerun};
\node[part, align=center] (emp) at (0,-1.8){Hilfswerk „Brücke der Hoffnung“ e.V.\\ Abteilung Praktika\\ Bonn, Deutschland};
\node[part, text width=4.5cm] (dat) at (2.2,-3.3) {Yaoundé, den 12. März 2025};
\node[part] (bet) at (0,-4.3) {\textbf{Bewerbung um ein Praktikum}};
\node[part] (anr) at (0,-5.3) {Sehr geehrte Damen und Herren,};
\node[part] (kor) at (0,-6.6) {mit großem Interesse habe ich Ihre Anzeige gelesen. Ich möchte mich um das Praktikum bewerben...};
\node[part] (gru) at (0,-8.1) {Mit freundlichen Grüßen};
\node[part] (unt) at (0,-9.0) {Merveille Ngo Bell};
\node[lab, right=1.6cm of abs] (labs) {der Absender};
\node[lab, right=1.6cm of emp] (lemp) {der Empfänger};
\node[lab, right=1.6cm of dat] (ldat) {Ort und Datum};
\node[lab, right=1.6cm of bet] (lbet) {der Betreff};
\node[lab, right=1.6cm of anr] (lanr) {die Anrede};
\node[lab, right=1.6cm of kor] (lkor) {der Textkörper};
\node[lab, right=1.6cm of gru] (lgru) {die Grußformel};
\node[lab, right=1.6cm of unt] (lunt) {die Unterschrift};
\draw[arr] (labs.west) -- (abs.east);
\draw[arr] (lemp.west) -- (emp.east);
\draw[arr] (ldat.west) -- (dat.east);
\draw[arr] (lbet.west) -- (bet.east);
\draw[arr] (lanr.west) -- (anr.east);
\draw[arr] (lkor.west) -- (kor.east);
\draw[arr] (lgru.west) -- (gru.east);
\draw[arr] (lunt.west) -- (unt.east);
\end{tikzpicture}
\legende{Les huit parties d'une lettre formelle allemande, de l'Absender à l'Unterschrift.}
\end{popfigure}

\begin{aretenir}[La date allemande]
La date s'écrit : ville + virgule + \all{den} + jour avec point + mois + année : \all{Yaoundé, den 12. März 2025}. Le point après le jour correspond à une terminaison ordinale (« le 12\ieme{} »). On utilise \all{den} (accusatif) car il répond implicitement à la question « Wann schreibst du? » (quand écris-tu?) → « Am 12. März » (datif) ou « Den 12. März » (accusatif adverbial de temps). On n'écrit jamais \all{am 12. März} dans l'en-tête de la lettre.
\end{aretenir}

\courssub{Compréhension détaillée : que dit vraiment Merveille ?}

Relisons le corps de la lettre de motivation (\all{Textkörper}). Il suit un plan en quatre mouvements, typique du \all{Bewerbungsschreiben} :

\begin{enumerate}
\item \textbf{Le déclencheur (der Anlass) :} \all{mit großem Interesse habe ich Ihre Anzeige auf Ihrer Internetseite gelesen} — « c'est avec un grand intérêt que j'ai lu votre annonce sur votre site internet ». La lettre commence toujours par dire \emph{où} l'on a vu l'offre. Structure : complément circonstanciel (\all{mit...}) + verbe (\all{habe}) + sujet (\all{ich}) + objet (\all{Ihre Anzeige}) + participe (\all{gelesen}).
\item \textbf{La demande (der Bewerbungswunsch) :} \all{Ich möchte mich um das Praktikum in Ihrer Organisation bewerben.} — « je voudrais poser ma candidature pour le stage dans votre organisation ». Note la construction : \all{sich bewerben um} + accusatif = postuler \emph{pour} quelque chose. \all{möchte} (Konjunktiv II de \all{mögen}) exprime la politesse.
\item \textbf{Les compétences (die Qualifikationen) :} Merveille présente son niveau (\all{ich lerne seit vier Jahren Deutsch} — \all{seit} + datif), ses langues et son expérience (\all{Vorsitzende des Deutschklubs} — génitif). Elle relie chaque qualité au poste : \all{Diese Erfahrung hat mir gezeigt, wie wichtig Teamarbeit und Kommunikation sind.} (subordonnée \all{wie}... verbe à la fin).
\item \textbf{La conclusion polie (der Schlusssatz) :} \all{Ich würde mich sehr freuen, von Ihnen zu hören.} — « je serais très heureuse d'avoir de vos nouvelles ». C'est LA formule de clôture standard d'une candidature : elle exprime l'espoir sans exiger. \all{würde} (Konjunktiv II de \all{werden}) + infinitif. \all{Für ein Gespräch stehe ich gern zur Verfügung.} (inversion sujet-verbe après complément prépositionnel). \all{Meinen Lebenslauf finden Sie im Anhang.} (topicalisation de l'objet accusatif).
\end{enumerate}

\begin{propriete}[sich bewerben : un verbe fort réfléchi avec préposition]
\all{sich bewerben um + Akkusativ} est un verbe fort (irrégulier, comme \all{werben}) : 
Présent : \all{ich bewerbe mich, du bewirbst dich, er/sie/es bewirbt sich, wir bewerben uns, ihr bewerbt euch, sie/Sie bewerben sich} ; 
Präteritum : \all{ich bewarb mich} ; 
Perfekt : \all{ich habe mich beworben} (pas de \all{ge-} car préfixe \all{be-}). 
Exemples : \all{Merveille bewirbt sich um ein Praktikum.} « Merveille postule pour un stage. » \all{Sie hat sich um die Stelle beworben.} « Elle a posé sa candidature pour le poste. » 
\textbf{Attention :} on ne dit jamais \all{bewerben für} — la préposition correcte est \all{um}.
\end{propriete}

\begin{attention}[Faux amis et pièges fréquents]
\begin{itemize}
\item \all{die Stelle} a deux sens : « le poste, l'emploi » (\all{die Stelle bekommen / erhalten} = obtenir le poste) et « l'endroit, le passage » (\all{an dieser Stelle} = à cet endroit). Dans une \all{Bewerbung}, \all{Stelle} signifie toujours « poste ».
\item \all{bekommen} ne veut pas dire « devenir » (\all{werden}) mais « recevoir, obtenir » : \all{du bekommst die Stelle} = « tu obtiens le poste ».
\item \all{sich interessieren für} + accusatif (pas \all{an} comme en anglais \emph{interested in}) : \all{Ich interessiere mich für Ihre Projekte.}
\item \all{seit} + datif (temps écoulé depuis un point passé) : \all{seit vier Jahren} (pas \all{vor}).
\end{itemize}
\end{attention}

\courssub{La lettre d'Amina : les marques de l'encouragement}

La lettre informelle d'Amina utilise des stratégies typiques de l'encouragement (\all{die Ermutigung}) :

\begin{itemize}
\item \textbf{Valoriser :} \all{Das ist eine tolle Idee!} « C'est une super idée ! » ; \all{Du bist stark.} « Tu es forte. »
\item \textbf{Rassurer :} \all{Aber keine Angst!} « Mais pas de panique ! » (ellipse de \all{hab}).
\item \textbf{Exprimer la confiance :} \all{Ich bin sicher, dass du die Stelle bekommst.} « Je suis sûre que tu obtiendras le poste. » Note la subordonnée avec \all{dass} : le verbe conjugué (\all{bekommst}) va \textbf{à la fin}.
\item \textbf{Promettre son soutien :} \all{dann probieren wir es zusammen noch einmal} « alors nous réessaierons ensemble ». \all{Ich glaube an dich.} (verbe \all{glauben an} + accusatif).
\end{itemize}

\begin{exemplebox}[Formules d'encouragement à réutiliser]
\all{Ich glaube an dich.} « Je crois en toi. »\\
\all{Du schaffst das!} « Tu vas y arriver ! » (\all{schaffen} = réussir, gérer).\\
\all{Kopf hoch!} « Courage, la tête haute ! » (idiome).\\
\all{Ich drücke dir die Daumen.} « Je croise les doigts pour toi. » (litt. « je te presse les pouces »).
\end{exemplebox}

\begin{savaistu}[La candidature à l'allemande : die Bewerbungsmappe]
En Allemagne, en Autriche et en Suisse, le dossier de candidature complet s'appelle \all{die Bewerbungsmappe}. Traditionnellement, on y joint une \textbf{photo} professionnelle sur le \all{Lebenslauf} (bien que cela devienne facultatif avec la loi anti-discrimination \all{AGG}), et la lettre de motivation (\all{Anschreiben}) doit tenir sur \textbf{une seule page A4} : les employeurs germanophones apprécient la concision, la clarté et l'absence de fautes. Au Cameroun, les élèves germanistes qui postulent à des bourses du DAAD (Deutscher Akademischer Austauschdienst, très actif à Yaoundé) ou à des stages au Goethe-Institut doivent maîtriser exactement ce formalisme : c'est un savoir-faire qui ouvre de vraies portes académiques et professionnelles !
\end{savaistu}

\courssub{Méthode : lire et comprendre une lettre allemande}

\begin{methode}[Quatre questions pour lire une lettre]
\begin{enumerate}
\item \textbf{Qui écrit ?} Regarde l'\cle{Absender} (en haut à gauche) et la signature (\cle{Unterschrift}).
\item \textbf{À qui ?} Regarde l'\cle{Empfänger} et l'\cle{Anrede} : \all{Sehr geehrte Damen und Herren} / \all{Sehr geehrte Frau Müller} = destinataire formel (vouvoiement \all{Sie}) ; \all{Liebe Merveille} / \all{Lieber Paul} = proche (tutoiement \all{du}).
\item \textbf{Quel est l'objectif ?} Le \cle{Betreff} et la première phrase le donnent : postuler (\all{sich bewerben um}), remercier (\all{sich bedanken für}), encourager (\all{ermutigen}), inviter (\all{einladen})...
\item \textbf{Quelles formules de politesse ?} Repère la \cle{Grußformel} et les tournures comme \all{Ich möchte...} / \all{Ich würde mich freuen...} / \all{Könnten Sie bitte...} : elles indiquent le degré de formalité et le ton de la lettre.
\end{enumerate}
\end{methode}

\courssub{Questions de compréhension}

\begin{exercice}[Verstehst du den Text?]
Réponds en allemand, par des phrases complètes :
\begin{enumerate}
\item Wer schreibt die erste E-Mail bzw. den ersten Brief?
\item An wen schreibt Merveille?
\item Warum schreibt sie? Was möchte sie?
\item Seit wann lernt Merveille Deutsch?
\item Welche Erfahrung hat sie in der Schule?
\item Wer schreibt den zweiten Brief? Ist der Brief formell oder informell?
\item Wie ermutigt Amina ihre Freundin? Nenne zwei Sätze.
\end{enumerate}
\end{exercice}

\begin{exoresolu}[Frage 3 — Warum schreibt Merveille?]
Énoncé : réponds à la question \all{Warum schreibt Merveille? Was möchte sie?} par une phrase complète en allemand, puis traduis ta réponse en français.
\tcblower
\textbf{Solution détaillée :} on cherche dans le texte la phrase qui exprime l'objectif : \all{Ich möchte mich um das Praktikum in Ihrer Organisation bewerben.} On la reformule à la troisième personne, car la question porte sur Merveille (elle) : attention à la conjugaison (\all{möchte} reste identique à la 3\ieme{} personne du singulier) et au pronom réfléchi (\all{sich}).\\
Réponse : \all{Merveille schreibt an die Organisation in Bonn, weil sie sich um das Praktikum bewerben möchte.}\\
Traduction : « Merveille écrit à l'organisation de Bonn parce qu'elle voudrait poser sa candidature pour le stage. »\\
Points de grammaire vérifiés : subordonnée avec \all{weil} → verbe conjugué (\all{möchte}) renvoyé à la fin, l'infinitif \all{bewerben} juste avant lui ; \all{sich bewerben um} + accusatif (\all{das Praktikum}).
\end{exoresolu}

\begin{corrige}[Exercice « Verstehst du den Text? »]
\begin{enumerate}
\item \all{Merveille Ngo Bell schreibt den ersten Brief.} (Merveille Ngo Bell écrit la première lettre.)
\item \all{Sie schreibt an das Hilfswerk „Brücke der Hoffnung“ in Bonn.} (Elle écrit à l'organisation caritative « Pont de l'Espoir » à Bonn.)
\item \all{Sie möchte sich um ein Praktikum bewerben.} (Elle voudrait postuler pour un stage.)
\item \all{Sie lernt seit vier Jahren Deutsch.} (Elle apprend l'allemand depuis quatre ans.)
\item \all{Sie ist die Vorsitzende des Deutschklubs und organisiert Spiele, Lieder und Theaterstücke.} (Elle est présidente du club d'allemand et organise jeux, chants et pièces de théâtre.)
\item \all{Amina schreibt den zweiten Brief. Er ist informell, weil sie „Liebe Merveille“ und „du“ schreibt.} (Amina écrit la deuxième lettre. Elle est informelle car elle utilise « Chère Merveille » et le tutoiement.)
\item \all{Sie sagt: „Das ist eine tolle Idee!“ und „Ich bin sicher, dass du die Stelle bekommst.“} (Elle dit : « C'est une super idée ! » et « Je suis sûre que tu obtiendras le poste. »)
\end{enumerate}
\end{corrige}

\courssub{Lecture orale : entraînement à la prononciation}

\begin{experience}[Lesen und nachsprechen — Lecture à voix haute]
Travaille à deux. L'élève A lit la lettre de Merveille à voix haute avec une intonation polie, posée et professionnelle ; l'élève B lit la lettre d'Amina avec un ton chaleureux, dynamique et bienveillant. Veillez particulièrement aux points suivants :
\begin{itemize}
\item \all{Bewerbung} [bəˈvɛʁbʊŋk] : « bé-vèr-boungk » (\all{w}=[v], \all{g} final=[k]) ;
\item \all{möchte} [ˈmœçtə] : « meu-r'té » (ö=[œ], ch=[ç] doux) ;
\item \all{würde} [ˈvʏʁdə] : « vur-dé » (ü=[ʏ], r roulé ou vocalisé) ;
\item \all{Sehr geehrte Damen und Herren} [zeːɐ̯ ɡəˈeːʁtə ˈdaːmən ʊnt ˈhɛʁən] : « zé-r gué-èr-té da-mèn ount hé-rèn » ;
\item \all{keine Angst} [ˈkaɪnə ʔaŋst] : « kaine angst » (diphthongue \all{ei}=[aɪ], \all{ng}=[ŋ]) ;
\item \all{Ich drücke dir die Daumen} [ɪç ˈdʁʏkə diːɐ̯ diː ˈdaʊmən] : « ik drücke dir di daou-men ».
\end{itemize}
Puis échangez les rôles. Le lecteur de la lettre formelle doit « sonner » comme une candidate sérieuse ; le lecteur de la lettre d'Amina doit « sonner » comme une vraie amie.
\end{experience}

\begin{aretenir}[Bilan de la section]
Une lettre allemande se construit en huit parties fixes : \all{Absender – Empfänger – Ort und Datum – Betreff – Anrede – Textkörper – Grußformel – Unterschrift}. La lettre formelle vouvoie (\all{Sehr geehrte Damen und Herren} ... \all{Mit freundlichen Grüßen}) ; la lettre informelle tutoie (\all{Liebe Merveille} ... \all{Viele liebe Grüße}). La candidature suit quatre mouvements : annonce de l'offre, demande avec \all{sich bewerben um + Akkusativ}, compétences, conclusion polie avec \all{Ich würde mich sehr freuen...}. L'encouragement valorise (\all{tolle Idee}), rassure (\all{keine Angst}), exprime la confiance (\all{Ich bin sicher, dass...}) et promet le soutien (\all{probieren wir es zusammen}). Le Konjunktiv II (\all{möchte, würde}) est la clé de la politesse formelle.
\end{aretenir}

\coursec{Lexique thématique : die Bewerbung}

Dans la section précédente, nous avons lu et commenté le texte \emph{Eine Bewerbung und ein Ermutigungsbrief} : Amina, élève de terminale à Yaoundé, écrit une lettre de candidature pour un stage, et son amie lui répond par une lettre d'encouragement. Ce texte contenait beaucoup de mots nouveaux : \all{die Bewerbung}, \all{der Lebenslauf}, \all{sich bewerben}, \all{die Grußformel}... Cette deuxième section a pour but de \cle{systématiser ce vocabulaire} : nous allons l'observer, le classer en champs lexicaux, apprendre ses articles et ses pluriels, découvrir les familles de mots, puis nous entraîner à l'employer dans des phrases. Un bon lexique, c'est la moitié de la réussite à l'examen : sans les mots justes, impossible de comprendre un texte ni de rédiger une lettre.

\courssub{Observer avant d'apprendre : les mots en contexte}

Avant tout tableau, lisons ce court texte qui réutilise presque tout le vocabulaire de la leçon. Soulignez mentalement les mots que vous reconnaissez déjà grâce au texte support.

\begin{textebox}[Aminas Bewerbung — la candidature d'Amina]
Amina Ngo Bell ist Schülerin in Yaoundé. In der Zeitung liest sie eine Anzeige: Ein deutsches Reisebüro in Douala sucht eine Praktikantin für die Sommerferien. Amina ist sofort motiviert. Sie schreibt einen Lebenslauf und ein Bewerbungsschreiben. In der Anrede schreibt sie: „Sehr geehrte Damen und Herren“. Am Ende steht die Grußformel: „Mit freundlichen Grüßen“, und darunter ihre Unterschrift.

Amina bewirbt sich um die Stelle als Praktikantin. Sie schickt die Bewerbung per E-Mail an den Arbeitgeber. Nach einer Woche antwortet der Chef: „Wir laden Sie zu einem Gespräch ein.“ Amina freut sich sehr über die Antwort. Ihre Freundin Chantal ermutigt sie: „Du bist fleißig, pünktlich und freundlich. Du schaffst das!“ Amina dankt ihrer Freundin für die netten Worte.

Beim Gespräch fragt der Chef: „Warum möchten Sie bei uns arbeiten?“ Amina antwortet höflich: „Ich spreche gut Deutsch, ich bin zuverlässig und ich habe schon Erfahrung mit Touristen.“ Der Chef lächelt: „Sie sind die beste Bewerberin. Die Stelle ist für Sie!“ Amina freut sich schon auf den Praktikumsplatz.
\end{textebox}

\textbf{Glossaire (ordre d'apparition) :}

\begin{itemize}
\item \all{die Anzeige (-n)} : l'annonce (dans un journal, sur Internet)
\item \all{das Reisebüro (-s)} : l'agence de voyages
\item \all{der Lebenslauf („-e)} : le CV
\item \all{das Bewerbungsschreiben (-)} : la lettre de motivation
\item \all{die Anrede (-n)} : la formule d'appel
\item \all{die Grußformel (-n)} : la formule de politesse finale
\item \all{die Unterschrift (-en)} : la signature
\item \all{sich bewerben um + Akk.} : postuler à, poser sa candidature pour
\item \all{der Arbeitgeber (-)} : l'employeur
\item \all{antworten} : répondre
\item \all{ermutigen} : encourager
\item \all{zuverlässig} : fiable, digne de confiance
\item \all{die Erfahrung (-en)} : l'expérience
\item \all{die Bewerberin (-nen)} : la candidate
\item \all{der Praktikumsplatz („-e)} : la place de stage
\end{itemize}

\textbf{Questions de compréhension (répondez en allemand, phrases complètes) :}

\begin{enumerate}
\item \all{Was liest Amina in der Zeitung?}
\item \all{Welche Dokumente schreibt sie?}
\item \all{Wie schickt sie die Bewerbung?}
\item \all{Wer ermutigt Amina?}
\item \all{Warum bekommt Amina die Stelle?}
\end{enumerate}

\courssub{Champ 1 — La candidature et les documents}

\begin{definition}[Le vocabulaire de la candidature]
\begin{itemize}
\item \all{die Bewerbung (-en)} : la candidature
\item \all{der Lebenslauf („-e)} : le CV (curriculum vitae)
\item \all{das Bewerbungsschreiben (-)} : la lettre de motivation
\item \all{die Anzeige (-n)} : l'annonce, l'offre d'emploi
\item \all{die Stelle (-n)} : le poste, l'emploi
\item \all{der Praktikumsplatz („-e)} : la place de stage
\end{itemize}
\end{definition}

Observez bien la notation entre parenthèses : \all{die Bewerbung (-en)} signifie que le pluriel est \all{die Bewerbungen} ; \all{der Lebenslauf („-e)} signifie que le pluriel est \all{die Lebensläufe} (avec Umlaut sur le \all{a}). En allemand, \textbf{on apprend toujours un nom avec son article et son pluriel}, jamais seul. C'est la règle d'or du vocabulaire allemand.

\begin{exemplebox}[Les documents en action]
\all{Amina liest eine Anzeige in der Zeitung.} « Amina lit une annonce dans le journal. »

\all{Sie schreibt ihren Lebenslauf und ein Bewerbungsschreiben.} « Elle rédige son CV et une lettre de motivation. »

\all{Die Stelle als Praktikantin ist interessant.} « Le poste de stagiaire est intéressant. »

\all{Er sucht einen Praktikumsplatz in Douala.} « Il cherche une place de stage à Douala. »
\end{exemplebox}

\begin{attention}[der Brief ou das Schreiben ?]
Le français n'a qu'un mot, « lettre », mais l'allemand en a deux : \all{der Brief (-e)} est la lettre en général (lettre à un ami, lettre personnelle), tandis que \all{das Schreiben (-)} désigne un courrier officiel, administratif. On dit donc \all{das Bewerbungsschreiben} (lettre de motivation, courrier officiel) mais \all{ein Brief an meine Freundin} (une lettre à mon amie). Autre piège : ne traduisez pas « candidature » par \all{die Kandidatur} dans ce contexte scolaire : le mot courant est \all{die Bewerbung}.
\end{attention}

\courssub{Champ 2 — Les personnes}

\begin{definition}[Les acteurs de la candidature]
\begin{itemize}
\item \all{der Bewerber (-) / die Bewerberin (-nen)} : le candidat / la candidate
\item \all{der Arbeitgeber (-)} : l'employeur (celui qui donne le travail)
\item \all{der Arbeitnehmer (-)} : l'employé, le salarié (celui qui reçoit le travail)
\item \all{der Chef (-s) / die Chefin (-nen)} : le chef / la cheffe, le patron
\item \all{der Praktikant (-en) / die Praktikantin (-nen)} : le stagiaire / la stagiaire
\end{itemize}
\end{definition}

\begin{propriete}[Le féminin avec -in]
En allemand, la plupart des noms de personnes forment leur féminin en ajoutant \cle{-in} au masculin ; le pluriel féminin est alors en \cle{-nen} :
\begin{center}
\begin{tabular}{|l|l|l|}
\hline
\rowcolor{popblueL} Masculin & Féminin & Pluriel féminin \\
\hline
der Bewerber & die Bewerberin & die Bewerberinnen \\
\hline
der Praktikant & die Praktikantin & die Praktikantinnen \\
\hline
der Chef & die Chefin & die Chefinnen \\
\hline
der Schüler & die Schülerin & die Schülerinnen \\
\hline
\end{tabular}
\end{center}
Attention au pluriel de \all{der Chef} : \all{die Chefs} (avec un simple \all{-s}, comme beaucoup de mots d'origine étrangère).
\end{propriete}

\begin{exemplebox}
\all{Der Bewerber schreibt einen Lebenslauf.} « Le candidat rédige un CV. »

\all{Die Bewerberin ist sehr motiviert.} « La candidate est très motivée. »

\all{Der Arbeitgeber liest die Bewerbung.} « L'employeur lit la candidature. »
\end{exemplebox}

\courssub{Champ 3 — Les actions (les verbes du thème)}

\begin{definition}[Les verbes essentiels]
\begin{itemize}
\item \all{sich bewerben um + Akkusativ} : postuler à, poser sa candidature pour
\item \all{schreiben (schrieb, hat geschrieben)} : écrire
\item \all{schicken} : envoyer
\item \all{antworten (auf + Akk.)} : répondre (à)
\item \all{ermutigen} : encourager
\item \all{danken (+ Dativ, für + Akk.)} : remercier (quelqu'un, de quelque chose)
\item \all{sich freuen über + Akk.} : se réjouir de (quelque chose de présent ou passé)
\item \all{sich freuen auf + Akk.} : se réjouir à l'idée de, attendre avec joie (quelque chose de futur)
\end{itemize}
\end{definition}

\begin{propriete}[sich bewerben : un verbe réfléchi irrégulier]
\all{sich bewerben} est un verbe \cle{réfléchi} et \cle{fort} (irrégulier) : \all{bewerben — bewarb — hat beworben}. Au présent : \all{ich bewerbe mich, du bewirbst dich, er/sie/es bewirbt sich, wir bewerben uns, ihr bewerbt euch, sie/Sie bewerben sich}. Le pronom réfléchi est à l'\textbf{accusatif} et se place juste après le verbe conjugué. Il est toujours suivi de la préposition \cle{um + accusatif}.
\end{propriete}

\begin{exemplebox}[Exemples traduits]
\all{Ich bewerbe mich um die Stelle als Praktikantin.} « Je postule au poste de stagiaire. »

\all{Sie ermutigt ihre Freundin.} « Elle encourage son amie. »

\all{Er schickt die Bewerbung per E-Mail.} « Il envoie la candidature par e-mail. »

\all{Der Chef antwortet schnell auf den Brief.} « Le chef répond vite à la lettre. »

\all{Ich danke Ihnen für Ihre Antwort.} « Je vous remercie de votre réponse. »

\all{Wir freuen uns über Ihre Bewerbung.} « Nous nous réjouissons de votre candidature. »

\all{Amina freut sich auf den Praktikumsplatz.} « Amina se réjouit à l'idée de cette place de stage. »
\end{exemplebox}

\begin{attention}[Le piège n° 1 : sich bewerben UM, jamais für]
Les francophones traduisent « postuler pour un poste » par \all{sich bewerben für die Stelle} : c'est \textbf{faux}. En allemand, on dit exclusivement \all{sich bewerben \textbf{um} + Akkusativ} : \all{Ich bewerbe mich um die Stelle.} Retenez aussi que \all{danken} se construit avec le \textbf{datif} de la personne : \all{Ich danke \textbf{dir}} (je te remercie), et non \all{ich danke dich}. Enfin, ne confondez pas \all{sich freuen über} (joie pour quelque chose qui existe déjà : la réponse reçue) et \all{sich freuen auf} (joie anticipée : le stage qui arrive).
\end{attention}

\courssub{Champ 4 — Les qualités (les adjectifs)}

Dans une lettre de motivation, on doit se présenter : voici les six adjectifs indispensables.

\begin{center}
\begin{tabularx}{\linewidth}{|X|X|X|}
\hline
\rowcolor{popblueL} Deutsch & Français & Beispiel \\
\hline
motiviert & motivé & \all{Ich bin sehr motiviert.} \\
\hline
pünktlich & ponctuel & \all{Sie ist immer pünktlich.} \\
\hline
fleißig & travailleur, appliqué & \all{Er ist ein fleißiger Schüler.} \\
\hline
freundlich & aimable, sympathique & \all{Die Chefin ist freundlich.} \\
\hline
zuverlässig & fiable, sérieux & \all{Ein zuverlässiger Mitarbeiter} \\
\hline
erfahren & expérimenté & \all{Sie ist sehr erfahren.} \\
\hline
\end{tabularx}
\end{center}

\begin{exemplebox}
\all{Ich bin motiviert, pünktlich und zuverlässig.} « Je suis motivé, ponctuel et fiable. » — phrase modèle à réutiliser dans vos propres lettres de motivation.
\end{exemplebox}

\courssub{Champ 5 — Les formules de la lettre}

\begin{definition}[Les éléments fixes d'une lettre]
\begin{itemize}
\item \all{die Anrede (-n)} : la formule d'appel (ex. \all{Sehr geehrte Damen und Herren,} / \all{Liebe Chantal,})
\item \all{der Betreff (-e)} : l'objet de la lettre (ex. \all{Bewerbung um einen Praktikumsplatz})
\item \all{die Grußformel (-n)} : la formule de politesse finale (ex. \all{Mit freundlichen Grüßen} / \all{Viele Grüße})
\item \all{die Unterschrift (-en)} : la signature
\end{itemize}
\end{definition}

\begin{exemplebox}[Formelles ou informelles ?]
Lettre formelle : \all{Sehr geehrte Damen und Herren,} ... \all{Mit freundlichen Grüßen} « Mesdames, Messieurs, ... Veuillez agréer mes salutations distinguées. »

Lettre informelle (à un ami) : \all{Lieber Paul,} / \all{Liebe Amina,} ... \all{Viele Grüße} ou \all{Deine Chantal} « Cher Paul, / Chère Amina, ... Amitiés. »
\end{exemplebox}

\courssub{Les familles de mots : un verbe cache souvent un nom}

L'allemand construit des familles entières autour d'une même racine. Apprendre la famille, c'est apprendre trois ou quatre mots pour le prix d'un.

\begin{center}
\begin{tabularx}{\linewidth}{|X|X|X|}
\hline
\rowcolor{popgreenL} Verbe & Nom & Autre membre de la famille \\
\hline
sich bewerben (postuler) & die Bewerbung (la candidature) & der Bewerber / die Bewerberin (le/la candidat(e)) \\
\hline
ermutigen (encourager) & die Ermutigung (l'encouragement) & mutig (courageux) ; der Mut (le courage) \\
\hline
arbeiten (travailler) & die Arbeit (le travail) & der Arbeitgeber (l'employeur) ; der Arbeitnehmer (le salarié) \\
\hline
schreiben (écrire) & das Schreiben (le courrier) & der Brief (la lettre) ; die Unterschrift (la signature) \\
\hline
danken (remercier) & der Dank (les remerciements) & dankbar (reconnaissant) \\
\hline
\end{tabularx}
\end{center}

Observez la logique : \all{der Arbeit\textbf{geber}} est celui qui \emph{donne} (\all{geben}) le travail, \all{der Arbeit\textbf{nehmer}} est celui qui le \emph{prend} (\all{nehmen}). De même, \all{er-mut-igen} signifie littéralement « donner du courage (\all{der Mut}) à quelqu'un » : encourager, c'est donner du courage !

\begin{popfigure}
\begin{center}\tcbox[colback=popink,colframe=popink,coltext=white,arc=9pt,boxsep=4pt,left=6pt,right=6pt]{\bfseries\small die Bewerbung}\end{center}
\vspace{-2pt}
\begin{tcbraster}[raster columns=3,raster equal height=rows,raster column skip=6pt,raster row skip=6pt,enhanced,blankest]
\begin{tcolorbox}[enhanced,colback=white,colframe=poporange!80,boxrule=0.9pt,arc=3pt,title={Personen},fonttitle=\bfseries\scriptsize,coltitle=white,colbacktitle=poporange!80,left=4pt,right=4pt,top=2pt,bottom=2pt,boxsep=1pt,toptitle=1.5pt,bottomtitle=1.5pt,halign=left]
\begin{itemize}[nosep,leftmargin=*,itemsep=1pt,font=\scriptsize]
\item der Bewerber
\item der Arbeitgeber
\item der Praktikant
\end{itemize}
\end{tcolorbox}
\begin{tcolorbox}[enhanced,colback=white,colframe=popgreen!70,boxrule=0.9pt,arc=3pt,title={Dokumente},fonttitle=\bfseries\scriptsize,coltitle=white,colbacktitle=popgreen!70,left=4pt,right=4pt,top=2pt,bottom=2pt,boxsep=1pt,toptitle=1.5pt,bottomtitle=1.5pt,halign=left]
\begin{itemize}[nosep,leftmargin=*,itemsep=1pt,font=\scriptsize]
\item der Lebenslauf
\item das Schreiben
\item die Anzeige
\end{itemize}
\end{tcolorbox}
\begin{tcolorbox}[enhanced,colback=white,colframe=poppurple!70,boxrule=0.9pt,arc=3pt,title={Aktionen},fonttitle=\bfseries\scriptsize,coltitle=white,colbacktitle=poppurple!70,left=4pt,right=4pt,top=2pt,bottom=2pt,boxsep=1pt,toptitle=1.5pt,bottomtitle=1.5pt,halign=left]
\begin{itemize}[nosep,leftmargin=*,itemsep=1pt,font=\scriptsize]
\item sich bewerben um
\item schicken
\item ermutigen
\end{itemize}
\end{tcolorbox}
\begin{tcolorbox}[enhanced,colback=white,colframe=poppink!80,boxrule=0.9pt,arc=3pt,title={Qualitäten},fonttitle=\bfseries\scriptsize,coltitle=white,colbacktitle=poppink!80,left=4pt,right=4pt,top=2pt,bottom=2pt,boxsep=1pt,toptitle=1.5pt,bottomtitle=1.5pt,halign=left]
\begin{itemize}[nosep,leftmargin=*,itemsep=1pt,font=\scriptsize]
\item motiviert
\item pünktlich
\item zuverlässig
\end{itemize}
\end{tcolorbox}
\begin{tcolorbox}[enhanced,colback=white,colframe=popgold!85,boxrule=0.9pt,arc=3pt,title={Formeln},fonttitle=\bfseries\scriptsize,coltitle=white,colbacktitle=popgold!85,left=4pt,right=4pt,top=2pt,bottom=2pt,boxsep=1pt,toptitle=1.5pt,bottomtitle=1.5pt,halign=left]
\begin{itemize}[nosep,leftmargin=*,itemsep=1pt,font=\scriptsize]
\item die Anrede
\item die Grußformel
\item die Unterschrift
\end{itemize}
\end{tcolorbox}
\end{tcbraster}

\legende{Carte mentale du champ lexical de \all{die Bewerbung} : cinq branches à mémoriser.}
\end{popfigure}

\begin{savaistu}[Pourquoi « Lebenslauf » ?]
\all{Der Lebenslauf} signifie littéralement « le cours de la vie » (\all{das Leben} = la vie, \all{der Lauf} = la course, le déroulement) : c'est le résumé du parcours d'une personne. Les Allemands disent aussi très souvent \all{der CV} (prononcé « tsé-fao »), abréviation du latin \emph{curriculum vitae} — exactement comme en français ! Autre curiosité : l'allemand adore les mots composés. \all{das Bewerbungsschreiben} = \all{Bewerbung} + \all{Schreiben} ; \all{der Praktikumsplatz} = \all{Praktikum} + \all{Platz}. Savoir découper ces mots longs, c'est comprendre le vocabulaire sans dictionnaire.
\end{savaistu}

\begin{exoresolu}[Complétez avec le mot qui convient]
Énoncé — Complétez chaque phrase avec le mot correct de la liste : \all{die Stelle, der Lebenslauf, die Grußformel, ermutigt, die Anzeige, sich bewerben}.

\begin{enumerate}
\item Amina liest \dots\ in der Zeitung.
\item Sie möchte \dots\ um die Stelle als Praktikantin.
\item Im \dots\ stehen ihre Schule und ihre Sprachen.
\item Am Ende des Briefes steht \dots\ : „Mit freundlichen Grüßen“.
\item Chantal \dots\ ihre Freundin vor dem Gespräch.
\item Der Chef sagt: „Diese \dots\ ist für Sie!“
\end{enumerate}
\tcblower
Solution détaillée :
\begin{enumerate}
\item \all{Amina liest \textbf{die Anzeige} in der Zeitung.} — On lit une \emph{annonce} dans un journal ; accusatif féminin, l'article reste \all{die}.
\item \all{Sie möchte \textbf{sich bewerben} um die Stelle als Praktikantin.} — Après \all{möchte}, l'infinitif se place en fin de phrase ; le verbe est réfléchi et suivi de \all{um}.
\item \all{Im \dots\ stehen ihre Schule und ihre Sprachen.} → \all{der Lebenslauf} ; attention : \all{im} = \all{in dem} (datif), donc \all{Im \textbf{Lebenslauf}}. C'est dans le CV qu'on indique son école et ses langues.
\item \all{Am Ende des Briefes steht \textbf{die Grußformel}.} — La formule finale de politesse.
\item \all{Chantal \textbf{ermutigt} ihre Freundin vor dem Gespräch.} — Verbe régulier : \all{sie ermutigt}.
\item \all{Der Chef sagt: „Diese \textbf{Stelle} ist für Sie!“} — Le poste, l'emploi.
\end{enumerate}
\end{exoresolu}

\begin{exercice}[À vous de jouer]
\begin{enumerate}
\item Traduisez en allemand : a) Je postule au poste de stagiaire. b) Elle encourage son amie. c) Nous nous réjouissons de votre réponse. d) Il envoie son CV à l'employeur.
\item Donnez le féminin et le pluriel féminin de : \all{der Bewerber, der Praktikant, der Arbeitgeber}.
\item Conjuguez \all{sich bewerben} au présent avec \all{ich}, \all{du}, \all{wir}.
\item Classez ces mots dans les cinq branches de la carte mentale : \all{die Anrede, fleißig, der Chef, schicken, das Bewerbungsschreiben, erfahren, danken, die Unterschrift}.
\end{enumerate}
\end{exercice}

\begin{corrige}[Exercice « À vous de jouer »]
\begin{enumerate}
\item a) \all{Ich bewerbe mich um die Stelle als Praktikant.} b) \all{Sie ermutigt ihre Freundin.} c) \all{Wir freuen uns über Ihre Antwort.} d) \all{Er schickt seinen Lebenslauf an den Arbeitgeber.}
\item \all{die Bewerberin (die Bewerberinnen) ; die Praktikantin (die Praktikantinnen) ; die Arbeitgeberin (die Arbeitgeberinnen)}.
\item \all{ich bewerbe mich ; du bewirbst dich ; wir bewerben uns}.
\item Formeln : \all{die Anrede, die Unterschrift} ; Qualitäten : \all{fleißig, erfahren} ; Personen : \all{der Chef} ; Aktionen : \all{schicken, danken} ; Dokumente : \all{das Bewerbungsschreiben}.
\end{enumerate}
\end{corrige}

\begin{aretenir}[L'essentiel du lexique]
\begin{itemize}
\item Chaque nom s'apprend avec son \cle{article} et son \cle{pluriel} : \all{die Bewerbung (-en)}, \all{der Lebenslauf („-e)}, \all{die Stelle (-n)}.
\item \all{sich bewerben} est réfléchi, irrégulier (\all{bewarb — hat beworben}) et toujours suivi de \cle{um + Akkusativ}.
\item Le féminin des noms de personnes se forme avec \cle{-in} (pluriel \cle{-nen}).
\item Les familles de mots : \all{bewerben → die Bewerbung, der Bewerber} ; \all{ermutigen → die Ermutigung, mutig} ; \all{arbeiten → der Arbeitgeber, der Arbeitnehmer}.
\item \all{der Brief} = lettre personnelle ; \all{das Schreiben} = courrier officiel.
\end{itemize}
\end{aretenir}

Maintenant que le vocabulaire est solide, nous pouvons passer à la grammaire : comment structurer une lettre formelle ou informelle en allemand ? C'est l'objet de la section suivante.

\coursec{Grammaire 1 : la lettre formelle et informelle}

Dans la section précédente, nous avons découvert le vocabulaire de la candidature : \all{die Bewerbung}, \all{der Lebenslauf}, \all{die Stelle}, \all{sich bewerben}. Nous savons maintenant nommer les pièces d'un dossier. Mais une candidature, c'est d'abord une \emph{lettre} : il faut donc savoir la construire correctement. Or l'allemand distingue très nettement deux registres : la lettre \cle{formelle} (\all{der formelle Brief}) et la lettre \cle{informelle} (\all{der informelle Brief}). C'est l'objet de cette section de grammaire.

\courssub{Observation : deux lettres, deux registres}

Reprenons les deux lettres de notre texte support et observons leurs débuts et leurs fins.

\begin{exemplebox}[Observer avant de comprendre]
\textbf{Lettre formelle (candidature) :}\par
\all{Sehr geehrte Damen und Herren, mit großem Interesse habe ich Ihre Anzeige gelesen. ... Mit freundlichen Grüßen}\par
« Mesdames, Messieurs, j'ai lu votre annonce avec un grand intérêt. … Je vous prie d'agréer mes salutations distinguées. »\par
\medskip
\textbf{Lettre informelle (encouragement) :}\par
\all{Liebe Amina, danke für deinen Brief! Ich freue mich sehr für dich. ... Viele Grüße, deine Sarah}\par
« Chère Amina, merci pour ta lettre ! Je suis très contente pour toi. … Amitiés, ta Sarah »
\end{exemplebox}

Observons les différences :

\begin{itemize}
\item \textbf{L'adresse initiale} (\all{die Anrede}) : \all{Sehr geehrte Damen und Herren} contre \all{Liebe Amina}.
\item \textbf{Le pronom personnel} : \all{Ihre} / \all{Sie} (vouvoiement) contre \all{deinen} / \all{dich} (tutoiement).
\item \textbf{La formule finale} (\all{die Grußformel}) : \all{Mit freundlichen Grüßen} contre \all{Viele Grüße}.
\item \textbf{La signature} (\all{die Unterschrift}) : nom complet dans la lettre formelle, simple prénom dans la lettre informelle.
\end{itemize}

\begin{methode}[Choisir le bon registre]
Avant d'écrire une lettre en allemand, pose-toi une seule question : \emph{à qui j'écris ?}
\begin{enumerate}
\item Destinataire \textbf{inconnu}, \textbf{supérieur hiérarchique}, \textbf{administration}, \textbf{entreprise} $\rightarrow$ registre \cle{formel} : vouvoiement \all{Sie}.
\item Destinataire \textbf{ami}, \textbf{camarade de classe}, \textbf{membre de la famille} $\rightarrow$ registre \cle{informel} : tutoiement \all{du}.
\item En cas de doute (un professeur, un correspondant adulte), choisis toujours le registre \textbf{formel} : on ne reproche jamais à quelqu'un d'être trop poli.
\end{enumerate}
\end{methode}

\courssub{La structure commune à toute lettre allemande}

Formelle ou informelle, toute lettre allemande suit la même architecture en cinq parties :

\begin{enumerate}
\item \textbf{L'en-tête} : adresse de l'expéditeur (lettre formelle), puis le \cle{lieu et la date} à droite : \all{Yaoundé, den 15. März 2025}.
\item \textbf{L'adresse} : \all{die Anrede}, suivie d'une \emph{virgule}.
\item \textbf{Le corps de la lettre} : \all{der Brieftext}, qui commence par une \emph{minuscule}.
\item \textbf{La formule de politesse finale} : \all{die Grußformel}.
\item \textbf{La signature} : \all{die Unterschrift}.
\end{enumerate}

Dans une lettre formelle (candidature, courrier administratif), on ajoute souvent avant l'Anrede une ligne d'\textbf{objet}, \all{der Betreff}, qui annonce le sujet : \all{Betreff: Bewerbung um die Stelle als Verkäuferin} « Objet : candidature au poste de vendeuse ».

\begin{propriete}[La date : « Yaoundé, den 15. März 2025 »]
La date s'écrit à droite, sous la forme : \textbf{ville + virgule + den + jour (ordinal) + mois + année}.\par
\all{Yaoundé, den 15. März 2025} = « Yaoundé, le 15 mars 2025 ».\par
Le mot \all{den} est l'\textbf{accusatif} de l'article \all{der} : la date est considérée comme un complément de temps (réponse implicite à la question « quand ? »). Le jour est un adjectif ordinal suivi d'un point : \all{den 1. Mai}, \all{den 23. Dezember}.\par
Variante sans article, possible dans les en-têtes : \all{Yaoundé, 15.03.2025}.
\end{propriete}

\begin{exemplebox}[Des dates à lire et à écrire]
\all{Douala, den 3. Februar 2025} « Douala, le 3 février 2025 »\par
\all{Bamenda, den 20. Juni 2025} « Bamenda, le 20 juin 2025 »\par
\all{Berlin, den 1. September 2025} « Berlin, le 1er septembre 2025 »
\end{exemplebox}

\begin{propriete}[La virgule après l'Anrede et la minuscule au début du corps]
En allemand, l'adresse (\all{die Anrede}) est suivie d'une \textbf{virgule}, et la première ligne du corps de la lettre commence par une \textbf{minuscule} (sauf si c'est un nom, qui garde toujours sa majuscule).\par
\all{Sehr geehrte Damen und Herren,}\par
\all{mit großem Interesse habe ich Ihre Anzeige gelesen.}\par
« Mesdames, Messieurs, j'ai lu votre annonce avec un grand intérêt. »\par
C'est l'inverse du français, où la première ligne commence par une majuscule après « Monsieur, ». Attention donc à ce réflexe francophone !
\end{propriete}

\begin{attention}[Minuscule après la virgule : l'erreur n° 1 des francophones]
En français on écrit : « Madame, \textbf{J}e vous écris… ». En allemand, c'est une faute : on écrit \all{Sehr geehrte Frau Müller, \textbf{i}ch schreibe Ihnen…}. Seuls les noms propres et les noms communs gardent leur majuscule : \all{Lieber Paul, \textbf{d}eine Schwester hat mir geschrieben.} mais \all{Lieber Paul, \textbf{D}eutschland gefällt mir sehr.}
\end{attention}

\courssub{La lettre formelle : vouvoiement et formules consacrées}

\begin{definition}[L'Anrede formelle]
Pour ouvrir une lettre formelle, on utilise :\par
\all{Sehr geehrte Damen und Herren,} « Mesdames, Messieurs, » (destinataire inconnu) ;\par
\all{Sehr geehrte Frau Müller,} « Madame Müller, » (femme nommée) ;\par
\all{Sehr geehrter Herr Müller,} « Monsieur Müller, » (homme nommé) ;\par
\all{Sehr geehrter Herr Direktor,} « Monsieur le Directeur, » (fonction).
\end{definition}

\begin{propriete}[L'accord de « geehrte/r »]
L'adjectif \all{geehrte} s'accorde avec le destinataire :\par
devant un nom \textbf{féminin} (\all{die Frau}) : \all{Sehr geehrte Frau Müller} ;\par
devant un nom \textbf{masculin} (\all{der Herr}) : \all{Sehr geehrter Herr Müller} ;\par
au \textbf{pluriel} : \all{Sehr geehrte Damen und Herren}.\par
Remarque : \all{der Herr} devient \all{Herren} au pluriel et prend aussi la forme \all{Herrn} à l'accusatif et au datif (\all{ich danke Herrn Müller} « je remercie M. Müller »).
\end{propriete}

\begin{attention}[Ne pas écrire « Sehr geehrter Frau Müller »]
Erreur très fréquente chez les francophones : oublier l'accord. Retiens le couple :\par
\all{Sehr geehrte Frau Müller} (féminin, en \textbf{-e}) — \all{Sehr geehrter Herr Müller} (masculin, en \textbf{-r}).\par
Astuce : l'article te donne la réponse — \all{die Frau} $\rightarrow$ \all{geehrte} ; \all{der Herr} $\rightarrow$ \all{geehrter}.
\end{attention}

\begin{attention}[« Sie » prend toujours une majuscule]
Le pronom de politesse \all{Sie} (vous) et ses formes (\all{Ihnen}, \all{Ihre}) prennent \textbf{toujours une majuscule}, même au milieu de la phrase : \all{Ich danke Ihnen für Ihre Antwort.} « Je vous remercie de votre réponse. »\par
En revanche, \all{sie} minuscule signifie « elle » ou « elles » : \all{Sie schreibt Frau Müller, weil sie krank ist.} « Elle écrit à Mme Müller parce qu'elle est malade. » La majuscule change le sens !
\end{attention}

\begin{propriete}[La Grußformel formelle]
La lettre formelle se termine par :\par
\all{Mit freundlichen Grüßen} « Je vous prie d'agréer mes salutations distinguées » (souvent abrégé \all{MfG} dans les courriels) ;\par
variante plus chaleureuse : \all{Mit freundlichen Grüßen und bestem Dank} « Avec mes salutations distinguées et mes meilleurs remerciements ».\par
Contrairement au français, il n'y a \textbf{pas de virgule} après la formule finale en allemand : on écrit \all{Mit freundlichen Grüßen} puis, à la ligne, la signature.
\end{propriete}

\courssub{La lettre informelle : tutoiement et chaleur}

\begin{definition}[L'Anrede informelle]
Pour ouvrir une lettre à un proche :\par
\all{Liebe Amina,} « Chère Amina, » (femme) ; \all{Lieber Paul,} « Cher Paul, » (homme) ;\par
\all{Liebe Familie Ngono,} « Chère famille Ngono, » ; \all{Hallo Sarah,} « Salut Sarah, » (très amical).\par
Même accord : \all{liebe} devant un prénom féminin, \all{lieber} devant un prénom masculin.
\end{definition}

\begin{propriete}[Le tutoiement « du »]
Dans la lettre informelle, on tutoie : \all{du}, \all{dich}, \all{dir}, \all{dein}. Dans les lettres (et courriels) personnels, l'usage traditionnel recommande d'écrire \all{Du}, \all{Dich}, \all{Dein} avec une \textbf{majuscule} par politesse affective — c'est facultatif depuis la réforme orthographique, mais très apprécié : \all{Ich freue mich sehr für Dich.} « Je suis très content pour toi. »
\end{propriete}

\begin{propriete}[Les Grußformeln informelles]
\all{Viele Grüße} « Amitiés » ; \all{Liebe Grüße} « Bien affectueusement » ; \all{Viele liebe Grüße} « Toutes mes amitiés » ; \all{Bis bald!} « À bientôt ! » ; \all{Deine Sarah} / \all{Dein Paul} « Ta Sarah / Ton Paul ».
\end{propriete}

\courssub{Tableau récapitulatif et comparaison avec le français}

\begin{center}
\begin{tabularx}{\linewidth}{|X|X|X|}
\hline
\rowcolor{popblueL} Élément & Lettre formelle & Lettre informelle \\
\hline
Anrede & Sehr geehrte Damen und Herren, / Sehr geehrte Frau Müller, / Sehr geehrter Herr Müller, & Liebe Amina, / Lieber Paul, / Hallo Sarah, \\
\hline
Pronom & Sie, Ihnen, Ihre (majuscule) & du, dich, dir, dein (Du recommandé) \\
\hline
Grußformel & Mit freundlichen Grüßen (MfG) & Viele Grüße / Liebe Grüße / Bis bald! \\
\hline
Signature & Prénom + NOM & Prénom seul (Deine Sarah) \\
\hline
\end{tabularx}
\end{center}

\begin{center}
\begin{tabularx}{\linewidth}{|X|X|X|}
\hline
\rowcolor{popblueL} Français & Deutsch & Remarque \\
\hline
Monsieur, Madame & Sehr geehrte Damen und Herren, & virgule en allemand \\
\hline
Madame Dupont, & Sehr geehrte Frau Dupont, & accord en -e \\
\hline
Monsieur Dupont, & Sehr geehrter Herr Dupont, & accord en -r \\
\hline
Je vous prie d'agréer… / Cordialement & Mit freundlichen Grüßen & sans virgule après \\
\hline
Amitiés / Bises & Viele Grüße / Liebe Grüße & registre informel \\
\hline
\end{tabularx}
\end{center}

\begin{popfigure}
\begin{tikzpicture}[
  box/.style={draw=popline, rounded corners=2pt, align=center, inner sep=5pt, font=\small},
  head/.style={box, font=\small\bfseries, text=white},
  lab/.style={font=\small\itshape, text=popmuted}
]
\node[head, fill=popblue, minimum width=5.6cm] (hf) at (0,0) {formeller Brief (Sie)};
\node[head, fill=poporange, minimum width=5.6cm] (hi) at (8.2,0) {informeller Brief (du)};
\node[box, fill=popblueL, minimum width=5.6cm, align=center] (a1) at (0,-1.5){Anrede:\\ Sehr geehrte Damen und Herren,\\ Sehr geehrte Frau Müller,};
\node[box, fill=poporangeL, minimum width=5.6cm, align=center] (a2) at (8.2,-1.5){Anrede:\\ Liebe Amina,\\ Lieber Paul,};
\node[box, fill=popblueL, minimum width=5.6cm, align=center] (p1) at (0,-3.2){Pronom:\\ Sie / Ihnen / Ihre\\ (majuscule !)};
\node[box, fill=poporangeL, minimum width=5.6cm, align=center] (p2) at (8.2,-3.2){Pronom:\\ du / dich / dein\\ (Du recommandé)};
\node[box, fill=popblueL, minimum width=5.6cm, align=center] (g1) at (0,-4.9){Grußformel:\\ Mit freundlichen Grüßen};
\node[box, fill=poporangeL, minimum width=5.6cm, align=center] (g2) at (8.2,-4.9){Grußformel:\\ Viele Grüße / Bis bald!};
\draw[-{Stealth}, popblue, thick] (hf) -- (a1);
\draw[-{Stealth}, popblue, thick] (a1) -- (p1);
\draw[-{Stealth}, popblue, thick] (p1) -- (g1);
\draw[-{Stealth}, poporange, thick] (hi) -- (a2);
\draw[-{Stealth}, poporange, thick] (a2) -- (p2);
\draw[-{Stealth}, poporange, thick] (p2) -- (g2);
\node[lab] at (4.1,-1.5) {ouverture};
\node[lab] at (4.1,-3.2) {pronom};
\node[lab] at (4.1,-4.9) {clôture};
\end{tikzpicture}
\legende{Les deux registres de la lettre allemande : formel (Sie) et informel (du).}
\end{popfigure}

\courssub{Exemples traduits}

\begin{exemplebox}[Phrases de lettres, registre par registre]
\all{Sehr geehrter Herr Direktor, ich danke Ihnen für Ihre Antwort.} « Monsieur le Directeur, je vous remercie de votre réponse. »\par
\all{Sehr geehrte Damen und Herren, ich bewerbe mich um die Stelle als Verkäuferin.} « Mesdames, Messieurs, je pose ma candidature au poste de vendeuse. »\par
\all{Ich freue mich auf Ihre Antwort. Mit freundlichen Grüßen} « Dans l'attente de votre réponse, je vous prie d'agréer mes salutations distinguées. »\par
\medskip
\all{Liebe Amina, danke für deinen Brief!} « Chère Amina, merci pour ta lettre ! »\par
\all{Lieber Paul, ich gratuliere dir zu deinem Erfolg!} « Cher Paul, je te félicite pour ta réussite ! »\par
\all{Viele Grüße und bis bald, deine Sarah.} « Amitiés et à bientôt, ta Sarah. »
\end{exemplebox}

\courssub{Exercice résolu}

\begin{exoresolu}[Formel ou informel ? Corrige et complète]
\textbf{Énoncé.} Voici le début d'une lettre de candidature écrite par un élève de Première. Il contient quatre erreurs de registre ou d'orthographe. Trouve-les et corrige-les :\par
\all{Sehr geehrter Frau Kamga, Ich habe ihre Anzeige gelesen und ich möchte mich bewerben. Viele Grüße, Junior Etoundi}
\tcblower
\textbf{Solution détaillée.}\par
1. \all{Sehr geehrter Frau Kamga} $\rightarrow$ \all{Sehr geehrte Frau Kamga} : \all{die Frau} est féminin, l'adjectif prend la terminaison \textbf{-e}.\par
2. \all{Ich habe…} avec majuscule après la virgule $\rightarrow$ \all{ich habe…} : après l'Anrede suivie d'une virgule, le corps commence par une \textbf{minuscule}.\par
3. \all{ihre Anzeige} $\rightarrow$ \all{Ihre Anzeige} : lettre formelle, donc vouvoiement ; \all{Ihre} prend la \textbf{majuscule}.\par
4. \all{Viele Grüße} $\rightarrow$ \all{Mit freundlichen Grüßen} : formule \textbf{informelle} interdite dans une candidature.\par
\textbf{Version corrigée :} \all{Sehr geehrte Frau Kamga, ich habe Ihre Anzeige gelesen und ich möchte mich bewerben. Mit freundlichen Grüßen, Junior Etoundi}
\end{exoresolu}

\begin{exercice}[À toi d'écrire les ouvertures et clôtures]
Rédige l'Anrede et la Grußformel adaptées : 1) une lettre au directeur d'une entreprise de Douala ; 2) une lettre à ta correspondante autrichienne Lisa ; 3) une lettre à une entreprise dont tu ne connais pas le destinataire ; 4) une lettre à ton oncle Pierre.
\end{exercice}

\begin{corrige}[Exercice : ouvertures et clôtures]
1. \all{Sehr geehrter Herr Direktor, … Mit freundlichen Grüßen} ; 2. \all{Liebe Lisa, … Viele Grüße} (ou \all{Liebe Grüße}) ; 3. \all{Sehr geehrte Damen und Herren, … Mit freundlichen Grüßen} ; 4. \all{Lieber Onkel Pierre, … Viele liebe Grüße}.
\end{corrige}

\begin{aretenir}[L'essentiel de la lettre allemande]
\begin{itemize}
\item Deux registres : \textbf{formel} (\all{Sie} + majuscule) et \textbf{informel} (\all{du}).
\item Anrede formelle : \all{Sehr geehrte Frau…} / \all{Sehr geehrter Herr…} / \all{Sehr geehrte Damen und Herren,} ; informelle : \all{Liebe…} / \all{Lieber…}.
\item Après la virgule de l'Anrede : \textbf{minuscule} au début du corps.
\item Date : \all{Yaoundé, den 15. März 2025} (accusatif avec \all{den}).
\item Clôture formelle : \all{Mit freundlichen Grüßen} ; informelle : \all{Viele Grüße}, \all{Bis bald!}.
\end{itemize}
\end{aretenir}

Maintenant que la structure et le registre de la lettre sont acquis, il reste à soigner le ton : comment formuler une demande polie dans le corps de la lettre ? C'est l'objet de la prochaine section de grammaire : la politesse avec \all{möchte} et \all{würde}.

\coursec{Grammaire 2 : la politesse avec « möchte » et « würde »}

Dans la section précédente, nous avons appris à construire une lettre formelle et une lettre informelle : l'en-tête, l'appel (\all{Sehr geehrte Damen und Herren}), la formule de politesse finale (\all{Mit freundlichen Grüßen}). Mais une lettre de motivation ne se juge pas seulement à sa structure : elle se juge aussi à son \cle{ton}. Or le ton, en allemand, passe par deux petites formes magiques : \all{möchte} et \all{würde}. C'est elles que nous étudions maintenant.

\courssub{Observation : pourquoi pas « Ich will » ?}

Lisons ces trois phrases, tirées de lettres de candidature :

\begin{exemplebox}[Trois façons de dire « je veux »]
\all{Ich will mich um die Stelle bewerben.} « Je veux postuler au poste. » (trop direct !)\par
\all{Ich möchte mich um die Stelle bewerben.} « Je voudrais postuler au poste. » (poli)\par
\all{Ich würde mich sehr freuen, von Ihnen zu hören.} « Je serais très heureux d'avoir de vos nouvelles. » (très poli)
\end{exemplebox}

Les trois phrases expriment la même volonté, mais elles ne produisent pas le même effet sur le lecteur. \all{Ich will} sonne comme un ordre ou une exigence d'enfant ; \all{Ich möchte} ressemble à notre « je voudrais » ; \all{Ich würde mich freuen} correspond à « j'aimerais / cela me ferait plaisir ». En français, nous passons du présent (« je veux ») au conditionnel (« je voudrais ») pour adoucir. L'allemand fait exactement la même chose : il passe de l'indicatif (\all{will}) au \cle{subjonctif II} (\all{möchte}, \all{würde}).

\begin{definition}[Le subjonctif II de politesse]
Le \cle{Konjunktiv II} (subjonctif II) est le mode de l'irréel, du souhait et de la politesse en allemand. Deux formes sont indispensables dès la classe de Première : \all{möchte} (subjonctif II de \all{mögen}, « aimer ») et \all{würde} (subjonctif II de \all{werden}). Elles permettent d'exprimer un souhait ou une demande sans paraître impoli.
\end{definition}

\courssub{« möchte » : le souhait poli}

\textbf{Formation.}\par
\all{möchte} est le subjonctif II du verbe modal \all{mögen} (aimer). On le forme à partir du prétérit \all{mochte} auquel on ajoute un tréma : \all{mochte} devient \all{möchte}. Ne cherche pas à mémoriser ce mécanisme pour l'instant : retiens simplement que \all{möchte} est une forme figée qui signifie « voudrais ». Attention : \all{möchte} est déjà une forme de politesse ; on ne dit donc jamais \all{ich würde mögen}.

\begin{center}
\begin{tabular}{|l|l|l|}
\hline
\rowcolor{popblueL} Personne & Forme & Exemple traduit \\
\hline
ich & möchte & \all{Ich möchte mich bewerben.} Je voudrais postuler. \\
\hline
du & möchtest & \all{Möchtest du einen Kaffee?} Voudrais-tu un café ? \\
\hline
er/sie/es & möchte & \all{Sie möchte einen Termin.} Elle souhaite un rendez-vous. \\
\hline
wir & möchten & \all{Wir möchten Sie einladen.} Nous voudrions vous inviter. \\
\hline
ihr & möchtet & \all{Möchtet ihr mitkommen?} Voudriez-vous venir avec nous ? \\
\hline
sie/Sie & möchten & \all{Möchten Sie Platz nehmen?} Voudriez-vous vous asseoir ? \\
\hline
\end{tabular}
\end{center}

\begin{attention}[Attention à la 3e personne]
À la 3e personne du singulier, la forme est \all{er möchte} SANS -t, contrairement au présent de l'indicatif (\all{er mag}). De même : \all{er würde}, jamais \all{er würdet}. C'est une particularité du subjonctif II : ich et er/sie/es ont toujours la même forme.
\end{attention}

\begin{propriete}[« möchte » + infinitif]
\all{möchte} exprime un souhait poli et se comporte comme un verbe modal : il occupe la \cle{position 2} dans la phrase déclarative et renvoie l'\cle{infinitif en fin de phrase}.\par
\all{Ich möchte mich um die Stelle bewerben.} « Je voudrais postuler au poste. »\par
\all{Ich möchte Sie zu einem Gespräch einladen.} « Je voudrais vous inviter à un entretien. »\par
On renforce souvent la courtoisie avec l'adverbe \all{gern} : \all{Ich möchte gern in Ihrer Firma arbeiten.} « Je voudrais volontiers travailler dans votre entreprise. »
\end{propriete}

\begin{exemplebox}[« möchte » dans la vie quotidienne et dans la lettre]
\all{Möchten Sie einen Termin?} « Souhaitez-vous un rendez-vous ? »\par
\all{Ich möchte mich für Ihre Einladung bedanken.} « Je voudrais vous remercier de votre invitation. »\par
\all{Wir möchten Sie bitten, uns bald zu antworten.} « Nous vous prions de bien vouloir nous répondre rapidement. »\par
\all{Möchtest du am Wochenende mit uns ins Stadion kommen?} « Voudrais-tu venir au stade avec nous ce week-end ? » (registre informel, avec un ami)
\end{exemplebox}

\courssub{« würde » + infinitif : le conditionnel allemand}

\textbf{Formation.}\par
\all{würde} est le subjonctif II du verbe \all{werden} (devenir). Combiné avec un infinitif placé en fin de phrase, il forme l'équivalent de notre \cle{conditionnel français} : « je ferais, j'aimerais, je serais ».

\begin{center}
\begin{tabular}{|l|l|l|}
\hline
\rowcolor{popgreenL} Personne & Forme & Exemple traduit \\
\hline
ich & würde & \all{Ich würde mich freuen.} Cela me ferait plaisir. \\
\hline
du & würdest & \all{Du würdest gut passen.} Tu conviendrais bien. \\
\hline
er/sie/es & würde & \all{Er würde gern kommen.} Il aimerait venir. \\
\hline
wir & würden & \all{Wir würden uns freuen.} Nous serions ravis. \\
\hline
ihr & würdet & \all{Ihr würdet es schaffen.} Vous y arriveriez. \\
\hline
sie/Sie & würden & \all{Würden Sie bitte warten?} Pourriez-vous attendre ? \\
\hline
\end{tabular}
\end{center}

\begin{propriete}[Emplois de « würde » + infinitif]
\begin{enumerate}
\item \cle{Hypothèse polie} dans une lettre : \all{Ich würde mich sehr freuen, von Ihnen zu hören.} « Je serais très heureux d'avoir de vos nouvelles. »
\item \cle{Demande adoucie} : \all{Würden Sie bitte das Formular ausfüllen?} « Pourriez-vous remplir le formulaire, s'il vous plaît ? »
\item \cle{Souhait irréel ou rêve} : \all{Ich würde gern in Deutschland studieren.} « J'aimerais étudier en Allemagne. »
\end{enumerate}
Dans tous les cas, \all{würde} est en position 2 et l'infinitif va en fin de phrase.
\end{propriete}

\begin{exemplebox}[La phrase modèle de la lettre de motivation]
\all{Ich würde mich über eine positive Antwort sehr freuen.} « Je me réjouirais beaucoup d'une réponse positive. »\par
Analyse de l'ordre des mots : \all{Ich} (sujet, position 1) — \all{würde} (verbe conjugué, position 2) — \all{mich über eine positive Antwort sehr} (milieu de phrase) — \all{freuen} (infinitif, dernière place). Cette phrase est la formule de conclusion la plus fréquente des lettres formelles allemandes : apprends-la par cœur.
\end{exemplebox}

\courssub{Tableau de conjugaison complet : möchte et würde}

\begin{popfigure}
\begin{tikzpicture}[
  personne/.style={rectangle, rounded corners=3pt, draw=popblue, fill=popblueL, font=\small, minimum width=1.5cm, minimum height=0.8cm},
  forme/.style={rectangle, rounded corners=3pt, draw=poporange, fill=poporangeL, font=\small\bfseries, minimum width=1.7cm, minimum height=0.8cm},
  forme2/.style={rectangle, rounded corners=3pt, draw=popgreen, fill=popgreenL, font=\small\bfseries, minimum width=1.7cm, minimum height=0.8cm},
  ex/.style={font=\footnotesize, align=left, anchor=west}
]
\node[personne] at (0,5) {ich};
\node[forme] at (2.2,5) {möchte};
\node[forme2] at (4.4,5) {würde};
\node[ex] at (5.8,5) {\all{Ich möchte/würde gern kommen.} « Je voudrais/j'aimerais venir. »};

\node[personne] at (0,4) {du};
\node[forme] at (2.2,4) {möchtest};
\node[forme2] at (4.4,4) {würdest};
\node[ex] at (5.8,4) {\all{Möchtest du Tee?} « Voudrais-tu du thé ? »};

\node[personne] at (0,3) {er/sie/es};
\node[forme] at (2.2,3) {möchte};
\node[forme2] at (4.4,3) {würde};
\node[ex] at (5.8,3) {\all{Sie würde gern arbeiten.} « Elle aimerait travailler. »};

\node[personne] at (0,2) {wir};
\node[forme] at (2.2,2) {möchten};
\node[forme2] at (4.4,2) {würden};
\node[ex] at (5.8,2) {\all{Wir möchten uns bedanken.} « Nous voudrions remercier. »};

\node[personne] at (0,1) {ihr};
\node[forme] at (2.2,1) {möchtet};
\node[forme2] at (4.4,1) {würdet};
\node[ex] at (5.8,1) {\all{Würdet ihr uns helfen?} « Pourriez-vous nous aider ? »};

\node[personne] at (0,0) {sie/Sie};
\node[forme] at (2.2,0) {möchten};
\node[forme2] at (4.4,0) {würden};
\node[ex] at (5.8,0) {\all{Möchten Sie einen Termin?} « Souhaitez-vous un rendez-vous ? »};
\end{tikzpicture}
\legende{Conjugaison de \all{möchte} (souhait poli) et \all{würde} (conditionnel) : remarque que ich et er/sie/es ont toujours la même forme.}
\end{popfigure}

\courssub{Comparaison français / allemand}

\begin{center}
\begin{tabularx}{\linewidth}{|X|X|X|}
\hline
\rowcolor{popblueL} Français & Deutsch & Remarque \\
\hline
je veux (trop direct) & \all{ich will} & à éviter dans une lettre formelle \\
\hline
je voudrais & \all{ich möchte} + infinitif & souhait poli, infinitif en fin de phrase \\
\hline
j'aimerais / cela me ferait plaisir & \all{ich würde mich freuen} & conditionnel, formule de lettre \\
\hline
pourriez-vous... ? & \all{würden Sie / könnten Sie...?} & demande très polie \\
\hline
je serais heureux de & \all{ich würde mich freuen, wenn...} & suivi d'une subordonnée avec \all{wenn} \\
\hline
\end{tabularx}
\end{center}

\courssub{Exceptions et limites : wäre, hätte, könnte}

\begin{propriete}[Les formes propres du subjonctif II]
En théorie, on peut construire le subjonctif II de tous les verbes avec \all{würde} + infinitif. Mais pour les verbes très fréquents — \all{sein}, \all{haben} et les verbes modaux — l'allemand préfère les \cle{formes propres} du subjonctif II :
\begin{itemize}
\item \all{sein} : \all{ich wäre} (je serais) et non \all{ich würde sein} ;
\item \all{haben} : \all{ich hätte} (j'aurais) et non \all{ich würde haben} ;
\item \all{können} : \all{ich könnte} (je pourrais) et non \all{ich würde können} ;
\item \all{es gibt} : \all{es gäbe} (il y aurait).
\end{itemize}
Exemples : \all{Ich hätte gern mehr Zeit.} « J'aimerais avoir plus de temps. » ; \all{Wäre das möglich?} « Serait-ce possible ? » ; \all{Könnten Sie mir bitte helfen?} « Pourriez-vous m'aider, s'il vous plaît ? » Ces formes sont très appréciées au Bac : elles montrent une maîtrise fine de la langue.
\end{propriete}

\begin{attention}[Les pièges des francophones]
\begin{enumerate}
\item \cle{Jamais « Ich will » dans une lettre de motivation.} \all{Ich will diese Stelle.} sonne arrogant, voire impoli. Écris toujours : \all{Ich möchte mich um diese Stelle bewerben.}
\item \cle{« würde » ne se traduit pas toujours par « voudrais ».} \all{Ich würde mich freuen} signifie « cela me ferait plaisir / je serais heureux », et non « je voudrais me réjouir ». \all{würde} + infinitif = conditionnel français, pas forcément « vouloir ».
\item \cle{N'oublie pas l'infinitif en fin de phrase.} Faux : \all{Ich möchte bewerben mich.} Juste : \all{Ich möchte mich bewerben.}
\item \cle{Pas de -t à la 3e personne.} Faux : \all{er möchtet}, \all{er würdet}. Juste : \all{er möchte}, \all{er würde}.
\item \cle{Le tréma est obligatoire.} \all{mochte} (prétérit : « j'aimais ») et \all{möchte} (« je voudrais ») sont deux formes différentes : sans tréma, le sens change.
\end{enumerate}
\end{attention}

\begin{savaistu}[La formule de politesse reine]
\all{Könnten Sie mir bitte...?} (subjonctif II de \all{können}) est la formule de demande la plus polie de la langue allemande : « Pourriez-vous, s'il vous plaît... ? ». Dans les pays germanophones, on l'entend partout : au guichet de la gare, à l'hôtel, au bureau. Un Allemand qui entend \all{Geben Sie mir das!} (« Donnez-moi ça ! ») fronce les sourcils ; avec \all{Könnten Sie mir das bitte geben?}, tout devient possible. Au Bac, glisser un \all{könnte}, un \all{wäre} ou un \all{hätte} dans ta production écrite est un signe d'excellence qui impressionne toujours le correcteur.
\end{savaistu}

\begin{exoresolu}[Transforme les phrases trop directes en phrases polies]
Énoncé : réécris chaque phrase avec \all{möchte} ou \all{würde} pour la rendre acceptable dans une lettre formelle.\par
1. \all{Ich will die Stelle haben.}\par
2. \all{Gib mir eine Antwort!}\par
3. \all{Ich will in Ihrer Firma arbeiten.}
\tcblower
Solution détaillée :\par
1. Le verbe \all{wollen} est trop direct. On le remplace par \all{möchte} + infinitif en fin de phrase : \all{Ich möchte mich um die Stelle bewerben.} « Je voudrais postuler au poste. » (On évite de dire qu'on « veut avoir » le poste : on dit qu'on souhaite postuler.)\par
2. L'impératif est brutal. On utilise \all{würde} + \all{sich freuen} : \all{Ich würde mich über eine Antwort sehr freuen.} « Je me réjouirais beaucoup d'une réponse. » Noter l'ordre : \all{würde} en position 2, \all{freuen} en dernière position.\par
3. \all{Ich möchte gern in Ihrer Firma arbeiten.} « Je voudrais volontiers travailler dans votre entreprise. » L'adverbe \all{gern} renforce encore la courtoisie.
\end{exoresolu}

\begin{exercice}[À ton tour]
1. Conjugue \all{möchte} et \all{würde} aux six personnes, de mémoire.\par
2. Traduis en allemand : a) « Je voudrais vous poser une question. » b) « Nous serions très heureux de vous rencontrer. » c) « Souhaitez-vous un rendez-vous lundi ? »\par
3. Rends polies ces phrases : \all{Ich will einen Termin.} / \all{Komm zu mir!}
\end{exercice}

\begin{corrige}[Exercice « À ton tour »]
1. Voir les tableaux de conjugaison ci-dessus.\par
2. a) \all{Ich möchte Sie etwas fragen.} ou \all{Ich möchte Ihnen eine Frage stellen.} b) \all{Wir würden uns sehr freuen, Sie kennenzulernen.} c) \all{Möchten Sie einen Termin am Montag?}\par
3. \all{Ich möchte gern einen Termin.} ; \all{Würden Sie bitte zu mir kommen?} ou \all{Könnten Sie bitte zu mir kommen?}
\end{corrige}

\begin{aretenir}[L'essentiel de la section]
\begin{itemize}
\item \all{möchte} = subjonctif II de \all{mögen} : « je voudrais » ; conjugaison : \all{ich möchte, du möchtest, er möchte, wir möchten, ihr möchtet, sie/Sie möchten}.
\item \all{würde} + infinitif = conditionnel : « je ferais, j'aimerais » ; conjugaison : \all{ich würde, du würdest, er würde, wir würden, ihr würdet, sie/Sie würden}.
\item Le verbe conjugué est en position 2, l'infinitif en fin de phrase : \all{Ich würde mich über eine Antwort sehr freuen.}
\item Jamais \all{ich will} dans une lettre formelle : toujours \all{ich möchte}.
\item Avec \all{sein}, \all{haben}, \all{können}, on préfère les formes propres : \all{wäre, hätte, könnte}.
\end{itemize}
\end{aretenir}

\coursec{Méthode du commentaire et commentaire modèle}

Dans la section précédente, nous avons appris à exprimer la politesse avec \all{möchte} et \all{würde}. Ces formes, nous les avons rencontrées dans les lettres de motivation et d'encouragement de Merveille et d'Amina. Il est maintenant temps de passer à une compétence essentielle pour l'examen et pour la vie scolaire : \cle{commenter un texte}. Commenter, c'est présenter, résumer, analyser et juger un texte de manière organisée. Cette section vous donne une méthode en cinq étapes, des phrases toutes prêtes (\all{Redemittel}), un commentaire modèle complet, puis des entraînements à l'écrit et à l'oral.

\courssub{Les cinq étapes du commentaire de texte}

\courssub{Observer avant de rédiger}

Lisez d'abord ce mini-commentaire rédigé par un élève de Première sur la lettre de motivation de Merveille :

\begin{exemplebox}[Un premier commentaire à observer]
\all{Der Text ist ein Bewerbungsschreiben. Merveille, eine Schülerin aus Yaoundé, schreibt an eine Organisation in Bonn. Sie möchte ein Praktikum machen. Der Brief hat eine Anrede, einen Hauptteil und eine Grußformel. Die Autorin benutzt höfliche Formen wie „Ich möchte" und „Ich würde mich freuen". Ich finde den Brief überzeugend, weil er klar und höflich ist.}

« Le texte est une lettre de motivation. Merveille, une élève de Yaoundé, écrit à une organisation à Bonn. Elle voudrait faire un stage. La lettre a une formule d'appel, une partie principale et une formule de politesse finale. L'auteure utilise des formes polies comme „Ich möchte" et „Ich würde mich freuen". Je trouve la lettre convaincante parce qu'elle est claire et polie. »
\end{exemplebox}

Observez l'ordre des idées : l'élève ne commence pas par son avis, il commence par \emph{présenter} le texte. Il ne recopie pas la lettre, il la \emph{résume}. Il ne dit pas seulement « j'aime », il \emph{justifie}. C'est exactement la démarche que nous allons formaliser.

\begin{propriete}[Les cinq étapes du commentaire]
Un commentaire de texte suit toujours le même ordre :
\begin{enumerate}
\item \cle{Présenter} : dire le type de texte, l'auteur, le destinataire, le thème. \all{Der Text ist ein Bewerbungsschreiben. Die Autorin ist... Sie schreibt an...}
\item \cle{Résumer} : donner l'essentiel en 3 ou 4 phrases, avec ses propres mots. \all{Sie möchte ein Praktikum machen und beschreibt ihre Motivation.}
\item \cle{Analyser} : décrire la structure (Anrede, Hauptteil, Grußformel) et les formules remarquables (politesse, temps, connecteurs).
\item \cle{Dégager l'intention} : dire ce que l'auteur veut obtenir. \all{Die Autorin möchte die Stelle bekommen.}
\item \cle{Donner un avis personnel justifié} : \all{Ich finde den Brief überzeugend, weil...}
\end{enumerate}
\end{propriete}

\begin{popfigure}
\begin{tikzpicture}[
  node distance=0.55cm,
  etape/.style={rectangle, rounded corners=3pt, draw=popblue, fill=popblueL, text width=4.6cm, align=center, minimum height=0.9cm, font=\small},
  etapeLast/.style={rectangle, rounded corners=3pt, draw=popgreen, fill=popgreenL, text width=4.6cm, align=center, minimum height=0.9cm, font=\small},
  fleche/.style={-{Stealth[length=2.5mm]}, thick, poporange}
]
\node[etape, align=center] (e1){\textbf{1. Présentation}\\ Typ, Autor, Empfänger, Thema};
\node[etape, below=of e1, align=center] (e2){\textbf{2. Résumé}\\ 3--4 Sätze, eigene Worte};
\node[etape, below=of e2, align=center] (e3){\textbf{3. Analyse}\\ Struktur und Formeln};
\node[etape, below=of e3, align=center] (e4){\textbf{4. Intention}\\ Was will der Autor?};
\node[etapeLast, below=of e4, align=center] (e5){\textbf{5. Avis personnel}\\ Meinung + Begründung mit \all{weil}};
\draw[fleche] (e1) -- (e2);
\draw[fleche] (e2) -- (e3);
\draw[fleche] (e3) -- (e4);
\draw[fleche] (e4) -- (e5);
\end{tikzpicture}
\legende{Organigramme : les cinq étapes du commentaire de texte, de la présentation à l'avis personnel.}
\end{popfigure}

\begin{attention}[L'avis personnel doit toujours être justifié]
En français comme en allemand, un avis sans justification vaut zéro point. La justification se construit avec \all{weil} + verbe conjugué \textbf{en dernière position} :
\all{Ich finde den Brief überzeugend, weil er klar und höflich ist.} « Je trouve la lettre convaincante parce qu'elle est claire et polie. »

Erreur typique du francophone : \all{*weil er ist klar}. Le verbe \all{ist} doit aller à la fin : \all{weil er klar ist}. Autre possibilité : \all{denn} + verbe en deuxième position : \all{Ich finde den Brief gut, denn er ist klar.}
\end{attention}

\courssub{Les phrases utiles (Redemittel) pour chaque étape}

Pour commenter un texte à l'écrit ou à l'oral, il faut disposer de phrases toutes prêtes. Les voici, classées par étape.

\begin{center}
\begin{tabularx}{\linewidth}{|l|X|X|}
\hline
\rowcolor{popblueL} Étape & Deutsch & Français \\
\hline
Présentation & \all{Der Text ist ein Bewerbungsschreiben / ein Ermutigungsbrief.} & Le texte est une lettre de motivation / d'encouragement. \\
\hline
Présentation & \all{Die Autorin ist eine Schülerin aus Yaoundé.} & L'auteure est une élève de Yaoundé. \\
\hline
Présentation & \all{Sie schreibt an eine Organisation in Bonn.} & Elle écrit à une organisation à Bonn. \\
\hline
Résumé & \all{Sie möchte ein Praktikum machen.} & Elle voudrait faire un stage. \\
\hline
Résumé & \all{Sie beschreibt ihre Motivation und ihre Kenntnisse.} & Elle décrit sa motivation et ses connaissances. \\
\hline
Analyse & \all{Der Brief hat drei Teile: Anrede, Hauptteil und Grußformel.} & La lettre a trois parties : l'appel, le corps et la formule finale. \\
\hline
Analyse & \all{Die Autorin benutzt höfliche Formen wie „Ich möchte" und „Ich würde mich freuen".} & L'auteure utilise des formes polies comme... \\
\hline
Intention & \all{Die Autorin möchte die Stelle bekommen.} & L'auteure voudrait obtenir le poste. \\
\hline
Avis & \all{Am Ende kann man sagen, dass der Brief überzeugend ist.} & Pour finir, on peut dire que la lettre est convaincante. \\
\hline
Avis & \all{Ich finde den Brief klar, weil er gut strukturiert ist.} & Je trouve la lettre claire parce qu'elle est bien structurée. \\
\hline
\end{tabularx}
\end{center}

\begin{attention}[Faux amis et pièges dans les Redemittel]
\begin{itemize}
\item \all{die Autorin} = l'auteure (femme) ; \all{der Autor} = l'auteur (homme). Choisissez selon la personne : Merveille et Amina sont des \all{Autorinnen}.
\item \all{bekommen} ne veut pas dire « devenir » mais « obtenir, recevoir » : \all{Sie möchte die Stelle bekommen.} « Elle voudrait obtenir le poste. » « Devenir » se dit \all{werden}.
\item \all{die Stelle} = le poste, l'emploi ; ne pas confondre avec \all{die Stelle} au sens géographique (l'endroit) — le contexte décide.
\item Dans \all{Am Ende kann man sagen, dass...}, le verbe de la subordonnée va à la fin : \all{dass der Brief überzeugend ist}.
\item \all{sich bewerben} s'emploie avec \all{bei} + datif (l'organisme) ou \all{um} + accusatif (le poste) : \all{sich bei einer Firma bewerben}, \all{sich um eine Stelle bewerben}.
\end{itemize}
\end{attention}

\courssub{Méthode : résumer un texte en trois phrases}

\begin{methode}[Comment résumer sans recopier]
Pour résumer une lettre (ou tout texte court) en trois phrases maximum :
\begin{enumerate}
\item \textbf{Qui ?} Identifier la personne qui écrit et à qui elle écrit : \all{Merveille schreibt an eine Organisation in Bonn.}
\item \textbf{Quoi et pourquoi ?} Dire ce qu'elle veut et pourquoi : \all{Sie möchte ein Praktikum machen, weil sie Deutsch lernt und die Kultur kennenlernen will.}
\item \textbf{Quelle demande ?} Terminer par la demande ou l'espoir exprimé : \all{Sie würde sich über eine positive Antwort freuen.}
\end{enumerate}
Règles d'or : utiliser ses \textbf{propres mots}, ne pas recopier de phrases entières du texte, ne pas donner son avis dans le résumé (l'avis vient à l'étape 5), employer des verbes de reformulation : \all{möchte}, \all{sich bewerben bei / um}, \all{beschreibt}, \all{bittet um}, \all{freut sich}.
\end{methode}

\begin{exemplebox}[Résumé guidé de la lettre de Merveille]
\all{Merveille, eine Schülerin aus Yaoundé, bewirbt sich bei einer Organisation in Bonn. Sie möchte ein Praktikum machen, weil sie Deutsch lernt und neue Erfahrungen sammeln will. Sie beschreibt ihre Kenntnisse und ihre Motivation. Am Ende würde sie sich über eine Einladung zu einem Gespräch freuen.}

« Merveille, une élève de Yaoundé, pose sa candidature auprès d'une organisation à Bonn. Elle voudrait faire un stage parce qu'elle apprend l'allemand et veut acquérir de nouvelles expériences. Elle décrit ses connaissances et sa motivation. À la fin, elle serait heureuse d'être invitée à un entretien. »
\end{exemplebox}

Remarquez les trois verbes clés de la reformulation : \all{sich bewerben bei} (postuler auprès de), \all{möchte} (voudrait), \all{sich freuen über} + accusatif (se réjouir de). Ces verbes permettent de résumer sans jamais recopier le texte.

\courssub{Commentaire modèle rédigé}

Voici maintenant un commentaire complet, rédigé selon les cinq étapes. Lisez-le attentivement : c'est le modèle à imiter à l'examen.

\begin{textebox}[Kommentar zum Bewerbungsbrief von Merveille]
Der Text ist eine formelle Bewerbung. Merveille, eine Schülerin aus Yaoundé, schreibt an eine Organisation in Bonn. Das Thema ist ein Praktikum in Deutschland.

Die Autorin möchte ein Praktikum machen und beschreibt ihre Motivation und ihre Kenntnisse. Sie lernt seit drei Jahren Deutsch und interessiert sich für die deutsche Kultur. Sie möchte neue Erfahrungen sammeln.

Der Brief hat eine klare Struktur: eine formelle Anrede („Sehr geehrte Damen und Herren"), einen Hauptteil mit der Motivation und eine höfliche Grußformel („Mit freundlichen Grüßen"). Die Autorin benutzt höfliche Formen wie „Ich möchte" und „Ich würde mich freuen". Diese Formen zeigen Respekt.

Die Autorin möchte die Stelle bekommen und den Leser überzeugen. Am Ende kann man sagen, dass der Brief sein Ziel erreicht.

Ich finde den Brief klar und überzeugend, weil er gut strukturiert und sehr höflich ist. Ich finde auch die Motivation der Autorin stark, weil sie konkrete Gründe nennt.
\end{textebox}

\begin{definition}[Glossaire du commentaire]
\begin{itemize}
\item \all{die Erfahrung, -en} : l'expérience
\item \all{sammeln} : rassembler, acquérir
\item \all{der Leser, -} (Pl. \all{die Leser}) : le lecteur
\item \all{überzeugen} : convaincre
\item \all{das Ziel, -e} : le but
\item \all{erreichen} : atteindre
\item \all{der Grund, die Gründe} : la raison
\item \all{nennen} : mentionner, citer
\end{itemize}
\end{definition}

\textbf{Questions de compréhension :}
\begin{enumerate}
\item \all{An wen schreibt Merveille?} (À qui Merveille écrit-elle ?)
\item \all{Welche höflichen Formen benutzt die Autorin?} (Quelles formes polies l'auteure utilise-t-elle ?)
\item \all{Warum findet der Kommentator den Brief überzeugend?} (Pourquoi le commentateur trouve-t-il la lettre convaincante ?)
\end{enumerate}

\begin{corrige}[Questions de compréhension]
\begin{enumerate}
\item \all{Sie schreibt an eine Organisation in Bonn.} « Elle écrit à une organisation à Bonn. »
\item \all{Sie benutzt „Ich möchte" und „Ich würde mich freuen".} « Elle utilise „Ich möchte" et „Ich würde mich freuen". »
\item \all{Er findet den Brief überzeugend, weil er gut strukturiert und höflich ist und weil die Autorin konkrete Gründe nennt.} « Il la trouve convaincante parce qu'elle est bien structurée, polie, et parce que l'auteure donne des raisons concrètes. »
\end{enumerate}
\end{corrige}

Analysons la construction de ce commentaire modèle : le premier paragraphe \emph{présente} (type, auteure, destinataire, thème) ; les deux paragraphes suivants \emph{résument} en trois phrases avec \all{möchte}, \all{beschreibt}, \all{interessiert sich} ; le quatrième \emph{analyse} la structure et les formules de politesse ; le cinquième dégage \emph{l'intention} ; le dernier donne un \emph{avis justifié} avec deux \all{weil}. Chaque étape est visible et courte : c'est exactement ce que l'examinateur attend.

\courssub{Exercice résolu}

\begin{exoresolu}[Résumer la lettre d'encouragement d'Amina]
Énoncé : Résumez en trois phrases la lettre d'encouragement qu'Amina écrit à son amie Clarisse avant l'examen d'allemand. Utilisez \all{möchte}, \all{sich freuen} et \all{weil}.

\tcblower

Solution détaillée : on applique la méthode « qui + quoi/pourquoi + quelle demande ».

\all{Amina schreibt einen Ermutigungsbrief an ihre Freundin Clarisse. Sie möchte ihr Mut machen, weil Clarisse vor der Deutschprüfung nervös ist. Sie würde sich freuen, wenn Clarisse ruhig bleibt und an sich glaubt.}

« Amina écrit une lettre d'encouragement à son amie Clarisse. Elle veut lui donner du courage parce que Clarisse est nerveuse avant l'examen d'allemand. Elle serait contente si Clarisse reste calme et croit en elle. »

Points de méthode : la phrase 1 répond à « qui écrit à qui » ; la phrase 2 répond à « quoi et pourquoi » avec \all{weil} + verbe final (\all{ist}) ; la phrase 3 exprime le souhait avec \all{würde sich freuen}. Aucune phrase n'est recopiée du texte d'origine.
\end{exoresolu}

\begin{exercice}[À votre tour : commenter la lettre d'Amina]
Rédigez un commentaire de 6 à 8 phrases sur la lettre d'encouragement d'Amina en suivant les cinq étapes : présentation (type de texte, auteure, destinataire), résumé en deux phrases, analyse d'une formule de politesse ou d'encouragement, intention de l'auteure, et avis personnel justifié avec \all{weil}. Aidez-vous du tableau des Redemittel.
\end{exercice}

\courssub{Entraînement oral : présenter la lettre en une minute}

À l'épreuve orale, on vous demandera souvent de présenter un document en une minute environ. Une minute, c'est environ huit à dix phrases simples : c'est exactement la longueur d'un commentaire en cinq étapes.

\begin{methode}[Présenter un texte à l'oral en une minute]
\begin{enumerate}
\item Préparez mentalement les cinq étapes : une phrase pour la présentation, deux pour le résumé, deux pour l'analyse, une pour l'intention, deux pour l'avis.
\item Parlez lentement et articulez ; commencez toujours par \all{Der Text ist...} ou \all{Das Dokument ist...}.
\item Utilisez des connecteurs simples pour marquer les étapes : \all{zuerst} (d'abord), \all{dann} (ensuite), \all{am Ende} (pour finir).
\item Terminez par votre avis justifié : \all{Ich finde den Brief..., weil...}.
\end{enumerate}
\end{methode}

\begin{experience}[Simulation d'épreuve orale en binôme]
Par groupes de deux : l'élève A présente la lettre d'Amina en une minute (type de texte, auteure, destinataire, résumé, une formule remarquable, avis). L'élève B chronomètre et note sur une grille : les cinq étapes sont-elles présentes ? L'avis est-il justifié avec \all{weil} ? Puis on échange les rôles avec la lettre de Merveille. Le professeur interroge ensuite deux ou trois binômes devant la classe.
\end{experience}

\begin{aretenir}[L'essentiel de la méthode]
\begin{itemize}
\item Un commentaire suit cinq étapes dans l'ordre : \cle{présentation}, \cle{résumé}, \cle{analyse}, \cle{intention}, \cle{avis justifié}.
\item Le résumé répond à trois questions : qui ? quoi et pourquoi ? quelle demande ? — en trois phrases maximum, sans recopier.
\item Les phrases utiles à mémoriser : \all{Der Text ist ein Bewerbungsschreiben.} / \all{Die Autorin möchte...} / \all{Am Ende kann man sagen, dass...} / \all{Ich finde den Brief überzeugend, weil...}.
\item Avec \all{weil}, le verbe conjugué va à la fin ; avec \all{denn}, il reste en deuxième position.
\item À l'oral, une minute = environ dix phrases simples, une ou deux par étape.
\end{itemize}
\end{aretenir}

Vous savez désormais commenter et résumer une lettre de motivation ou d'encouragement, à l'écrit comme à l'oral. La dernière section vous proposera une production guidée complète et un entraînement à la traduction pour consolider l'ensemble de la leçon.

\coursec{Production guidée et entraînement à la traduction}

Après avoir appris à lire et à commenter une lettre de motivation, il est temps de passer à la pratique : c'est en \emph{écrivant} et en \emph{traduisant} que les règles deviennent des réflexes. Cette dernière étape de la leçon te guide pas à pas : d'abord une production écrite guidée avec un plan en cinq blocs, puis une variante informelle (la lettre d'encouragement), enfin un entraînement à la traduction (thème et version) avec correction collective des erreurs typiques des francophones.

\courssub{Production écrite guidée : la lettre de motivation}

\begin{methode}[Rédiger une lettre de motivation en 5 étapes]
Pour rédiger une courte lettre de motivation (8 à 10 lignes) pour un stage ou une bourse, suis toujours le même plan :
\begin{enumerate}
\item \textbf{Anrede formelle} : commence par \all{Sehr geehrte Damen und Herren,} si tu ne connais pas le destinataire, ou \all{Sehr geehrter Herr Müller,} / \all{Sehr geehrte Frau Ngono,} si tu connais son nom. La phrase suivante commence par une \emph{minuscule}.
\item \textbf{Se présenter} : nom, âge, classe, école. Exemple : \all{Mein Name ist Amina Mbarga und ich besuche die Klasse Première A am Gymnasium Yaoundé.}
\item \textbf{Expliquer sa motivation} : pourquoi ce stage, cette bourse, cette entreprise ? Exemple : \all{Ich interessiere mich sehr für Sprachen und für Ihre Schule.}
\item \textbf{Demander poliment} avec \all{möchte} ou \all{würde} : \all{Ich möchte mich um dieses Stipendium bewerben.} / \all{Ich würde mich sehr über eine positive Antwort freuen.}
\item \textbf{Grußformel} : termine par \all{Mit freundlichen Grüßen} (sans virgule !), puis ta signature.
\end{enumerate}
\end{methode}

\begin{popfigure}
\begin{tikzpicture}[node distance=0.45cm,
  bloc/.style={draw=popblue, fill=popblueL, rounded corners, thick, text width=10.5cm, align=left, font=\small, inner sep=6pt},
  num/.style={circle, draw=popdark, fill=poporange, text=white, font=\bfseries\small, inner sep=2.5pt}]
\node[bloc] (b1) {\textbf{1. Anrede} — \all{Sehr geehrte Damen und Herren,}};
\node[num, left=0.25cm of b1] {1};
\node[bloc, below=of b1] (b2) {\textbf{2. Présentation} — \all{Mein Name ist …, ich besuche die Klasse …}};
\node[num, left=0.25cm of b2] {2};
\node[bloc, below=of b2] (b3) {\textbf{3. Motivation} — \all{Ich interessiere mich für …, weil …}};
\node[num, left=0.25cm of b3] {3};
\node[bloc, below=of b3] (b4) {\textbf{4. Demande polie} — \all{Ich möchte mich um … bewerben. Ich würde mich über … freuen.}};
\node[num, left=0.25cm of b4] {4};
\node[bloc, below=of b4] (b5) {\textbf{5. Grußformel} — \all{Mit freundlichen Grüßen} + signature};
\node[num, left=0.25cm of b5] {5};
\draw[-{Stealth}, thick, popdark] (b1.south) -- (b2.north);
\draw[-{Stealth}, thick, popdark] (b2.south) -- (b3.north);
\draw[-{Stealth}, thick, popdark] (b3.south) -- (b4.north);
\draw[-{Stealth}, thick, popdark] (b4.south) -- (b5.north);
\end{tikzpicture}
\legende{Plan-type de la lettre de motivation en 5 blocs numérotés}
\end{popfigure}

\begin{propriete}[La place de l'infinitif après « möchte » et « würde »]
Après \all{möchte} et \all{würde}, le verbe conjugué occupe la \textbf{deuxième place} de la phrase et l'\textbf{infinitif va en toute fin de phrase}, même si la phrase est longue. C'est la règle de la « cadre » (\all{Satzklammer}) : le verbe conjugué ouvre le cadre, l'infinitif le ferme.
\begin{itemize}
\item \all{Ich möchte mich um diese Stelle bewerben.} « Je voudrais postuler à ce poste. »
\item \all{Ich würde mich sehr über eine Einladung zu einem Gespräch freuen.} « Je serais très content d'être invité à un entretien. » (l'infinitif \all{freuen} reste à la fin malgré le long complément)
\end{itemize}
\textbf{Attention :} avec un verbe \emph{réfléchi}, le pronom se place juste après le verbe conjugué (ou après le sujet), et l'infinitif du verbe principal va à la fin : \all{Ich möchte mich bei Ihnen bewerben.} (le pronom \all{mich} suit \all{möchte}, l'infinitif \all{bewerben} ferme la phrase). Pour un verbe à particule séparable (ex. \all{anrufen}), c'est le verbe \emph{entier} qui va à la fin : \all{Ich möchte dich anrufen.}
\end{propriete}

\begin{exemplebox}[Un modèle complet de lettre de motivation]
Observe ce modèle, puis repère les 5 blocs du plan :

\all{Sehr geehrte Damen und Herren,}

\all{mein Name ist Junior Etoundi, ich bin 17 Jahre alt und ich besuche die Klasse Première A am Gymnasium Douala. Ich lerne seit drei Jahren Deutsch und ich interessiere mich sehr für die deutsche Sprache und Kultur.}

\all{Ich möchte mich um ein Praktikum in Ihrer Firma bewerben, weil ich später als Übersetzer arbeiten möchte. Ich bin motiviert, pünktlich und ich arbeite gern im Team.}

\all{Ich würde mich sehr über eine positive Antwort freuen.}

\all{Mit freundlichen Grüßen}

\all{Junior Etoundi}

« Mesdames, Messieurs, je m'appelle Junior Etoundi, j'ai 17 ans et je suis en classe de Première A au lycée de Douala. J'apprends l'allemand depuis trois ans et je m'intéresse beaucoup à la langue et à la culture allemandes. Je voudrais postuler à un stage dans votre entreprise, parce que je voudrais plus tard travailler comme traducteur. Je suis motivé, ponctuel et j'aime travailler en équipe. Je serais très content de recevoir une réponse positive. Cordialement, Junior Etoundi. »
\end{exemplebox}

\begin{attention}[Le pronom réfléchi de « sich bewerben » s'accorde !]
Dans \all{sich bewerben} (postuler), le pronom réfléchi change selon la personne : \all{ich bewerbe \textbf{mich}}, \all{du bewirbst \textbf{dich}}, \all{er/sie bewirbt \textbf{sich}}, \all{wir bewerben \textbf{uns}}, \all{ihr bewerbt \textbf{euch}}, \all{sie/Sie bewerben \textbf{sich}}. Erreur typique du francophone : oublier le pronom (\all{Ich möchte um diese Stelle bewerben.}) ou le laisser invariable (\all{ich bewerbe sich}). De plus, on postule \emph{à} quelque chose : \all{sich bewerben \textbf{um} + Akkusativ} : \all{Ich bewerbe mich um die Stelle.} « Je postule au poste. »
\end{attention}

\courssub{Variante : la lettre d'encouragement (registre informel)}

La lettre d'encouragement (\all{der Ermutigungsbrief}) s'adresse à un(e) ami(e) : on utilise donc le registre \emph{informel}, avec \all{du}, une Anrede familière et une Grußformel amicale.

\begin{center}
\begin{tabularx}{\linewidth}{|X|X|X|}
\hline
\rowcolor{popblueL} \textbf{Élément} & \textbf{Lettre formelle} & \textbf{Lettre informelle} \\
\hline
Anrede & \all{Sehr geehrte Damen und Herren,} & \all{Liebe Sandra,} / \all{Lieber Paul,} \\
\hline
Pronom & \all{Sie} (majuscule) & \all{du}, \all{dich}, \all{dein} \\
\hline
Ton & poli, distant (\all{möchte}, \all{würde}) & chaleureux, direct (\all{Du schaffst das!}) \\
\hline
Grußformel & \all{Mit freundlichen Grüßen} & \all{Viele Grüße} / \all{Deine Sandra} / \all{Bis bald!} \\
\hline
\end{tabularx}
\end{center}

\begin{exemplebox}[Modèle de lettre d'encouragement]
\all{Liebe Aïcha,}

\all{du hast nächste Woche deine Deutschprüfung. Keine Angst! Du hast viel gelernt und du bist gut vorbereitet. Ich denke an dich und ich drücke dir die Daumen. Du schaffst das! Nach der Prüfung feiern wir zusammen.}

\all{Viele liebe Grüße}

\all{Deine Larissa}

« Chère Aïcha, tu as ton examen d'allemand la semaine prochaine. Pas de panique ! Tu as beaucoup travaillé et tu es bien préparée. Je pense à toi et je croise les doigts pour toi. Tu vas réussir ! Après l'examen, on fêtera ça ensemble. Grosses bises, ta Larissa. »

Expressions utiles : \all{Ich drücke dir die Daumen.} « Je croise les doigts pour toi. » / \all{Du schaffst das!} « Tu vas y arriver ! » / \all{Keine Angst!} « Pas de panique ! » / \all{Viel Erfolg!} « Bonne chance ! »
\end{exemplebox}

\begin{experience}[À toi d'écrire !]
Choisis \textbf{une} des deux tâches et rédige 8 à 10 lignes :
\begin{enumerate}
\item Tu rédiges une lettre de motivation pour une bourse d'études (\all{das Stipendium}) en Allemagne. Respecte le plan en 5 blocs et utilise au moins une fois \all{möchte} et une fois \all{würde}.
\item Ton ami Rodrigue passe le concours d'entrée à l'ENAM la semaine prochaine. Écris-lui une lettre d'encouragement au registre informel (\all{du}) avec au moins deux expressions d'encouragement.
\end{enumerate}
Échange ensuite ta lettre avec ton voisin et évalue-la avec la grille ci-dessous.
\end{experience}

\begin{methode}[Grille d'auto-évaluation]
Relis ta lettre et coche chaque critère :
\begin{enumerate}
\item \textbf{Anrede correcte ?} (\all{Sehr geehrte…} ou \all{Liebe/Lieber…} selon le destinataire ; minuscule en début de phrase suivante)
\item \textbf{« möchte » / « würde » employés} pour la demande polie (lettre formelle) ?
\item \textbf{Infinitif en fin de phrase} après \all{möchte} / \all{würde} ?
\item \textbf{Pronom réfléchi accordé} dans \all{sich bewerben} (\all{ich bewerbe mich}) ?
\item \textbf{Grußformel adaptée} au registre (\all{Mit freundlichen Grüßen} / \all{Viele Grüße}) et signature ?
\end{enumerate}
Un « oui » à chaque question = une lettre réussie !
\end{methode}

\courssub{Entraînement à la traduction : thème}

\begin{exoresolu}[Thème : du français vers l'allemand]
Traduis en allemand :
\begin{enumerate}
\item Je voudrais postuler à ce poste.
\item Je serais content de recevoir une réponse.
\item Nous postulons à une bourse en Allemagne.
\item Vous trouverez mon CV en pièce jointe.
\item Je vous remercie de votre réponse.
\end{enumerate}
\tcblower
\textbf{Solutions détaillées :}
\begin{enumerate}
\item \all{Ich möchte mich um diese Stelle bewerben.} — \all{möchte} en 2\textsuperscript{e} position, pronom réfléchi \all{mich} juste après, préposition \all{um} + Akkusativ (\all{diese Stelle}), infinitif \all{bewerben} en toute fin.
\item \all{Ich würde mich über eine Antwort freuen.} — construction idiomatique \all{sich freuen über} + Akkusativ ; \all{würde} en 2\textsuperscript{e} position, infinitif \all{freuen} à la fin.
\item \all{Wir bewerben uns um ein Stipendium in Deutschland.} — pronom réfléchi \all{uns} (1\textsuperscript{re} personne du pluriel) ; \all{das Stipendium} est neutre : \all{ein Stipendium}.
\item \all{Meinen Lebenslauf finden Sie im Anhang.} ou \all{Sie finden meinen Lebenslauf im Anhang.} — \all{der Lebenslauf} à l'accusatif masculin : \all{meine\textbf{n} Lebenslauf}.
\item \all{Ich danke Ihnen für Ihre Antwort.} — \all{danken} se construit avec le \emph{datif} : \all{Ihnen}, pas \all{Sie} ; \all{für} + Akkusativ : \all{Ihre Antwort}.
\end{enumerate}
\end{exoresolu}

\begin{exercice}[Thème : à ton tour]
Traduis en allemand :
\begin{enumerate}
\item Je voudrais faire un stage dans votre entreprise.
\item Tu réussiras ton examen, j'en suis sûr !
\item Elle postule à un poste de traductrice.
\item Je serais très content d'une invitation à un entretien.
\item Nous vous remercions beaucoup.
\end{enumerate}
\end{exercice}

\begin{corrige}[Exercice : thème]
\begin{enumerate}
\item \all{Ich möchte ein Praktikum in Ihrer Firma machen.} (\all{das Praktikum}, neutre ; \all{Ihrer Firma} : datif féminin après \all{in})
\item \all{Du schaffst deine Prüfung, da bin ich sicher!} (registre informel : \all{du})
\item \all{Sie bewirbt sich um eine Stelle als Übersetzerin.} (pronom \all{sich}, 3\textsuperscript{e} personne)
\item \all{Ich würde mich sehr über eine Einladung zu einem Gespräch freuen.} (infinitif \all{freuen} en toute fin malgré la longueur)
\item \all{Wir danken Ihnen sehr.} (datif \all{Ihnen})
\end{enumerate}
\end{corrige}

\courssub{Entraînement à la traduction : version}

\begin{textebox}[Extrait d'une lettre formelle à traduire]
\all{Sehr geehrter Herr Direktor,}

\all{ich habe Ihre Anzeige in der Zeitung gelesen und ich möchte mich um die Stelle als Verkäufer bewerben. Ich bin 18 Jahre alt und ich habe mein Abitur dieses Jahr gemacht. Ich spreche Französisch, Englisch und ein bisschen Deutsch. Ich bin freundlich und ich arbeite gern mit Kunden.}

\all{Ich würde mich sehr freuen, wenn Sie mich zu einem Gespräch einladen würden. Meinen Lebenslauf finden Sie im Anhang.}

\all{Mit freundlichen Grüßen}

\all{Chantal Ngo Bell}

\textbf{Glossaire :} die Anzeige, -n (l'annonce) ; die Zeitung, -en (le journal) ; der Verkäufer, - (le vendeur) ; das Abitur (le baccalauréat) ; der Kunde, -n (le client) ; einladen (inviter) ; der Anhang, „e (la pièce jointe).
\end{textebox}

\begin{exoresolu}[Version : traduis l'extrait en français]
Traduis le texte ci-dessus en français soigné.
\tcblower
« Monsieur le Directeur, j'ai lu votre annonce dans le journal et je voudrais postuler au poste de vendeur. J'ai 18 ans et j'ai passé mon baccalauréat cette année. Je parle français, anglais et un peu allemand. Je suis aimable et j'aime travailler avec les clients. Je serais très heureuse que vous m'invitiez à un entretien. Vous trouverez mon CV en pièce jointe. Cordialement, Chantal Ngo Bell. »

Points de vigilance : \all{ich möchte mich … bewerben} = « je voudrais postuler » (et non « je veux me postuler ») ; \all{ich würde mich freuen, wenn…} = « je serais heureuse que… » (subjonctif en français) ; \all{Mit freundlichen Grüßen} se traduit par « Cordialement » ou « Veuillez agréer… », jamais mot à mot.
\end{exoresolu}

\courssub{Correction collective : les erreurs typiques}

Voici les fautes les plus fréquentes des élèves francophones. Sache les repérer et les corriger :

\begin{center}
\begin{tabularx}{\linewidth}{|X|X|X|}
\hline
\rowcolor{popblueL} \textbf{Erreur typique} & \textbf{Correction} & \textbf{Règle} \\
\hline
\all{Ich möchte bewerben mich um die Stelle.} & \all{Ich möchte mich um die Stelle bewerben.} & L'infinitif en toute fin de phrase. \\
\hline
\all{Ich bewerbe um die Stelle.} & \all{Ich bewerbe \textbf{mich} um die Stelle.} & Ne pas oublier le pronom réfléchi. \\
\hline
\all{Wir bewerben sich um ein Stipendium.} & \all{Wir bewerben \textbf{uns}…} & Le pronom s'accorde avec la personne. \\
\hline
\all{Sehr geehrte Damen und Herren, Ich möchte…} & \all{…Herren,} puis \all{ich möchte…} & Minuscule après l'Anrede. \\
\hline
\all{Mit freundlichen Grüßen,} (virgule) & \all{Mit freundlichen Grüßen} & Pas de virgule après la Grußformel. \\
\hline
\all{Ich danke Sie.} & \all{Ich danke \textbf{Ihnen}.} & \all{danken} + datif. \\
\hline
\all{Lieber Paul,} puis \all{Sie} & \all{Lieber Paul,} puis \all{du} & Cohérence du registre : pas de mélange \all{du}/\all{Sie}. \\
\hline
\end{tabularx}
\end{center}

\begin{attention}[Le mélange des registres, faute grave]
La faute la plus pénalisée à l'examen est le mélange des registres : commencer par \all{Liebe Frau Ngono,} et continuer avec \all{Sie}, ou écrire \all{Sehr geehrte Damen und Herren} et terminer par \all{Viele Grüße}. Choisis ton registre au départ (formel = \all{Sie} + \all{Mit freundlichen Grüßen} ; informel = \all{du} + \all{Viele Grüße}) et tiens-toi-y jusqu'à la signature.
\end{attention}

\begin{savaistu}[La lettre au Bac : des points faciles à prendre]
Au baccalauréat camerounais, la production écrite demande très souvent une lettre (lettre à un ami, lettre de motivation, lettre au directeur). Les correcteurs attribuent des points précis à la \emph{forme} : une Anrede correcte, une Grußformel adaptée et une signature rapportent des points « faciles », indépendamment du contenu. Un élève qui maîtrise le plan en 5 blocs et les formules \all{möchte} / \all{würde} part donc avec une longueur d'avance. Entraîne-toi à rédiger une lettre complète en 15 minutes : c'est un excellent investissement !
\end{savaistu}

\begin{aretenir}[L'essentiel de la section]
\begin{itemize}
\item La lettre de motivation suit un plan en 5 blocs : Anrede → présentation → motivation → demande polie → Grußformel.
\item \all{möchte} et \all{würde} : verbe conjugué en 2\textsuperscript{e} position, infinitif en toute fin de phrase.
\item \all{sich bewerben um} + Akkusativ, avec un pronom réfléchi accordé : \all{ich bewerbe mich}, \all{wir bewerben uns}.
\item Registre formel : \all{Sehr geehrte…}, \all{Sie}, \all{Mit freundlichen Grüßen} ; registre informel : \all{Liebe/Lieber…}, \all{du}, \all{Viele Grüße}. Jamais de mélange !
\item En traduction : \all{danken} + datif, \all{sich freuen über} + Akkusativ, et on ne traduit jamais les formules de politesse mot à mot.
\end{itemize}
\end{aretenir}

\newpage

\coursec{Activité d'intégration}

\coursec{Leçon ALA19 : Les lettres de motivation et d'encouragement}

\courssub{Objectifs de l'activité}

À la fin de cette activité d'intégration, tu seras capable de :

\begin{itemize}
\item comprendre une petite annonce de stage rédigée en allemand et en dégager les informations essentielles (qui cherche ? quoi ? quand ? conditions ?) ;
\item rédiger une lettre de motivation formelle complète (10 à 12 lignes) en respectant la structure de la lettre officielle : \cle{die Anrede}, \cle{der Betreff}, présentation, motivation, \cle{die Grußformel} ;
\item employer correctement le vocabulaire de la candidature : \all{die Bewerbung}, \all{der Lebenslauf}, \all{die Stelle}, \all{sich bewerben um} ;
\item utiliser les formes de politesse \all{möchte} (au moins deux fois) et \all{würde} (au moins une fois) dans des phrases correctes ;
\item rédiger en réponse une courte lettre d'encouragement informelle avec le tutoiement et une formule de clôture amicale ;
\item lire une lettre à voix haute avec une prononciation claire et évaluer la production d'un camarade selon des critères simples.
\end{itemize}

\courssub{Situation}

Tu es élève en classe de Première A au Lycée bilingue de Yaoundé. En feuilletant le journal mural du club d'allemand, tu découvres une annonce affichée par ton professeur. Le Goethe-Institut Kamerun propose des stages d'été aux lycéens motivés. Tu as toujours rêvé de travailler dans un environnement germanophone : c'est l'occasion parfaite de mettre en pratique tout ce que tu as appris dans cette leçon. Tu décides de poser ta candidature. Ta camarade Amina, elle aussi, hésite : elle trouve son allemand « pas assez bon ». Après avoir écrit ta lettre, tu lui écriras un mot pour l'encourager.

\begin{textebox}[Annonce — Das Goethe-Institut Kamerun sucht Praktikanten]
\textbf{Praktikum im Sommer — Bewirb dich jetzt!}

Das Goethe-Institut Kamerun in Yaoundé sucht für die Monate Juli und August motivierte Schülerinnen und Schüler (ab 16 Jahren) für ein Praktikum.

\textbf{Deine Aufgaben:} Hilfe bei der Sommerschule für Kinder, Arbeit in der Bibliothek, Organisation von Filmnachmittagen.

\textbf{Wir erwarten:} gute Deutschkenntnisse (Niveau A2), Freude am Kontakt mit Kindern, Pünktlichkeit.

\textbf{Wir bieten:} ein Zertifikat, Deutschkurse gratis, ein nettes Team.

Interesse? Schreibe eine Bewerbung mit Lebenslauf an: \texttt{frau.mballa\@goethe.de}

Anmeldeschluss: 15. Juni.
\end{textebox}

\begin{definition}[Petit glossaire de l'annonce]
\begin{itemize}
\item \all{das Praktikum, die Praktika} : le stage ; \all{der Praktikant, -en} / \all{die Praktikantin, -nen} : le/la stagiaire
\item \all{die Aufgabe, -n} : la tâche ; \all{die Deutschkenntnisse} (pluriel) : les connaissances en allemand
\item \all{pünktlich} : ponctuel(le) ; \all{das Zertifikat, -e} : le certificat
\item \all{gratis} : gratuit ; \all{der Anmeldeschluss} : la date limite d'inscription
\item \all{sich bewerben} : poser sa candidature ; \all{die Bewerbung, -en} : la candidature, la lettre de motivation
\end{itemize}
\end{definition}

\courssub{Consignes}

Travaille d'abord seul(e), puis en binôme. Suis les étapes dans l'ordre.

\begin{enumerate}
\item \textbf{Comprendre l'annonce.} Lis le document deux fois, puis réponds en allemand par des phrases complètes : \all{Wer sucht Praktikanten? Wann beginnt das Praktikum? Was sind die Aufgaben? Was muss man schicken? Bis wann?}
\item \textbf{Repérer le lexique utile.} Souligne dans l'annonce cinq mots ou expressions du thème de la candidature et recopie-les dans ton cahier avec leur article et leur pluriel (pour les noms).
\item \textbf{Préparer le plan de ta lettre.} Note en quatre points : (a) l'en-tête (expéditeur, destinataire, lieu et date) ; (b) \all{der Betreff} et \all{die Anrede} formelle ; (c) le corps en trois moments : présentation, motivation, disponibilité ; (d) \all{die Grußformel} et la signature.
\item \textbf{Rédiger ta lettre de motivation} (10 à 12 lignes). Contraintes obligatoires : Anrede formelle \all{Sehr geehrte Frau Mballa,} ; au moins \textbf{deux} phrases avec \all{möchte} ; au moins \textbf{une} phrase avec \all{würde} ; Grußformel \all{Mit freundlichen Grüßen}. Attention à la place du verbe et à la majuscule des noms.
\item \textbf{Relecture en binôme.} Échange ta lettre avec ton voisin. Vérifie à l'aide de la grille : Anrede correcte ? deux fois \all{möchte} ? une fois \all{würde} ? Grußformel ? verbe en position 2 ? majuscules ? Signale deux erreurs et un point fort.
\item \textbf{Écrire la lettre d'encouragement.} Ta camarade Amina hésite à postuler. Rédige-lui une courte lettre informelle (5 à 6 lignes) : Anrede avec \all{Liebe Amina,}, tutoiement (\all{du}, \all{dein}), au moins deux expressions d'encouragement (\all{Du schaffst das!}, \all{Ich glaube an dich!}, \all{Kopf hoch!}), et une formule de clôture amicale (\all{Viele Grüße} / \all{Deine} + prénom).
\item \textbf{Présentation orale.} Les deux ou trois meilleures lettres sont lues à voix haute devant la classe, corrigées collectivement, puis affichées au tableau d'allemand.
\end{enumerate}

\courssub{Ressources à mobiliser}

\begin{aretenir}[La structure de la lettre formelle]
\begin{enumerate}
\item En-tête : expéditeur (nom, adresse) en haut à gauche, destinataire en dessous, lieu et date à droite : \all{Yaoundé, den 12. Mai 2025}.
\item \all{Der Betreff} : l'objet, sans verbe : \all{Bewerbung um ein Praktikum}.
\item \all{Die Anrede} : \all{Sehr geehrte Frau Mballa,} — la phrase suivante commence par une \textbf{minuscule}.
\item Le corps : présentation — motivation — disponibilité et remerciement.
\item \all{Die Grußformel} : \all{Mit freundlichen Grüßen} (sans virgule !), puis la signature.
\end{enumerate}
\end{aretenir}

\begin{propriete}[La politesse avec \all{möchte} et \all{würde}]
\all{möchte} (subjonctif II de \all{mögen}) exprime un souhait poli ; l'infinitif se place à la fin : \all{Ich möchte mich um das Praktikum bewerben.} « Je voudrais poser ma candidature pour le stage. »

\all{würde} + infinitif exprime l'hypothèse ou la politesse : \all{Ich würde mich sehr über eine positive Antwort freuen.} « Je me réjouirais beaucoup d'une réponse positive. »

Conjugaison de \all{möchte} (Präsens Konjunktiv II) : \all{ich möchte}, \all{du möchtest}, \all{er/sie/es möchte}, \all{wir möchten}, \all{ihr möchtet}, \all{sie/Sie möchten}.
\end{propriete}

\begin{attention}[Erreurs fréquentes des francophones]
\begin{itemize}
\item Pas de virgule après \all{Mit freundlichen Grüßen} (contrairement au français).
\item Après la virgule de l'Anrede, la phrase commence par une minuscule : \all{Sehr geehrte Frau Mballa, ich möchte...}
\item \all{sich bewerben um} + accusatif : \all{Ich bewerbe mich um die Stelle.} Ne dis pas \all{bewerben für}.
\item Faux ami : \all{die Bewerbung} = la candidature, pas « la boisson » (\all{das Getränk}) ; \all{das Praktikum} = le stage.
\item Lettre informelle : tutoiement obligatoire ; majuscule possible à \all{Du} / \all{Dein} dans les lettres (usage traditionnel poli).
\end{itemize}
\end{attention}

\begin{center}
\begin{tabularx}{\linewidth}{|X|X|}\hline
\rowcolor{popblueL} \textbf{Français} & \textbf{Deutsch} \\ \hline
Je voudrais poser ma candidature & Ich möchte mich bewerben \\ \hline
J'apprends l'allemand depuis trois ans & Ich lerne seit drei Jahren Deutsch \\ \hline
Je me réjouirais d'une réponse & Ich würde mich über eine Antwort freuen \\ \hline
Tu vas réussir ! & Du schaffst das ! \\ \hline
Je crois en toi & Ich glaube an dich \\ \hline
\end{tabularx}
\end{center}

\courssub{Proposition de production et corrigé commenté}

\begin{exoresolu}[Modèle de lettre de motivation]
\textbf{Énoncé.} Rédige ta lettre de motivation pour le stage du Goethe-Institut Kamerun (10 à 12 lignes) avec Anrede formelle, deux fois \all{möchte}, une fois \all{würde}, et la Grußformel.

\tcblower
\textbf{Production modèle.}

Marlyse Etoundi, Quartier Nkolbisson, Yaoundé

Goethe-Institut Kamerun, Frau Mballa, Yaoundé

Yaoundé, den 12. Mai 2025

\textbf{Bewerbung um ein Praktikum für den Sommer}

Sehr geehrte Frau Mballa,

ich heiße Marlyse Etoundi und ich bin 16 Jahre alt. Ich besuche die Klasse Première A am Lycée bilingue in Yaoundé und ich lerne seit drei Jahren Deutsch. Ich habe Ihre Anzeige gelesen und ich möchte mich um das Praktikum im Sommer bewerben.

Ich mag Kinder sehr und ich arbeite gern im Team. In der Schule helfe ich schon bei der Deutsch-AG. Ich möchte meine Deutschkenntnisse verbessern und neue Erfahrungen machen. Ich bin pünktlich und motiviert, und im Juli und August habe ich Zeit.

Ich würde mich sehr über eine positive Antwort freuen. Meinen Lebenslauf finden Sie im Anhang.

Mit freundlichen Grüßen

Marlyse Etoundi

\medskip
\textbf{Commentaire des choix de langue.}
\begin{itemize}
\item \textbf{Structure respectée} : en-tête complet, Betreff sans verbe, Anrede formelle avec minuscule après la virgule, corps en trois paragraphes (présentation — motivation — disponibilité), Grußformel sans virgule.
\item \textbf{Politesse} : deux occurrences de \all{möchte} (\all{ich möchte mich ... bewerben} ; \all{Ich möchte meine Deutschkenntnisse verbessern}) avec l'infinitif rejeté en fin de phrase ; une occurrence de \all{würde} dans la formule de clôture classique \all{Ich würde mich sehr über eine positive Antwort freuen.}
\item \textbf{Lexique du thème} : \all{die Bewerbung} (dans le Betreff), \all{sich bewerben um} + accusatif, \all{das Praktikum}, \all{die Deutschkenntnisse}, \all{der Lebenslauf}, \all{der Anhang} (la pièce jointe).
\item \textbf{Grammaire} : \all{seit} + datif avec le présent (\all{ich lerne seit drei Jahren Deutsch}) ; verbe toujours en deuxième position ; vouvoiement avec majuscule (\all{Ihre Anzeige}).
\end{itemize}

\medskip
\textbf{Modèle de lettre d'encouragement (informelle).}

Liebe Amina,

du schaffst das! Dein Deutsch ist wirklich gut, und du bist super motiviert. Schreib einfach die Bewerbung — ich helfe dir gern. Ich glaube an dich! Kopf hoch und bis morgen!

Viele Grüße

Marlyse

\medskip
\textbf{Commentaire.} Anrede informelle avec \all{Liebe}, tutoiement (\all{du}, \all{dein}, \all{dir}), deux expressions d'encouragement (\all{Du schaffst das!}, \all{Ich glaube an dich!}), impératif correct (\all{Schreib}), clôture amicale \all{Viele Grüße} sans virgule.
\end{exoresolu}

\courssub{Bilan de l'activité}

Cette activité t'a permis de mobiliser dans une situation de communication authentique l'ensemble des savoir-faire de la leçon :

\begin{itemize}
\item \textbf{Compréhension de l'écrit} : tu as lu une annonce réelle de type germanophone et tu en as extrait les informations utiles avant d'agir ;
\item \textbf{Lexique} : tu as réutilisé activement le vocabulaire de la candidature (\all{die Bewerbung}, \all{der Lebenslauf}, \all{die Stelle}, \all{sich bewerben um}) dans ta propre production ;
\item \textbf{Grammaire} : tu as appliqué la structure de la lettre formelle et informelle, ainsi que les formes de politesse \all{möchte} et \all{würde}, en surveillant la place du verbe et les majuscules ;
\item \textbf{Expression écrite et orale} : tu as produit deux textes de registres différents (formel et amical), évalué la lettre d'un camarade avec une grille de critères, et lu une production devant la classe.
\end{itemize}

\begin{savaistu}[Et après le lycée ?]
Savoir rédiger une \all{Bewerbung} est une compétence très concrète : les universités allemandes, autrichiennes et suisses, le DAAD (service allemand d'échanges universitaires) et le Goethe-Institut proposent chaque année des stages et des bourses aux jeunes Camerounais germanophones. Une lettre de motivation claire, polie et sans faute est souvent la première étape vers ces opportunités. Garde précieusement le modèle rédigé aujourd'hui dans ton portfolio d'allemand !
\end{savaistu}

\newpage

\coursec{Méthodes et exercices résolus}

\begin{textebox}[Texte support : Eine Bewerbung und ein Ermutigungsbrief]
\textbf{Text A : Lettre de motivation formelle (Bewerbung)}

\all{Yaoundé, den 15. Mai 2025}

\all{Sehr geehrte Damen und Herren,}

\all{ich heiße Amina Mbarga und wohne in Yaoundé. Ich habe Ihre Anzeige in der Zeitung gelesen und möchte mich um die Stelle als Verkäuferin bewerben. Ich habe zwei Jahre Erfahrung in einem Geschäft und ich spreche gut Französisch, Englisch und ein bisschen Deutsch. Mein Lebenslauf liegt bei. Ich freue mich auf Ihre Antwort.}

\all{Mit freundlichen Grüßen}

\all{Amina Mbarga}

\vspace{0.5cm}
\textbf{Text B : Lettre d'encouragement informelle (Ermutigungsbrief)}

\all{Douala, den 10. Juni 2025}

\all{Lieber Eric,}

\all{ich habe gehört, dass du die Prüfung nicht bestanden hast. Kopf hoch! Das ist nicht das Ende. Du bist intelligent und du arbeitest viel. Gib nicht auf und lerne jeden Tag ein bisschen Deutsch. Ich helfe dir gern. Du schaffst das beim nächsten Mal!}

\all{Viele Grüße}

\all{Dein Freund Samuel}

\vspace{0.3cm}
\textbf{Glossaire :}
\begin{itemize}
\item \all{die Anzeige, -n} : l'annonce (presse, Internet)
\item \all{sich bewerben um + Akk} : postuler à / pour
\item \all{die Stelle, -n} : le poste, la place
\item \all{die Verkäuferin, -nen} : la vendeuse
\item \all{die Erfahrung, -en} : l'expérience
\item \all{der Lebenslauf, die Lebensläufe} : le CV (curriculum vitæ)
\item \all{beiliegen} (trennbar, \all{liegt ... bei}) : être joint (document)
\item \all{sich freuen auf + Akk} : se réjouir de (attendre avec plaisir)
\item \all{Kopf hoch!} : Courage ! / Tiens bon ! (litt. « la tête haute »)
\item \all{nicht bestanden} : non réussi (Partizip II de \all{nicht bestehen})
\item \all{aufgeben} (trennbar, \all{gib auf}) : abandonner
\item \all{beim nächsten Mal} : la prochaine fois (contraction \all{bei dem})
\end{itemize}
\end{textebox}

\begin{definition}[Lexique thématique : \cle{die Bewerbung}]
\begin{center}
\begin{tabularx}{\linewidth}{|X|X|X|}
\hline
\rowcolor{popblueL} \textbf{Deutsch (Artikel + Pluriel)} & \textbf{Français} & \textbf{Construction / Exemple} \\
\hline
\all{die Anzeige, -n} & l'annonce & \all{eine Anzeige schalten / lesen} \\
\all{die Bewerbung, -en} & la candidature & \all{eine Bewerbung schreiben / senden} \\
\all{der Lebenslauf, die Lebensläufe} & le CV & \all{der Lebenslauf liegt bei} \\
\all{die Stelle, -n} & le poste, la place & \all{sich um eine Stelle bewerben} \\
\all{die Erfahrung, -en} & l'expérience & \all{Erfahrung haben / sammeln} \\
\all{die Anrede, -n} & la formule d'appel & \all{Sehr geehrte Damen und Herren} \\
\all{die Grußformel, -n} & la formule de politesse (fin) & \all{Mit freundlichen Grüßen} \\
\all{der Bewerber, die Bewerber} & le candidat & \all{der Bewerber wird eingeladen} \\
\all{der Arbeitgeber, -} & l'employeur & \all{der Arbeitgeber antwortet} \\
\hline
\end{tabularx}
\end{center}

\textbf{Familles de mots à retenir :}
\begin{itemize}
\item \all{bewerben} (verbe) $\rightarrow$ \all{die Bewerbung} (nom) $\rightarrow$ \all{der Bewerber / die Bewerberin} (agent).
\item \all{arbeiten} $\rightarrow$ \all{die Arbeit} $\rightarrow$ \all{der Arbeitgeber} / \all{der Arbeitnehmer}.
\item \all{anzeigen} $\rightarrow$ \all{die Anzeige}.
\end{itemize}

\textbf{Constructions fixes (à apprendre par cœur) :}
\begin{itemize}
\item \all{sich bewerben um + Akkusativ} : \all{Ich bewerbe mich um die Stelle.}
\item \all{sich freuen auf + Akkusativ} (attente future) : \all{Ich freue mich auf das Gespräch.}
\item \all{sich freuen über + Akkusativ} (joie présente/passée) : \all{Ich freue mich über die Zusage.}
\item \all{helfen + Dativ} : \all{Ich helfe dir / Ihnen.}
\end{itemize}
\end{definition}

\begin{propriete}[Grammaire 1 : \cle{La lettre formelle} et \cle{la lettre informelle}]
\textbf{1. Structure commune (obligatoire) :}
\begin{enumerate}
\item \textbf{Lieu et date} : en haut à droite. Format : \all{Ort, den TT. Monat Jahr} (ex. \all{Yaoundé, den 15. Mai 2025}).
\item \textbf{Formule d'appel (\cle{die Anrede})} : suivie d'une virgule.
\item \textbf{Corps de la lettre} : paragraphes séparés par une ligne blanche.
\item \textbf{Formule de politesse finale (\cle{die Grußformel})} : sans virgule à la fin.
\item \textbf{Signature} : prénom + nom (formel) ou prénom seul (informel).
\end{enumerate}

\textbf{2. Lettre formelle (geschäftlich / offiziell) :}
\begin{center}
\begin{tabular}{|l|l|}
\hline
\textbf{Élément} & \textbf{Allemand} \\
\hline
Appel (inconnu) & \all{Sehr geehrte Damen und Herren,} \\
Appel (connu, femme) & \all{Sehr geehrte Frau Ngono,} \\
Appel (connu, homme) & \all{Sehr geehrter Herr Ngono,} \\
Pronoms & \cle{\textbf{Sie, Ihnen, Ihr, Ihre}} (toujours majuscule) \\
Corps & Vouvoiement, style soutenu, pas de familier \\
Formule finale standard & \all{Mit freundlichen Grüßen} \\
Formule finale plus formelle & \all{Mit freundlichen Grüßen und besten Empfehlungen} \\
\hline
\end{tabular}
\end{center}

\textbf{3. Lettre informelle (privat / freundschaftlich) :}
\begin{center}
\begin{tabular}{|l|l|}
\hline
\textbf{Élément} & \textbf{Allemand} \\
\hline
Appel (homme) & \all{Lieber Paul,} \\
Appel (femme) & \all{Liebe Sandra,} \\
Pronoms & \cle{\textbf{du, dir, dein, dich}} (minuscule) \\
Corps & Tutoiement, ton personnel, contractions (\all{beim, zum}) \\
Formule finale amicale & \all{Viele Grüße} / \all{Liebe Grüße} \\
Signature & \all{Dein Freund / Deine Freundin} + prénom \\
\hline
\end{tabular}
\end{center}

\textbf{4. Règle d'or majuscule / minuscule après l'appel :}
Après la virgule de l'appel (\all{Sehr geehrte Damen und Herren,} ou \all{Lieber Paul,}), la phrase commence par une \textbf{minuscule} (sauf nom propre) :
\all{Sehr geehrte Damen und Herren, ich beziehe mich auf Ihre Anzeige.}
\all{Lieber Paul, ich danke dir für deinen Brief.}
\end{propriete}

\begin{propriete}[Grammaire 2 : \cle{La politesse} avec \cle{möchte} et \cle{würde} (Konjunktiv II)]
Le \cle{Konjunktiv II} exprime l'irréel, le souhait ou la politesse. En A1/A2, on maîtrise deux formes figées : \all{möchte} (de \all{mögen}) et \all{würde} (de \all{werden}).

\textbf{1. \all{möchte} = « je voudrais / j'aimerais » (volonté polie) }
Conjugaison (Präsens Konjunktiv II) :
\begin{center}
\begin{tabular}{|l|l|l|l|}
\hline
\all{ich möchte} & \all{wir möchten} & \all{ich möchte nicht} & \textbf{Négation} \\
\all{du möchtest} & \all{ihr möchtet} & & \\
\all{er/sie/es möchte} & \all{sie/Sie möchten} & & \\
\hline
\end{tabular}
\end{center}
\textbf{Structure :} \all{Subjekt + möchte + ... + Infinitif (à la fin).}
\all{Ich möchte mich um die Stelle bewerben.}
\all{Möchten Sie einen Kaffee trinken?}

\textbf{2. \all{würde} + infinitif = « je serais / je ferais / j'aimerais » (conditionnel) }
Conjugaison :
\begin{center}
\begin{tabular}{|l|l|}
\hline
\all{ich würde} & \all{wir würden} \\
\all{du würdest} & \all{ihr würdet} \\
\all{er/sie/es würde} & \all{sie/Sie würden} \\
\hline
\end{tabular}
\end{center}
\textbf{Structure :} \all{Subjekt + würde + ... + Infinitif (à la fin).}
Souvent associé à \cle{gern} pour adoucir : \all{Ich würde gern kommen.}

\textbf{3. Place des verbes (Règle fondamentale) :}
\begin{itemize}
\item Phrase affirmative : Verbe modal conjugué (\all{möchte/würde}) en \textbf{Position 2}, Infinitif à la \textbf{fin}.
\item Question fermée (sans mot interrogatif) : Verbe modal en \textbf{Position 1} (inversion), Infinitif à la \textbf{fin}.
\item Subordonnée (\all{dass, weil, wenn...}) : Verbe modal conjugué à la \textbf{fin}, Infinitif juste avant.
\end{itemize}

\textbf{4. Faux ami majeur :}
\all{bekommen} = \textbf{recevoir} / obtenir (\all{Ich bekomme ein Geschenk.}).
\all{werden} = \textbf{devenir} / \textbf{aller} (futur) (\all{Ich möchte Lehrerin werden.}).
\end{propriete}

\courssub{Méthode 1 : Comprendre une lettre de motivation (lecture globale puis détaillée)}

\begin{methode}[Lire et comprendre un texte sur la candidature]
Pour réussir les questions de compréhension au Bac, suis toujours la même démarche :
\begin{enumerate}
\item \textbf{Première lecture (globale)} : lis le texte entier sans dictionnaire. Repère le type de document (lettre formelle ? lettre amicale ?), l'expéditeur (\cle{\all{der Absender}}), le destinataire (\cle{\all{der Empfänger}}) et le but (\cle{\all{das Ziel}}) : postuler, encourager, remercier.
\item \textbf{Repérer les marqueurs de la lettre formelle} : \all{Sehr geehrte Damen und Herren}, \all{Mit freundlichen Grüßen}, la date, l'objet (\cle{\all{Betreff}}). Pour la lettre informelle : \all{Lieber...}, \all{Viele Grüße}.
\item \textbf{Deuxième lecture (détaillée)} : souligne les mots-clés du thème : \cle{\all{die Stelle}}, \cle{\all{sich bewerben}}, \cle{\all{der Lebenslauf}}, \cle{\all{die Erfahrung}}. Entoure les verbes conjugués pour identifier le temps (Präsens, Perfekt, Futur).
\item \textbf{Répondre aux questions} : réponds par une phrase complète en allemand, en reprenant les mots de la question. Ne recopie jamais tout le texte : reformule avec tes mots.
\item \textbf{Vérifier} : verbe en position 2, majuscule aux noms, accord du sujet et du verbe.
\end{enumerate}
\textbf{Réflexe} : une question avec \cle{\all{Warum}} (pourquoi) exige une réponse avec \cle{\all{weil}} (verbe à la fin) ou \cle{\all{denn}} (verbe en position 2).
\end{methode}

\begin{exoresolu}[Compréhension de texte : une lettre de motivation]
\textbf{Texte.} \all{Sehr geehrte Damen und Herren, ich heiße Amina Mbarga und wohne in Yaoundé. Ich habe Ihre Anzeige in der Zeitung gelesen und möchte mich um die Stelle als Verkäuferin bewerben. Ich habe zwei Jahre Erfahrung in einem Geschäft und ich spreche gut Französisch, Englisch und ein bisschen Deutsch. Mein Lebenslauf liegt bei. Ich freue mich auf Ihre Antwort. Mit freundlichen Grüßen, Amina Mbarga.}

\textbf{Énoncé.} 1) Wer schreibt den Brief? 2) Um welche Stelle bewirbt sie sich? 3) Warum kann sie gut für die Stelle arbeiten? (Réponds par des phrases complètes.)
\tcblower
\textbf{Raisonnement.} Question 1 : on cherche l'expéditeur (\all{der Absender}) ; le texte dit \all{ich heiße Amina Mbarga}. Question 2 : le mot-clé est \all{die Stelle als Verkäuferin}. Question 3 : question avec \all{Warum} : on répond avec \all{weil} + verbe à la fin ; les arguments sont l'expérience (\all{zwei Jahre Erfahrung}) et les langues.

\textbf{Solution rédigée.}
1) \all{Amina Mbarga schreibt den Brief. Sie wohnt in Yaoundé.} (Amina Mbarga écrit la lettre. Elle habite à Yaoundé.)

2) \all{Sie bewirbt sich um die Stelle als Verkäuferin.} (Elle postule pour le poste de vendeuse.) Attention : \cle{\all{sich bewerben um}} + accusatif ; au Präsens, 3\up{e} personne : \all{sie bewirbt sich}.

3) \all{Sie kann gut für die Stelle arbeiten, weil sie zwei Jahre Erfahrung hat und drei Sprachen spricht.} (Elle peut bien travailler pour ce poste parce qu'elle a deux ans d'expérience et parle trois langues.) Justification : subordonnée avec \cle{\all{weil}} : le verbe conjugué \all{hat} va à la fin ; coordination \all{und} : le deuxième verbe \all{spricht} reste en position 2 de sa proposition.

\textbf{Conclusion.} Chaque réponse est une phrase complète, avec verbe en position 2 (ou à la fin après \all{weil}), majuscule aux noms (\all{Brief, Stelle, Erfahrung}) et reprise des mots de la question.
\end{exoresolu}

\courssub{Méthode 2 : Employer le vocabulaire de la candidature}

\begin{methode}[Maîtriser le lexique de \cle{\all{die Bewerbung}}]
\begin{enumerate}
\item \textbf{Apprends chaque nom avec son article et son pluriel} : \cle{\all{die Bewerbung, -en}} ; \cle{\all{der Lebenslauf, die Lebensläufe}} ; \cle{\all{die Stelle, -n}} ; \cle{\all{die Anzeige, -n}} ; \cle{\all{die Erfahrung, -en}} ; \cle{\all{die Anrede, -n}} ; \cle{\all{die Grußformel, -n}}.
\item \textbf{Retiens les constructions fixes} : \cle{\all{sich bewerben um}} + accusatif (postuler à) ; \cle{\all{sich freuen auf}} + accusatif (attendre avec plaisir) ; \cle{\all{sich freuen über}} + accusatif (se réjouir de).
\item \textbf{Utilise les familles de mots} : \all{bewerben} $\rightarrow$ \all{die Bewerbung} $\rightarrow$ \all{der Bewerber} (le candidat) ; \all{arbeiten} $\rightarrow$ \all{die Arbeit} $\rightarrow$ \all{der Arbeitgeber} (l'employeur).
\item \textbf{Distingue formel et informel} : dans une lettre formelle, on vouvoie : \cle{\all{Sie, Ihre, Ihnen}} (toujours avec majuscule) ; dans une lettre amicale : \all{du, dein, dir}.
\item \textbf{Mémorise les formules toutes faites} : \all{Ich möchte mich um die Stelle als ... bewerben.} / \all{Ich freue mich auf Ihre Antwort.}
\end{enumerate}
\end{methode}

\begin{exoresolu}[Lexique : compléter une lettre]
\textbf{Énoncé.} Complète avec : \all{Stelle, Bewerbung, Lebenslauf, Erfahrung, Anzeige}.

\all{Liebe Frau Kamdem, ich habe Ihre (1) \dots\ im Internet gesehen. Ich möchte mich um die (2) \dots\ als Sekretärin bewerben. Ich habe drei Jahre (3) \dots\ in einem Büro in Douala. Mein (4) \dots\ liegt bei. Ich hoffe, dass meine (5) \dots\ Sie interessiert.}
\tcblower
\textbf{Raisonnement.} (1) On « voit » une annonce dans un journal ou sur Internet : \cle{\all{die Anzeige}}. (2) Après \cle{\all{sich bewerben um die}}, il faut le poste : \cle{\all{die Stelle}}. (3) Avec « trois ans de... », c'est l'expérience : \cle{\all{die Erfahrung}} (ici sans pluriel, sens général). (4) Le document joint qui décrit le parcours : \cle{\all{der Lebenslauf}}. (5) Le dossier complet envoyé à l'employeur : \cle{\all{die Bewerbung}}.

\textbf{Solution rédigée.}
1) \all{Anzeige} (l'annonce) — 2) \all{Stelle} (le poste) — 3) \all{Erfahrung} (l'expérience) — 4) \all{Lebenslauf} (le CV) — 5) \all{Bewerbung} (la candidature).

Phrase complète vérifiée : \all{Ich habe Ihre Anzeige im Internet gesehen. Ich möchte mich um die Stelle als Sekretärin bewerben.} (J'ai vu votre annonce sur Internet. Je voudrais postuler pour le poste de secrétaire.)

\textbf{Conclusion.} Vérifie toujours le genre demandé par l'article de la phrase (\all{Ihre} peut cacher un féminin ou un pluriel) et la majuscule obligatoire à chaque nom.
\end{exoresolu}

\courssub{Méthode 3 : Rédiger une lettre formelle et une lettre informelle}

\begin{methode}[Construire une lettre correcte]
\begin{enumerate}
\item \textbf{Structure obligatoire} : lieu et date en haut à droite (\all{Yaoundé, den 15. Mai}) ; formule d'appel (\cle{\all{die Anrede}}) ; corps de la lettre ; formule de politesse (\cle{\all{die Grußformel}}) ; signature.
\item \textbf{Lettre formelle} : \all{Sehr geehrte Damen und Herren,} ou \all{Sehr geehrte Frau Ngono,} ; vouvoiement \cle{\all{Sie}} avec majuscule ; conclusion : \all{Mit freundlichen Grüßen}.
\item \textbf{Lettre informelle} (lettre d'encouragement à un ami) : \all{Lieber Paul,} / \all{Liebe Sandra,} ; tutoiement \all{du} ; conclusion : \all{Viele Grüße} ou \all{Dein Freund / Deine Freundin}.
\item \textbf{Après l'appel}, la phrase commence par une minuscule (sauf nom) : \all{Lieber Paul, ich danke dir für deinen Brief.}
\item \textbf{Pour encourager}, utilise : \cle{\all{Kopf hoch!}} (courage !), \all{Ich glaube an dich.} (je crois en toi), \cle{\all{Du schaffst das!}} (tu vas y arriver !), \cle{\all{Gib nicht auf!}} (n'abandonne pas !).
\end{enumerate}
\end{methode}

\begin{exoresolu}[Rédaction : une lettre d'encouragement]
\textbf{Énoncé.} Ton ami Eric a échoué à un examen d'allemand. Écris-lui une courte lettre informelle (5-6 phrases) pour l'encourager. Utilise : \all{Lieber...}, \all{Kopf hoch}, \all{du schaffst das}, \all{Viele Grüße}.
\tcblower
\textbf{Raisonnement.} Lettre informelle : tutoiement, appel \all{Lieber Eric,}, minuscule après la virgule, verbe en position 2 dans chaque phrase, impératif avec \all{du} pour encourager (\all{Gib nicht auf!}).

\textbf{Solution rédigée.}

\all{Yaoundé, den 15. Mai}

\all{Lieber Eric,}

\all{ich habe gehört, dass du die Prüfung nicht bestanden hast. Kopf hoch! Das ist nicht das Ende. Du bist intelligent und du arbeitest viel. Gib nicht auf und lerne jeden Tag ein bisschen Deutsch. Ich helfe dir gern. Du schaffst das beim nächsten Mal!}

\all{Viele Grüße}

\all{Dein Freund Samuel}

(J'ai entendu que tu n'as pas réussi l'examen. Courage ! Ce n'est pas la fin. Tu es intelligent et tu travailles beaucoup. N'abandonne pas et apprends un peu d'allemand chaque jour. Je t'aide volontiers. Tu vas y arriver la prochaine fois !)

\textbf{Justifications grammaticales.} Subordonnée avec \all{dass} : verbe \all{hast} à la fin. \all{helfen} + datif : \all{ich helfe dir}. Impératif : \all{Gib} (verbe fort \all{geben}, impératif irrégulier avec changement e $\rightarrow$ i). Après \all{Lieber Eric,} : minuscule à \all{ich}.

\textbf{Conclusion.} La lettre respecte la structure : date, appel, corps, formule finale, signature.
\end{exoresolu}

\courssub{Méthode 4 : Traduire avec la politesse (\cle{\all{möchte}} et \cle{\all{würde}})}

\begin{methode}[Traduire « je voudrais » et exprimer la politesse]
\begin{enumerate}
\item \textbf{« Je voudrais »} se traduit par \cle{\all{ich möchte}} + infinitif à la fin : \all{Ich möchte mich bewerben.} (Je voudrais postuler.) Jamais \all{ich will} (trop direct, impoli).
\item \textbf{Conjugaison de \cle{\all{möchte}} (Konjunktiv II de \all{mögen})} : \all{ich möchte, du möchtest, er/sie/es möchte, wir möchten, ihr möchtet, sie/Sie möchten}.
\item \textbf{« Je serais / il serait / j'aimerais »} : \cle{\all{ich würde}} + infinitif : \all{Ich würde gern in Deutschland arbeiten.} (J'aimerais travailler en Allemagne.) Conjugaison : \all{ich würde, du würdest, er würde, wir würden, ihr würdet, sie würden}.
\item \textbf{Place du verbe} : le modal conjugué est en position 2, l'infinitif va à la fin de la phrase : \all{Ich möchte heute einen Brief schreiben.}
\item \textbf{Pour adoucir}, ajoute \cle{\all{gern}} : \all{Ich würde gern mehr erfahren.} (J'aimerais en savoir plus.)
\end{enumerate}
\end{methode}

\begin{exoresolu}[Thème : traduction du français vers l'allemand]
\textbf{Énoncé.} Traduis en allemand : 1) Je voudrais postuler pour ce poste. 2) Nous aimerions travailler dans votre entreprise. 3) Seriez-vous d'accord pour un entretien ? (utilise \all{würden Sie einverstanden sein...})
\tcblower
\textbf{Raisonnement.} Phrase 1 : « je voudrais » = \cle{\all{ich möchte}} ; « postuler pour » = \cle{\all{sich bewerben um}} + accusatif ; l'infinitif réflexif \all{mich bewerben} va à la fin. Phrase 2 : « nous aimerions » = \cle{\all{wir würden gern}} + infinitif \all{arbeiten} ; « dans votre entreprise » = \all{in Ihrem Unternehmen} (datif après \all{in} pour le lieu, \all{das Unternehmen} $\rightarrow$ \all{Ihrem}). Phrase 3 : question sans mot interrogatif : inversion, le verbe \all{würden} en position 1 ; \all{einverstanden sein} = être d'accord ; \all{mit einem Gespräch} (datif après \all{mit}).

\textbf{Solution rédigée.}
1) \all{Ich möchte mich um diese Stelle bewerben.} (modal en position 2, infinitif réflexif à la fin, \all{um} + accusatif : \all{diese Stelle}.)

2) \all{Wir würden gern in Ihrem Unternehmen arbeiten.} (\all{würden} + \all{gern} pour la politesse ; \all{in} + datif : \all{Ihrem Unternehmen}.)

3) \all{Würden Sie mit einem Gespräch einverstanden sein?} (inversion du modal en tête de question ; \all{Sie} avec majuscule ; infinitif \all{sein} à la fin.)

\textbf{Conclusion.} La règle d'or : modal conjugué en position 2 (ou 1 dans une question), infinitif renvoyé à la fin de la phrase.
\end{exoresolu}

\courssub{Méthode 5 : Résumer et produire un court paragraphe}

\begin{methode}[Résumer un texte et écrire un paragraphe guidé]
\begin{enumerate}
\item \textbf{Pour résumer} : relève l'idée principale de chaque partie (qui ? quoi ? pourquoi ?) et reformule en 2-3 phrases avec tes propres mots, au Präsens.
\item \textbf{Utilise des connecteurs simples} : \cle{\all{zuerst}} (d'abord), \cle{\all{dann}} (ensuite), \cle{\all{auch}} (aussi), \cle{\all{deshalb}} (c'est pourquoi), \all{und}, \all{aber}.
\item \textbf{Attention à \cle{\all{deshalb}}} : c'est un connecteur qui compte comme position 1 : le verbe suit immédiatement : \all{Sie hat Erfahrung, deshalb bewirbt sie sich.} (et non « deshalb sie bewirbt sich »).
\item \textbf{Pour un paragraphe de production} (ex. ta propre mini-lettre de motivation) : 1 phrase d'introduction (qui tu es), 2 phrases d'arguments (expérience, qualités), 1 phrase de conclusion (politesse).
\item \textbf{Relis-toi} : verbe conjugué en position 2, infinitifs à la fin, majuscules aux noms, \cle{\all{Sie}} avec majuscule dans une lettre formelle.
\end{enumerate}
\end{methode}

\begin{exoresolu}[Production guidée : mini-lettre de motivation]
\textbf{Énoncé.} Tu t'appelles Nadège Fotso, tu habites à Douala et tu as vu une annonce pour un poste de guide touristique. Écris 4-5 phrases de candidature avec : \all{sich bewerben um}, \all{Erfahrung}, \all{ich möchte}, \all{ich freue mich auf}.
\tcblower
\textbf{Raisonnement.} Plan : présentation $\rightarrow$ annonce vue $\rightarrow$ candidature avec \all{möchte} + infinitif à la fin $\rightarrow$ argument (expérience, langues) $\rightarrow$ formule de clôture. Vérifier : \all{sich bewerben um} + accusatif ; \all{sich freuen auf} + accusatif (\all{Ihre Antwort}).

\textbf{Solution rédigée.}

\all{Sehr geehrte Damen und Herren,}

\all{ich heiße Nadège Fotso und wohne in Douala. Ich habe Ihre Anzeige gelesen und ich möchte mich um die Stelle als Reiseleiterin bewerben. Ich habe ein Jahr Erfahrung mit Touristen und ich spreche Französisch, Englisch und Deutsch. Ich freue mich auf Ihre Antwort.}

\all{Mit freundlichen Grüßen}

\all{Nadège Fotso}

(Je m'appelle Nadège Fotso et j'habite à Douala. J'ai lu votre annonce et je voudrais postuler pour le poste de guide touristique. J'ai un an d'expérience avec les touristes et je parle français, anglais et allemand. J'attends votre réponse avec plaisir.)

\textbf{Justifications.} \all{möchte} en position 2, infinitif \all{bewerben} à la fin ; \all{um die Stelle} (accusatif) ; \all{auf Ihre Antwort} (accusatif après \all{sich freuen auf}) ; \all{Ihre} avec majuscule (vouvoiement).

\textbf{Conclusion.} Un paragraphe court mais complet : présentation, motivation, arguments, politesse finale.
\end{exoresolu}

\begin{attention}[Les 5 erreurs qui coûtent des points au Bac]
\begin{enumerate}
\item \textbf{Oublier la majuscule aux noms} : écrire \all{die stelle} au lieu de \cle{\all{die Stelle}}. En allemand, \textbf{TOUS} les noms prennent une majuscule : \cle{\all{der Lebenslauf, die Bewerbung, das Gespräch}}.
\item \textbf{Mal placer le verbe avec \cle{\all{möchte}}} : « Ich möchte bewerben mich » est faux. Le modal est en position 2 et l'infinitif à la fin : \cle{\all{Ich möchte mich um die Stelle bewerben.}}
\item \textbf{Faux ami \cle{\all{bekommen}}} : il ne veut pas dire « devenir » mais « recevoir ». « Devenir » se dit \cle{\all{werden}} : \all{Ich möchte Lehrerin werden.} (Je voudrais devenir professeure.)
\item \textbf{Calquer le français « je suis d'accord pour »} : ne traduis pas mot à mot. Dis \all{Ich bin einverstanden.} ou \all{Ich bin mit einem Gespräch einverstanden.} De même, « postuler pour » = \cle{\all{sich bewerben um}} (+ accusatif), jamais \all{für}.
\item \textbf{Oublier la majuscule au vouvoiement et la minuscule après l'appel} : dans une lettre formelle, \cle{\all{Sie, Ihnen, Ihre}} ont toujours une majuscule ; après \all{Sehr geehrte Damen und Herren,} la phrase commence par une minuscule : \all{ich habe Ihre Anzeige gelesen.}
\end{enumerate}
\end{attention}

\newpage

\coursec{Exercices}

\courssub{Je vérifie mes connaissances}

\begin{exercice}[QCM : la lettre de motivation]
Entourez la bonne réponse.
\begin{enumerate}
\item \all{Die Bewerbung} signifie : a) la réponse b) la candidature c) l'invitation.
\item \all{Der Lebenslauf} est : a) le CV b) la lettre c) le diplôme.
\item Une lettre formelle commence par : a) \all{Lieber Peter,} b) \all{Sehr geehrte Damen und Herren,} c) \all{Hallo zusammen,}
\item \all{Ich möchte mich bewerben} signifie : a) je veux partir b) je voudrais postuler c) je dois répondre.
\item La formule de politesse finale d'une lettre formelle est : a) \all{Viele Grüße} b) \all{Bis bald} c) \all{Mit freundlichen Grüßen}.
\end{enumerate}
\end{exercice}

\begin{exercice}[Vrai ou faux — justifiez]
Répondez par \all{richtig} ou \all{falsch} et corrigez la phrase si elle est fausse.
\begin{enumerate}
\item On utilise \all{du} dans une lettre formelle.
\item \all{Möchte} est plus poli que \all{will}.
\item \all{Die Stelle} signifie « le poste ».
\item Dans une lettre informelle, on écrit \all{Sehr geehrte Damen und Herren}.
\item \all{Sich bewerben} est un verbe réfléchi : \all{ich bewerbe mich}.
\end{enumerate}
\end{exercice}

\begin{exercice}[Vocabulaire : trouvez le mot]
Complétez chaque définition avec le mot allemand correct (avec son article).
\begin{enumerate}
\item Le document qui présente les études et les expériences : \all{der \dots}
\item La lettre qu'on envoie pour obtenir un travail : \all{die \dots}
\item La formule du début de la lettre (\all{Sehr geehrte\dots}) : \all{die \dots}
\item La formule de la fin (\all{Mit freundlichen Grüßen}) : \all{die \dots}
\item Le travail qu'on cherche : \all{die \dots}
\item Envoyer une candidature (verbe réfléchi) : \all{sich \dots}
\end{enumerate}
\end{exercice}

\begin{exercice}[Conjugaison : \all{möchte} et \all{würde}]
Complétez le tableau avec les formes correctes.

\begin{center}
\begin{tabular}{|l|l|l|}
\hline
\rowcolor{popblueL} Personne & \all{möchte} & \all{würde} \\
\hline
ich & \dots & \dots \\
\hline
du & \dots & \dots \\
\hline
er/sie/es & \dots & \dots \\
\hline
wir & \dots & \dots \\
\hline
ihr & \dots & \dots \\
\hline
sie/Sie & \dots & \dots \\
\hline
\end{tabular}
\end{center}
\end{exercice}

\courssub{Je m'entraîne}

\begin{exercice}[Texte à trous]
Complétez le texte avec les mots suivants : \all{die Bewerbung, der Lebenslauf, die Stelle, die Anrede, die Grußformel, sich bewerben}.

\begin{textebox}[Ein Brief an die Firma]
Amina möchte \all{(1)} bei einer Firma in Douala. Sie schreibt einen Brief. Am Anfang steht \all{(2)} : „Sehr geehrte Damen und Herren“. Dann schreibt sie, dass sie \all{(3)} als Sekretärin interessant findet. Sie schickt auch \all{(4)} mit ihren Noten und Erfahrungen. Am Ende steht \all{(5)} : „Mit freundlichen Grüßen“. So ist \all{(6)} fertig.
\end{textebox}
\end{exercice}

\begin{exercice}[De la phrase directe à la phrase polie]
Transformez les phrases directes en phrases polies avec \all{möchte} ou \all{würde}.

Exemple : \all{Ich will den Job.} $\rightarrow$ \all{Ich möchte den Job.}
\begin{enumerate}
\item \all{Ich will mich bei Ihnen bewerben.}
\item \all{Gib mir bitte Informationen über die Stelle.}
\item \all{Ich will ein Praktikum in Deutschland machen.}
\item \all{Schick mir bitte den Vertrag.}
\item \all{Ich will Sie bald besuchen.}
\end{enumerate}
\end{exercice}

\begin{exercice}[Appariement : formel ou informel ?]
Classez les expressions suivantes dans le tableau : lettre formelle (F) ou lettre informelle (I).

\all{Liebe Mama,} — \all{Sehr geehrter Herr Mbarga,} — \all{Mit freundlichen Grüßen} — \all{Viele liebe Grüße} — \all{Ich möchte mich bewerben} — \all{Kommst du morgen?} — \all{Ich würde mich über Ihre Antwort freuen} — \all{Bis bald!}

\begin{center}
\begin{tabularx}{\linewidth}{|X|X|}
\hline
\rowcolor{popblueL} Lettre formelle (F) & Lettre informelle (I) \\
\hline
 & \\
\hline
 & \\
\hline
 & \\
\hline
 & \\
\hline
\end{tabularx}
\end{center}
\end{exercice}

\begin{exercice}[Reformulation : lettre informelle $\rightarrow$ lettre formelle]
Réécrivez cette lettre informelle en lettre formelle adressée au directeur d'une entreprise : changez l'adresse (\all{du} $\rightarrow$ \all{Sie}), la formule d'appel et la formule de politesse finale.

\begin{textebox}[Brief an einen Freund]
Lieber Thomas,\\
ich will mich bei deiner Firma bewerben. Ich will ein Praktikum machen. Schick mir bitte die Adresse.\\
Viele Grüße\\
Paul
\end{textebox}
\end{exercice}

\courssub{Je me prépare au Bac}

\begin{exercice}[Type Bac — Compréhension écrite]
Lisez la lettre de motivation suivante, puis répondez aux questions en allemand, par des phrases complètes.

\begin{textebox}[Eine Bewerbung aus Yaoundé]
Yaoundé, den 12. März\\
\\
Sehr geehrte Damen und Herren,\\
ich heiße Sandrine Etoundi und ich bin 17 Jahre alt. Ich besuche das Gymnasium in Yaoundé und ich lerne seit drei Jahren Deutsch. Ich möchte mich um ein Praktikum in Ihrer Firma bewerben, weil ich sehr gut mit Computern arbeite und gern mit Menschen spreche. Ich bin pünktlich, freundlich und motiviert.\\
Ich würde mich sehr freuen, wenn ich bei Ihnen lernen könnte. Mein Lebenslauf und meine Zeugnisse liegen bei. Für weitere Informationen stehe ich Ihnen gern zur Verfügung.\\
Ich danke Ihnen für Ihre Aufmerksamkeit und warte auf Ihre Antwort.\\
\\
Mit freundlichen Grüßen\\
Sandrine Etoundi
\end{textebox}

Questions :
\begin{enumerate}
\item \all{Wer schreibt diesen Brief ?}
\item \all{An wen schreibt sie ?}
\item \all{Warum schreibt sie diesen Brief ?}
\item \all{Welche Qualitäten hat Sandrine ? Nennen Sie drei.}
\item \all{Was schickt sie mit dem Brief ?}
\item \all{Wie endet die formelle Grußformel ?}
\end{enumerate}
\end{exercice}

\begin{exercice}[Type Bac — Thème et version]
\textbf{A. Thème} — Traduisez en allemand :
\begin{enumerate}
\item Je voudrais postuler à ce poste.
\item Je serais très heureux de recevoir votre réponse.
\item Veuillez trouver ci-joint mon CV.
\item Je vous remercie de votre attention.
\end{enumerate}

\textbf{B. Version} — Traduisez en français :

\begin{textebox}[Auszug aus einem formellen Brief]
Sehr geehrte Frau Ngo Bassa,\\
ich möchte mich um die Stelle als Verkäuferin bewerben. Ich würde gern in Ihrem Geschäft arbeiten, weil ich Kontakt mit Kunden liebe. Ich würde mich sehr über eine positive Antwort freuen.\\
Mit freundlichen Grüßen
\end{textebox}
\end{exercice}

\begin{exercice}[Type Bac — Production écrite]
Rédigez une lettre de motivation d'environ 80 mots en allemand.

\begin{textebox}[Sujet]
Sie möchten ein Praktikum in Deutschland machen. Schreiben Sie einen Bewerbungsbrief (etwa 80 Wörter) an eine deutsche Firma. Vergessen Sie nicht : die Anrede, höfliche Formen mit \all{möchte} und \all{würde}, Ihre Qualitäten und die Grußformel.
\end{textebox}

Grille de critères :
\begin{itemize}
\item Formule d'appel formelle correcte (\all{Sehr geehrte Damen und Herren}) : 2 points.
\item Emploi de \all{möchte} et de \all{würde} : 4 points.
\item Présentation de soi et de ses qualités : 4 points.
\item Verbe \all{sich bewerben} correctement employé : 2 points.
\item Formule de politesse finale (\all{Mit freundlichen Grüßen}) : 2 points.
\item Correction grammaticale et ordre des mots : 6 points.
\end{itemize}
\end{exercice}

\newpage

\coursec{Corrigés des exercices}

\begin{corrige}[Exercice 1 — QCM : la lettre de motivation]
\begin{enumerate}
\item \textbf{b) la candidature} — \all{die Bewerbung} désigne la candidature (lettre + CV envoyés pour obtenir un poste).
\item \textbf{a) le CV} — \all{der Lebenslauf} (pluriel : \all{die Lebensläufe}) est le curriculum vitae.
\item \textbf{b) \all{Sehr geehrte Damen und Herren,}} — c'est la formule d'appel d'une lettre formelle. \all{Lieber Peter,} et \all{Hallo zusammen,} sont des formules informelles.
\item \textbf{b) je voudrais postuler} — \all{möchte} exprime le souhait poli ; \all{sich bewerben} signifie « postuler, poser sa candidature ».
\item \textbf{c) \all{Mit freundlichen Grüßen}} — c'est la formule de politesse finale des lettres formelles. \all{Viele Grüße} et \all{Bis bald} sont réservées aux lettres informelles.
\end{enumerate}
\end{corrige}

\begin{corrige}[Exercice 2 — Vrai ou faux — justifiez]
\begin{enumerate}
\item \all{falsch} — Dans une lettre formelle, on utilise \all{Sie} (avec majuscule). \all{du} est réservé aux lettres informelles (amis, famille).
\item \all{richtig} — \all{möchte} (« je voudrais ») est la forme polie de \all{wollen} (« vouloir ») ; \all{will} est trop direct dans une lettre formelle.
\item \all{richtig} — \all{die Stelle} (pluriel : \all{die Stellen}) signifie bien « le poste, l'emploi ».
\item \all{falsch} — Dans une lettre informelle, on écrit par exemple \all{Lieber Thomas,} ou \all{Liebe Mama,}. \all{Sehr geehrte Damen und Herren} est la formule formelle.
\item \all{richtig} — \all{sich bewerben} est bien un verbe réfléchi : \all{ich bewerbe mich}, \all{du bewirbst dich}, \all{er bewirbt sich}. Attention au changement de voyelle radicale \textbf{e → i} aux 2\ieme{} et 3\ieme{} personnes du singulier : \all{du bewirbst}, \all{er bewirbt}.
\end{enumerate}
\end{corrige}

\begin{corrige}[Exercice 3 — Vocabulaire : trouvez le mot]
\begin{enumerate}
\item \all{der Lebenslauf} (pluriel : \all{die Lebensläufe}) — le CV.
\item \all{die Bewerbung} (pluriel : \all{die Bewerbungen}) — la candidature / la lettre de motivation.
\item \all{die Anrede} (pluriel : \all{die Anreden}) — la formule d'appel.
\item \all{die Grußformel} (pluriel : \all{die Grußformeln}) — la formule de politesse finale.
\item \all{die Stelle} (pluriel : \all{die Stellen}) — le poste.
\item \all{sich bewerben} — postuler (conjugaison : \all{ich bewerbe mich}, \all{du bewirbst dich}, \all{er bewirbt sich} ; parfait : \all{ich habe mich beworben}).
\end{enumerate}
\textbf{Point de vigilance :} n'oubliez pas l'article devant chaque nom ; en allemand, le genre fait partie du mot et doit être mémorisé avec lui.
\end{corrige}

\begin{corrige}[Exercice 4 — Conjugaison : \all{möchte} et \all{würde}]
\begin{center}
\begin{tabular}{|l|l|l|}
\hline
\rowcolor{popblueL} Personne & \all{möchte} & \all{würde} \\
\hline
ich & möchte & würde \\
\hline
du & möchtest & würdest \\
\hline
er/sie/es & möchte & würde \\
\hline
wir & möchten & würden \\
\hline
ihr & möchtet & würdet \\
\hline
sie/Sie & möchten & würden \\
\hline
\end{tabular}
\end{center}
\textbf{Justification :} \all{möchte} est le Konjunktiv II de \all{mögen} et \all{würde} celui de \all{werden} ; ce sont des formes de politesse (présent du subjonctif II) figées. Remarquez que les formes de la 1\iere{} et de la 3\ieme{} personne du singulier sont identiques (\all{ich möchte} / \all{er möchte}), comme au Präteritum. Ces deux verbes modaux se placent en 2\ieme{} position et renvoient l'infinitif à la fin de la phrase : \all{Ich würde mich sehr freuen.} (« Je me réjouirais beaucoup. »)
\end{corrige}

\begin{corrige}[Exercice 5 — Texte à trous]
\begin{enumerate}
\item \all{sich bewerben} — après \all{möchte}, l'infinitif va à la fin : \all{Amina möchte sich bewerben}.
\item \all{die Anrede} — la formule d'appel « Sehr geehrte Damen und Herren ».
\item \all{die Stelle} — le poste de secrétaire (ici sujet ou complément d'objet direct selon la phrase, forme inchangée au féminin).
\item \all{den Lebenslauf} — attention, accusatif : \all{Sie schickt auch den Lebenslauf} (le mot de la liste \all{der Lebenslauf} passe à l'accusatif \all{den}).
\item \all{die Grußformel} — la formule finale « Mit freundlichen Grüßen ».
\item \all{die Bewerbung} — la candidature est ainsi terminée : \all{So ist die Bewerbung fertig}.
\end{enumerate}
\textbf{Point de vigilance :} à la question 4, le nominatif \all{der Lebenslauf} devient \all{den Lebenslauf} à l'accusatif (complément direct de \all{schickt}).
\end{corrige}

\begin{corrige}[Exercice 6 — De la phrase directe à la phrase polie]
\begin{enumerate}
\item \all{Ich möchte mich bei Ihnen bewerben.} (« Je voudrais postuler chez vous. ») — \all{möchte} remplace \all{will} ; l'infinitif \all{bewerben} reste à la fin.
\item \all{Ich würde Sie bitten, mir Informationen über die Stelle zu geben.} ou \all{Ich möchte Sie bitten, mir Informationen über die Stelle zu geben.} (« Je vous prierais de me donner des informations sur le poste. ») — dans une lettre formelle, on ne tutoie pas : \all{Gib mir} devient une demande polie avec \all{Sie} + \all{bitten} + infinitif avec \all{zu}.
\item \all{Ich möchte ein Praktikum in Deutschland machen.} (« Je voudrais faire un stage en Allemagne. »)
\item \all{Ich würde Sie bitten, mir den Vertrag zu schicken.} ou \all{Ich möchte Sie bitten, mir den Vertrag zu schicken.} (« Je vous prierais de m'envoyer le contrat. ») — l'impératif direct \all{Schick mir} est transformé en demande polie.
\item \all{Ich würde Sie gern bald besuchen.} ou \all{Ich möchte Sie gern bald besuchen.} (« J'aimerais vous rendre visite bientôt. »)
\end{enumerate}
\textbf{Point de vigilance :} avec \all{möchte} et \all{würde}, l'infinitif se place toujours en fin de phrase ; \all{würde} + \all{gern} correspond au français « j'aimerais ». La construction \all{jemanden bitten, etwas zu tun} (demander à qqn de faire qqch) exige la virgule avant l'infinitif avec \all{zu}.
\end{corrige}

\begin{corrige}[Exercice 7 — Appariement : formel ou informel ?]
\begin{center}
\begin{tabularx}{\linewidth}{|X|X|}
\hline
\rowcolor{popblueL} Lettre formelle (F) & Lettre informelle (I) \\
\hline
Sehr geehrter Herr Mbarga, & Liebe Mama, \\
\hline
Mit freundlichen Grüßen & Viele liebe Grüße \\
\hline
Ich möchte mich bewerben & Kommst du morgen? \\
\hline
Ich würde mich über Ihre Antwort freuen & Bis bald! \\
\hline
\end{tabularx}
\end{center}
\textbf{Justification :} les indices du registre formel sont la formule \all{Sehr geehrte/r}, le vouvoiement \all{Ihre} (avec majuscule) et les tournures de politesse avec \all{möchte} / \all{würde}. Les indices du registre informel sont \all{Liebe/Lieber}, le tutoiement (\all{du}, \all{Kommst du}) et les salutations familières (\all{Viele liebe Grüße}, \all{Bis bald!}).
\end{corrige}

\begin{corrige}[Exercice 8 — Reformulation : lettre informelle $\rightarrow$ lettre formelle]
\textbf{Proposition de corrigé :}
\begin{textebox}[Brief an den Direktor]
Sehr geehrter Herr Direktor,\\
ich möchte mich bei Ihrer Firma bewerben. Ich möchte gern ein Praktikum machen. Ich würde Sie bitten, mir die Adresse zu schicken.\\
Mit freundlichen Grüßen\\
Paul
\end{textebox}
\textbf{Justification des transformations :}
\begin{itemize}
\item \all{Lieber Thomas,} $\rightarrow$ \all{Sehr geehrter Herr Direktor,} : formule d'appel formelle.
\item \all{ich will} $\rightarrow$ \all{ich möchte} : politesse obligatoire dans une lettre formelle.
\item \all{deiner Firma} $\rightarrow$ \all{Ihrer Firma} : le possessif de la 2\ieme{} personne passe du tutoiement (\all{dein}) au vouvoiement (\all{Ihr}, avec majuscule) ; \all{bei} impose le datif (\all{Ihrer}).
\item \all{Schick mir bitte die Adresse.} $\rightarrow$ \all{Ich würde Sie bitten, mir die Adresse zu schicken.} : l'impératif familier devient une demande polie ; notez la construction \all{bitten} + infinitif avec \all{zu} en fin de proposition (et la virgule).
\item \all{Viele Grüße} $\rightarrow$ \all{Mit freundlichen Grüßen} : formule finale formelle.
\end{itemize}
\textbf{Point de vigilance :} \all{Sie} et \all{Ihr/Ihre} prennent toujours une majuscule dans une lettre, même au milieu de la phrase.
\end{corrige}

\begin{corrige}[Exercice 9 — Type Bac : Compréhension écrite]
\textbf{Réponses (en phrases complètes) :}
\begin{enumerate}
\item \all{Sandrine Etoundi schreibt diesen Brief.} (« C'est Sandrine Etoundi qui écrit cette lettre. »)
\item \all{Sie schreibt an eine Firma. / Sie schreibt an die Damen und Herren einer Firma.} (« Elle écrit à une entreprise. »)
\item \all{Sie schreibt diesen Brief, weil sie sich um ein Praktikum in der Firma bewerben möchte.} (« Elle écrit parce qu'elle voudrait postuler à un stage dans l'entreprise. ») — notez la subordonnée avec \all{weil} : le verbe conjugué (\all{möchte}) va à la fin.
\item \all{Sandrine ist pünktlich, freundlich und motiviert.} (« Sandrine est ponctuelle, aimable et motivée. »)
\item \all{Sie schickt ihren Lebenslauf und ihre Zeugnisse mit dem Brief.} (« Elle joint son CV et ses bulletins/diplômes. »)
\item \all{Die formelle Grußformel endet mit „Mit freundlichen Grüßen“.} (« La formule de politesse finale est « Mit freundlichen Grüßen ». »)
\end{enumerate}
\textbf{Barème indicatif :} 1 point par réponse correcte au sens (6 points), 1 point pour la correction grammaticale globale (accords, place du verbe, majuscules), soit 7 points.
\textbf{Points de vigilance :} répondre par des phrases complètes avec verbe conjugué en 2\ieme{} position ; reprendre le vocabulaire du texte sans le recopier intégralement ; ne pas oublier la majuscule aux noms (\all{Brief}, \all{Firma}, \all{Lebenslauf}, \all{Zeugnisse}).
\end{corrige}

\begin{corrige}[Exercice 10 — Type Bac : Thème et version]
\textbf{A. Thème — Traductions proposées :}
\begin{enumerate}
\item \all{Ich möchte mich um diese Stelle bewerben.} — \all{sich bewerben um} + accusatif (« postuler à ») ; l'infinitif \all{bewerben} va à la fin.
\item \all{Ich würde mich sehr über Ihre Antwort freuen.} — \all{sich freuen über} + accusatif (« se réjouir de ») ; \all{Ihre} avec majuscule (vouvoiement formel).
\item \all{Anbei finden Sie meinen Lebenslauf.} ou \all{Mein Lebenslauf liegt bei.} — \all{beiliegen} est un verbe à particule séparable : \all{liegt bei} (3\ieme{} pers. sg).
\item \all{Ich danke Ihnen für Ihre Aufmerksamkeit.} — \all{danken} se construit avec le datif : \all{Ihnen} (et non \all{Sie}) ; \all{für} + accusatif (\all{Ihre Aufmerksamkeit}).
\end{enumerate}
\textbf{B. Version — Traduction proposée :}

« Madame Ngo Bassa,\\
je voudrais postuler au poste de vendeuse. J'aimerais travailler dans votre magasin, parce que j'adore le contact avec les clients. Je me réjouirais beaucoup de recevoir une réponse positive.\\
Je vous prie d'agréer, Madame, l'expression de mes salutations distinguées. »

\textbf{Barème indicatif :} Thème : 1,5 point par phrase (6 points). Version : 4 points (fidélité du sens 2 points, qualité du français 2 points).
\textbf{Points de vigilance :} \all{die Stelle als Verkäuferin} = « le poste de vendeuse » ; \all{sich freuen über} = « se réjouir de » ; la formule \all{Mit freundlichen Grüßen} se traduit par une formule de politesse française complète, pas mot à mot.
\end{corrige}

\begin{corrige}[Exercice 11 — Type Bac : Production écrite]
\textbf{Proposition de rédaction modèle (environ 85 mots) :}
\begin{textebox}[Musterbrief]
Douala, den 15. April\\
\\
Sehr geehrte Damen und Herren,\\
mein Name ist Jean Fotso und ich bin 17 Jahre alt. Ich besuche das Gymnasium in Douala und lerne seit drei Jahren Deutsch. Ich möchte mich um ein Praktikum in Ihrer Firma in Deutschland bewerben. Ich bin fleißig, pünktlich und freundlich, und ich arbeite gern im Team. Ich würde mich sehr freuen, wenn ich bei Ihnen lernen und meine Deutschkenntnisse verbessern könnte. Mein Lebenslauf liegt bei. Ich danke Ihnen für Ihre Aufmerksamkeit.\\
\\
Mit freundlichen Grüßen\\
Jean Fotso
\end{textebox}
\textbf{Barème indicatif (20 points) :}
\begin{itemize}
\item Formule d'appel formelle correcte : 2 points.
\item Emploi correct de \all{möchte} et de \all{würde} : 4 points.
\item Présentation de soi et de ses qualités : 4 points.
\item \all{sich bewerben} correctement employé (pronom réfléchi + préposition \all{um}) : 2 points.
\item Formule finale \all{Mit freundlichen Grüßen} : 2 points.
\item Correction grammaticale et ordre des mots : 6 points.
\end{itemize}
\textbf{Points de vigilance :}
\begin{itemize}
\item Après \all{möchte} et \all{würde}, l'infinitif se place en fin de phrase : \all{Ich möchte mich ... bewerben}.
\item Dans la subordonnée avec \all{wenn}, le verbe conjugué (\all{könnte}) va à la fin ; les infinitifs le précèdent : \all{wenn ich bei Ihnen lernen und meine Deutschkenntnisse verbessern könnte}.
\item Majuscule à tous les noms : \all{Praktikum}, \all{Firma}, \all{Lebenslauf}, \all{Aufmerksamkeit}, \all{Deutschkenntnisse}.
\item Ne pas écrire \all{ich will} : trop direct pour une lettre de motivation.
\item Vérifier l'accord du pronom réfléchi : \all{ich bewerbe \textbf{mich}}, pas \all{sich}.
\end{itemize}
\end{corrige}

\newpage

\coursec{Fiche bilan}

\begin{popfigure}
\begin{center}\tcbox[colback=popblue,colframe=popblue,coltext=white,arc=9pt,boxsep=4pt,left=6pt,right=6pt]{\bfseries\small Briefe: Motivation und Ermutigung}\end{center}
\vspace{-2pt}
\begin{tcbraster}[raster columns=3,raster equal height=rows,raster column skip=6pt,raster row skip=6pt,enhanced,blankest]
\begin{tcolorbox}[enhanced,colback=poporange,colframe=poporange,coltext=white,boxrule=0.9pt,arc=3pt,left=4pt,right=4pt,top=3pt,bottom=3pt,boxsep=1pt,halign=left]\bfseries\scriptsize Texte: Bewerbung + Ermutigungsbrief\end{tcolorbox}
\begin{tcolorbox}[enhanced,colback=popgreen,colframe=popgreen,coltext=white,boxrule=0.9pt,arc=3pt,left=4pt,right=4pt,top=3pt,bottom=3pt,boxsep=1pt,halign=left]\bfseries\scriptsize Lexique: die Bewerbung, der Lebenslauf\end{tcolorbox}
\begin{tcolorbox}[enhanced,colback=poppurple,colframe=poppurple,coltext=white,boxrule=0.9pt,arc=3pt,left=4pt,right=4pt,top=3pt,bottom=3pt,boxsep=1pt,halign=left]\bfseries\scriptsize Lettre formelle vs informelle\end{tcolorbox}
\begin{tcolorbox}[enhanced,colback=popgold,colframe=popgold,coltext=white,boxrule=0.9pt,arc=3pt,left=4pt,right=4pt,top=3pt,bottom=3pt,boxsep=1pt,halign=left]\bfseries\scriptsize Politesse: möchte, würde\end{tcolorbox}
\begin{tcolorbox}[enhanced,colback=poppink,colframe=poppink,coltext=white,boxrule=0.9pt,arc=3pt,left=4pt,right=4pt,top=3pt,bottom=3pt,boxsep=1pt,halign=left]\bfseries\scriptsize Méthode du commentaire\end{tcolorbox}
\begin{tcolorbox}[enhanced,colback=popblue!70,colframe=popblue!70,coltext=white,boxrule=0.9pt,arc=3pt,left=4pt,right=4pt,top=3pt,bottom=3pt,boxsep=1pt,halign=left]\bfseries\scriptsize Production + traduction\end{tcolorbox}
\end{tcbraster}

\legende{Carte mentale de la leçon : les lettres de motivation et d'encouragement.}
\end{popfigure}

\begin{bilan}
\textbf{1. Lexique clé (à connaître avec article et pluriel)}
\begin{itemize}
  \item \all{die Bewerbung, -en} : la candidature ; \all{sich bewerben um + Akk.} : postuler à
  \item \all{der Lebenslauf, „-e} : le CV ; \all{die Stelle, -n} : le poste
  \item \all{die Anrede, -n} : la formule d'appel ; \all{die Grußformel, -n} : la formule de politesse finale
  \item \all{der Absender, -} : l'expéditeur ; \all{der Empfänger, -} : le destinataire
  \item \all{die Ermutigung, -en} : l'encouragement ; \all{jdn. ermutigen} : encourager qqn
  \item \all{die Erfahrung, -en} : l'expérience ; \all{sich interessieren für + Akk.} : s'intéresser à
\end{itemize}

\textbf{2. Grammaire à connaître par cœur}
\begin{center}
\begin{tabularx}{\linewidth}{|X|X|}
\hline
\rowcolor{popblueL} Lettre formelle & Lettre informelle \\
\hline
\all{Sehr geehrte Frau Müller,} / \all{Sehr geehrte Damen und Herren,} & \all{Lieber Paul,} / \all{Liebe Anna,} \\
\hline
\all{Mit freundlichen Grüßen} (MfG) & \all{Viele Grüße} / \all{Dein(e) ...} \\
\hline
Vouvoiement : \all{Sie, Ihnen, Ihr} (majuscule !) & Tutoiement : \all{du, dir, dein} \\
\hline
\all{Ich möchte mich um die Stelle bewerben.} & \all{Ich möchte dich ermutigen.} \\
\hline
\end{tabularx}
\end{center}

\textbf{3. La politesse : \all{möchte} et \all{würde}}
\begin{itemize}
  \item \all{möchte} = « je voudrais » (forme polie de \all{mögen}) : \all{ich/er möchte, du möchtest, wir möchten}. Verbe à l'infinitif en fin de phrase : \all{Ich möchte in Deutschland studieren.}
  \item \all{würde} + infinitif = conditionnel de politesse : \all{Ich würde mich über eine Antwort freuen.} « Je serais heureux d'une réponse. »
  \item Jamais \all{ich will} dans une lettre de motivation : trop direct !
\end{itemize}

\textbf{4. Méthodes express}
\begin{itemize}
  \item \textbf{Commentaire} : 1) identifier le type de lettre (formelle/informelle) ; 2) relever Anrede et Grußformel ; 3) repérer le but (sich bewerben / ermutigen) ; 4) relever 2-3 marques de politesse ; 5) conclure sur l'effet produit.
  \item \textbf{Traduction (thème)} : repérer d'abord le temps, le sujet, la place du verbe (2\up{e} position), puis traduire les formules figées mot à mot apprises par cœur.
\end{itemize}
\end{bilan}

\begin{aretenir}[Checklist avant le Bac]
\begin{itemize}
  \item $\square$ Je connais le vocabulaire avec article et pluriel : \all{die Bewerbung, der Lebenslauf, die Stelle, die Anrede, die Grußformel}.
  \item $\square$ Je sais conjuguer \all{sich bewerben} au Präsens et au Perfekt (\all{ich habe mich beworben}).
  \item $\square$ Je distingue sans hésiter lettre formelle (\all{Sehr geehrte/r ...}) et lettre informelle (\all{Liebe/r ...}).
  \item $\square$ Je mets une majuscule à \all{Sie, Ihnen, Ihr} dans une lettre formelle.
  \item $\square$ Je sais écrire la formule finale : \all{Mit freundlichen Grüßen} (formel) / \all{Viele Grüße} (informel).
  \item $\square$ Je conjugue \all{möchte} correctement : \all{ich möchte, du möchtest, er möchte, wir möchten}.
  \item $\square$ J'utilise \all{würde} + infinitif en fin de phrase pour une demande polie.
  \item $\square$ Je place le verbe conjugué en 2\up{e} position et l'infinitif/participe en fin de phrase.
  \item $\square$ Je n'oublie pas la majuscule à tous les noms allemands.
  \item $\square$ Je sais rédiger 4-5 phrases d'une lettre de motivation simple : présentation, motivation, expérience, formule de politesse.
  \item $\square$ Je sais écrire un court message d'encouragement : \all{Gib nicht auf! Ich glaube an dich!}
  \item $\square$ Je relis ma production : ordre des mots, cas après \all{für/um} (accusatif), orthographe des Umlaute.
\end{itemize}
\end{aretenir}

\findecours

\end{document}
